
\Data\ID{1003061206}
\Updated{2022-04-23 05:04:00}
\Level{A}
\SubArea{Analytic geometry in a plane}
\LANGen{
\Question{
Find the value of a parameter \( a \), such that the line \( ax-6y-15=0 \) is parallel to the line \( y=-\frac23 x+4 \).}
{\( a=-4 \)}{\( a=4 \)}{\( a=-\frac23 \)}{\( a=-2 \)}{}{}}
\LANGpl{
\Question{
Wyznacz wartość parametru \( a \) tak, aby prosta \( ax-6y-15=0 \) była równoległa do prostej  \( y=-\frac23 x+4 \).}
{\( a=-4 \)}{\( a=4 \)}{\( a=-\frac23 \)}{\( a=-2 \)}{}{}}
\LANGcs{
\Question{
Určete hodnotu parametru \( a \) tak, aby přímka \( ax-6y-15=0 \) byla rovnoběžná s přímkou \( y=-\frac23 x+4 \).}
{\( a=-4 \)}{\( a=4 \)}{\( a=-\frac23 \)}{\( a=-2 \)}{}{}}
\LANGsk{
\Question{
Určte hodnotu parametra \( a \) tak, aby priamka \( ax-6y-15=0 \) bola rovnobežná s priamkou \( y=-\frac23 x+4 \).}
{\( a=-4 \)}{\( a=4 \)}{\( a=-\frac23 \)}{\( a=-2 \)}{}{}}
\LANGes{
\Question{
Halla el valor de un parámetro \( a \), suponiendo que la recta \( ax-6y-15=0 \)  es paralela a la recta \( y=-\frac23 x+4 \).}
{\( a=-4 \)}{\( a=4 \)}{\( a=-\frac23 \)}{\( a=-2 \)}{}{}}
\End

\Data\ID{1003061208}
\Updated{2020-07-15 03:07:19}
\Level{A}
\SubArea{Analytic geometry in a plane}
\LANGen{
\Question{
We are given a straight line \( q=\left\{[1+3t;2-2t]\text{, }t\in\mathbb{R}\right\} \). Specify the value of a parameter \( a \) such that the line \( 5x+ay+1=0 \) is parallel to \( q \).}
{\( a=7.5 \)}{\( a=2.5 \)}{\( a=-7.5 \)}{\( a=-2.5 \)}{}{}}
\LANGpl{
\Question{
Dana jest prosta  \( q=\left\{[1+3t;2-2t]\text{, }t\in\mathbb{R}\right\} \). Wyznacz wartość parametru  \( a \) tak, aby prosta \( 5x+ay+1=0 \) była równoległa do  \( q \).}
{\( a=7{,}5 \)}{\( a=2{,}5 \)}{\( a=-7{,}5 \)}{\( a=-2{,}5 \)}{}{}}
\LANGcs{
\Question{
Je dána přímka \( q=\left\{[1+3t;2-2t]\text{, }t\in\mathbb{R}\right\} \). Určete hodnotu parametru \( a \) tak, aby přímka daná rovnicí \( 5x+ay+1=0 \) byla rovnoběžná s přímkou \( q \).}
{\( a=7{,}5 \)}{\( a=2{,}5 \)}{\( a=-7{,}5 \)}{\( a=-2{,}5 \)}{}{}}
\LANGsk{
\Question{
Je daná priamka \( q=\left\{[1+3t;2-2t]\text{, }t\in\mathbb{R}\right\} \). Určte hodnotu parametra \( a \) tak, aby priamka daná rovnicou \( 5x+ay+1=0 \) bola rovnobežná s priamkou \( q \).}
{\( a=7{,}5 \)}{\( a=2{,}5 \)}{\( a=-7{,}5 \)}{\( a=-2{,}5 \)}{}{}}
\LANGes{
\Question{
Dada la recta \( q=\left\{[1+3t;2-2t]\text{, }t\in\mathbb{R}\right\} \). Halla el valor de un parámetro  \( a \) suponiendo que la recta \( 5x+ay+1=0 \) es paralela a la recta \( q \).}
{\( a=7.5 \)}{\( a=2.5 \)}{\( a=-7.5 \)}{\( a=-2.5 \)}{}{}}
\End

\Data\ID{1003061304}
\Updated{2020-07-15 03:07:12}
\Level{A}
\SubArea{Analytic geometry in a plane}
\LANGen{
\Question{
Determine the relative position of the lines
\( p\colon4x-3y+9=0  \) and
\[ \begin{aligned} q\colon x&=6+3t, \\
       y&=11+4t, \end{aligned} \]
 where \( t\in\mathbb{R}\).}
{identical lines, \( p=q \)}{parallel different lines, \( p\parallel q;\ p\neq q \)}{intersecting lines, \( p\cap q=\{[0;3]\} \)}{intersecting lines, \( p\cap q=\{[6;11]\} \)}{}{}}
\LANGpl{
\Question{
Określ wzajemne położenie prostych:
\( p\colon4x-3y+9=0  \) i
\[ \begin{aligned} q\colon x&=6+3t, \\
       y&=11+4t, \end{aligned} \]
gdzie \( t\in\mathbb{R}\).}
{proste identyczne, \( p=q \)}{proste równoległe nie pokrywające się, \( p\parallel q;\ p\neq q \)}{proste przecinające się, \( p\cap q=\{[0;3]\} \)}{proste przecinające się, \( p\cap q=\{[6;11]\} \)}{}{}}
\LANGcs{
\Question{
Vyšetřete vzájemnou polohu přímek
\( p\colon4x-3y+9=0  \) a
\[ \begin{aligned} q\colon x&=6+3t, \\
       y&=11+4t, \end{aligned} \]
 kde \( t\in\mathbb{R}\).}
{totožné přímky, \( p=q \)}{různé rovnoběžky, \( p\parallel q;\ p\neq q \)}{různoběžky, \( p\cap q=\{[0;3]\} \)}{různoběžky, \( p\cap q=\{[6;11]\} \)}{}{}}
\LANGsk{
\Question{
Vyšetrite vzájomnú polohu priamok
\( p\colon4x-3y+9=0  \) a
\[ \begin{aligned} q\colon x&=6+3t, \\
       y&=11+4t, \end{aligned} \]
 kde \( t\in\mathbb{R}\).}
{totožné priamky, \( p=q \)}{rôzne rovnobežky, \( p\parallel q;\ p\neq q \)}{rôznobežky, \( p\cap q=\{[0;3]\} \)}{rôznobežky, \( p\cap q=\{[6;11]\} \)}{}{}}
\LANGes{
\Question{
Determina la posición relativa de las rectas 
\( p\colon4x-3y+9=0  \) y
\[ \begin{aligned} q\colon x&=6+3t, \\
       y&=11+4t, \end{aligned} \]
donde \( t\in\mathbb{R}\).}
{rectas coincidentes, \( p=q \)}{rectas paralelas, \( p\parallel q;\ p\neq q \)}{rectas secantes, \( p\cap q=\{[0;3]\} \)}{rectas secantes, \( p\cap q=\{[6;11]\} \)}{}{}}
\End

\Data\ID{1003061305}
\Updated{2022-04-23 05:04:00}
\Level{A}
\SubArea{Analytic geometry in a plane}
\LANGen{
\Question{
Determine the relative position of the lines \( p\colon 4x+6y-5=0 \) and \( q\colon y=-\frac23 x-6 \).}
{parallel different lines, \( p\parallel q;\ p\neq q \)}{identical lines, \( p=q \)}{intersecting lines, \( p\cap q=\left\{\left[0;\frac54\right]\right\} \)}{intersecting lines, \( p\cap q=\left\{\left[0;\frac56\right]\right\} \)}{}{}}
\LANGpl{
\Question{
Określ wzajemne położenie prostych \( p\colon 4x+6y-5=0 \) i \( q\colon y=-\frac23 x-6 \).}
{proste równoległe, nie pokrywające się, \( p\parallel q;\ p\neq q \)}{proste identyczne, \( p=q \)}{proste przecinające się, \( p\cap q=\left\{\left[0;\frac54\right]\right\} \)}{proste przecinające się, \( p\cap q=\left\{\left[0;\frac56\right]\right\} \)}{}{}}
\LANGcs{
\Question{
Vyšetřete vzájemnou polohu přímek \( p\colon 4x+6y-5=0 \) a \( q\colon y=-\frac23 x-6 \).}
{různé rovnoběžky, \( p\parallel q;\ p\neq q \)}{totožné přímky, \( p=q \)}{různoběžky, \( p\cap q=\left\{\left[0;\frac54\right]\right\} \)}{různoběžky, \( p\cap q=\left\{\left[0;\frac56\right]\right\} \)}{}{}}
\LANGsk{
\Question{
Vyšetrite vzájomnú polohu priamok \( p\colon 4x+6y-5=0 \) a \( q\colon y=-\frac23 x-6 \).}
{rôzne rovnobežky, \( p\parallel q;\ p\neq q \)}{totožné priamky, \( p=q \)}{rôznobežky, \( p\cap q=\left\{\left[0;\frac54\right]\right\} \)}{rôznobežky, \( p\cap q=\left\{\left[0;\frac56\right]\right\} \)}{}{}}
\LANGes{
\Question{
Determina la posición relativa de las rectas \( p\colon 4x+6y-5=0 \) y \( q\colon y=-\frac23 x-6 \).}
{rectas paralelas, \( p\parallel q;\ p\neq q \)}{rectas coincidentes, \( p=q \)}{rectas secantes, \( p\cap q=\left\{\left[0;\frac54\right]\right\} \)}{rectas secantes, \( p\cap q=\left\{\left[0;\frac56\right]\right\} \)}{}{}}
\End

\Data\ID{1003061306}
\Updated{2020-07-15 03:07:12}
\Level{A}
\SubArea{Analytic geometry in a plane}
\LANGen{
\Question{
Determine the relative position of the lines \( p\colon 2x-3y+7=0 \) and
\[ \begin{aligned} q\colon   x& =2+t, \\
       y& = -3t, \end{aligned} \]
where \( t\in\mathbb{R} \).}
{intersecting lines, \( p\cap q=\left\{\left[1;3\right]\right\} \)}{identical lines, \( p=q \)}{parallel different lines, \( p\parallel q;\ p\neq q \)}{intersecting lines, \( p\cap q=\left\{\left[7;7\right]\right\} \)}{}{}}
\LANGpl{
\Question{
Określ wzajemne położenie prostych \( p\colon 2x-3y+7=0 \) i
\[ \begin{aligned} q\colon   x& =2+t, \\
       y& = -3t, \end{aligned} \]
gdzie \( t\in\mathbb{R} \).}
{proste przecinające się, \( p\cap q=\left\{\left[1;3\right]\right\} \)}{proste identyczne, \( p=q \)}{proste równoległe, nie pokrywające się, \( p\parallel q;\ p\neq q \)}{proste przecinające się, \( p\cap q=\left\{\left[7;7\right]\right\} \)}{}{}}
\LANGcs{
\Question{
Vyšetřete vzájemnou polohu přímek \( p\colon 2x-3y+7=0 \) a
\[ \begin{aligned} q\colon   x& =2+t, \\
       y& = -3t, \end{aligned} \]
kde \( t\in\mathbb{R} \).}
{různoběžky, \( p\cap q=\left\{\left[1;3\right]\right\} \)}{totožné přímky, \( p=q \)}{různé rovnoběžky, \( p\parallel q;\ p\neq q \)}{různoběžky, \( p\cap q=\left\{\left[7;7\right]\right\} \)}{}{}}
\LANGsk{
\Question{
Vyšetrite vzájomnú polohu priamok \( p\colon 2x-3y+7=0 \) a
\[ \begin{aligned} q\colon   x& =2+t, \\
       y& = -3t, \end{aligned} \]
kde \( t\in\mathbb{R} \).}
{rôznobežky, \( p\cap q=\left\{\left[1;3\right]\right\} \)}{totožné priamky, \( p=q \)}{rôzne rovnobežky, \( p\parallel q;\ p\neq q \)}{rôznobežky, \( p\cap q=\left\{\left[7;7\right]\right\} \)}{}{}}
\LANGes{
\Question{
Determina la posición relativa de las rectas \( p\colon 2x-3y+7=0 \) y
\[ \begin{aligned} q\colon   x& =2+t, \\
       y& = -3t, \end{aligned} \]
donde \( t\in\mathbb{R} \).}
{rectas secantes, \( p\cap q=\left\{\left[1;3\right]\right\} \)}{rectas coincidentes, \( p=q \)}{rectas paralelas, \( p\parallel q;\ p\neq q \)}{rectas secantes, \( p\cap q=\left\{\left[7;7\right]\right\} \)}{}{}}
\End

\Data\ID{1103061201}
\Updated{2021-12-05 07:12:31}
\Level{A}
\SubArea{Analytic geometry in a plane}
\IMG{1103061201.pdf}
\LANGen{
\Question{
From the following list, choose parametric equations, that do not represent the straight line passing through the points \( A \) and \( B \) (see the picture).}
{$\begin{aligned}
p\colon x&=2+4t, \\
y&=6+2t;\  t\in\mathbb{R}
\end{aligned}$}{$\begin{aligned}
p\colon x&=2+2t, \\
y&=1+t;\ t\in\mathbb{R}
\end{aligned}$}{$\begin{aligned}
p\colon x&=6+4t, \\
y&=3+2t;\ t\in\mathbb{R}
\end{aligned}$}{$\begin{aligned}
p\colon x&=2-2t, \\
y&=1-t;\ t\in\mathbb{R}
\end{aligned}$}{$\begin{aligned}
p\colon x&=4+4t, \\
y&=2+2t;\ t\in\mathbb{R}
\end{aligned}$}{}}
\LANGpl{
\Question{
Z poniższej listy wybierz równania parametryczne, które  nie określają prostej przechodzącej przez punkty \( A \) i \( B \) (spójrz na rysunek).}
{$\begin{aligned}
p\colon x&=2+4t, \\
y&=6+2t;\  t\in\mathbb{R}
\end{aligned}$}{$\begin{aligned}
p\colon x&=2+2t, \\
y&=1+t;\ t\in\mathbb{R}
\end{aligned}$}{$\begin{aligned}
p\colon x&=6+4t, \\
y&=3+2t;\ t\in\mathbb{R}
\end{aligned}$}{$\begin{aligned}
p\colon x&=2-2t, \\
y&=1-t;\ t\in\mathbb{R}
\end{aligned}$}{$\begin{aligned}
p\colon x&=4+4t, \\
y&=2+2t;\ t\in\mathbb{R}
\end{aligned}$}{}}
\LANGcs{
\Question{
Z nabízených možností vyberte parametrické rovnice, které nevyjadřují přímku procházející body \( A \) a \( B \)  (viz obrázek).}
{$\begin{aligned}
p\colon x&=2+4t, \\
y&=6+2t;\  t\in\mathbb{R}
\end{aligned}$}{$\begin{aligned}
p\colon x&=2+2t, \\
y&=1+t;\ t\in\mathbb{R}
\end{aligned}$}{$\begin{aligned}
p\colon x&=6+4t, \\
y&=3+2t;\ t\in\mathbb{R}
\end{aligned}$}{$\begin{aligned}
p\colon x&=2-2t, \\
y&=1-t;\ t\in\mathbb{R}
\end{aligned}$}{$\begin{aligned}
p\colon x&=4+4t, \\
y&=2+2t;\ t\in\mathbb{R}
\end{aligned}$}{}}
\LANGsk{
\Question{
Z ponúknutých možností vyberte parametrické rovnice, ktoré nevyjadrujú priamku prechádzajúcu bodmi \( A \) a \( B \)  (viď obrázok).}
{$\begin{aligned}
p\colon x&=2+4t, \\
y&=6+2t;\  t\in\mathbb{R}
\end{aligned}$}{$\begin{aligned}
p\colon x&=2+2t, \\
y&=1+t;\ t\in\mathbb{R}
\end{aligned}$}{$\begin{aligned}
p\colon x&=6+4t, \\
y&=3+2t;\ t\in\mathbb{R}
\end{aligned}$}{$\begin{aligned}
p\colon x&=2-2t, \\
y&=1-t;\ t\in\mathbb{R}
\end{aligned}$}{$\begin{aligned}
p\colon x&=4+4t, \\
y&=2+2t;\ t\in\mathbb{R}
\end{aligned}$}{}}
\LANGes{
\Question{
De la siguiente lista, elige las ecuaciones paramétricas, que   no representan  la  recta que pasa por los puntos \( A \) y \( B \) (mira la imagen).}
{$\begin{aligned}
p\colon x&=2+4t, \\
y&=6+2t;\  t\in\mathbb{R}
\end{aligned}$}{$\begin{aligned}
p\colon x&=2+2t, \\
y&=1+t;\ t\in\mathbb{R}
\end{aligned}$}{$\begin{aligned}
p\colon x&=6+4t, \\
y&=3+2t;\ t\in\mathbb{R}
\end{aligned}$}{$\begin{aligned}
p\colon x&=2-2t, \\
y&=1-t;\ t\in\mathbb{R}
\end{aligned}$}{$\begin{aligned}
p\colon x&=4+4t, \\
y&=2+2t;\ t\in\mathbb{R}
\end{aligned}$}{}}
\End

\Data\ID{1103061202}
\Updated{2022-04-23 05:04:00}
\Level{A}
\SubArea{Analytic geometry in a plane}
\IMG{1103061202.pdf}
\LANGen{
\Question{
A straight line \( p \) is given by the point \( A \) and the normal vector \( \vec{n} \) (see the picture). Find its equation in general form.}
{\( p\colon 2x-5y-6=0 \)}{\( p\colon 2x+5y-6=0 \)}{\( p\colon 5x-2y-15=0 \)}{\( p\colon 5x+2y-15=0 \)}{}{}}
\LANGpl{
\Question{
Dana jest prosta \( p \) określona punktem \( A \) oraz wektorem normalnym \( \vec{n} \) (spójrz na rysunek). Wyznacz ogólne równanie prostej.}
{\( p\colon 2x-5y-6=0 \)}{\( p\colon 2x+5y-6=0 \)}{\( p\colon 5x-2y-15=0 \)}{\( p\colon 5x+2y-15=0 \)}{}{}}
\LANGcs{
\Question{
Přímka \( p \) je dána bodem \( A \) a normálovým vektorem \( \vec{n} \) (viz obrázek). Určete její obecnou rovnici.}
{\( p\colon 2x-5y-6=0 \)}{\( p\colon 2x+5y-6=0 \)}{\( p\colon 5x-2y-15=0 \)}{\( p\colon 5x+2y-15=0 \)}{}{}}
\LANGsk{
\Question{
Priamka \( p \) je daná bodom \( A \) a normálovým vektorom \( \vec{n} \) (viď obrázok). Určte jej všeobecnú rovnicu.}
{\( p\colon 2x-5y-6=0 \)}{\( p\colon 2x+5y-6=0 \)}{\( p\colon 5x-2y-15=0 \)}{\( p\colon 5x+2y-15=0 \)}{}{}}
\LANGes{
\Question{
La recta \( p \) viene dada por el punto \( A \) y el vector normal \( \vec{n} \) (mira la imagen). Determina su ecuación general.}
{\( p\colon 2x-5y-6=0 \)}{\( p\colon 2x+5y-6=0 \)}{\( p\colon 5x-2y-15=0 \)}{\( p\colon 5x+2y-15=0 \)}{}{}}
\End

\Data\ID{1103061203}
\Updated{2021-12-05 07:12:17}
\Level{A}
\SubArea{Analytic geometry in a plane}
\IMG{1103061203.pdf}
\LANGen{
\Question{
A straight line \( p \) is given by the point \( A \) and the direction angle \( \varphi \) (see the picture). Choose the equation of the line \( p \) in the slope-intercept form.}
{\( p\colon y=-\sqrt3x+3 \)}{\( p\colon y=\sqrt3x+3 \)}{\( p\colon y=1.7x+3 \)}{\( p\colon y=-1.7x+3 \)}{}{}}
\LANGpl{
\Question{
Dana jest prosta  \( p \) określona punktem \( A \) oraz kątem kierunkowym \( \varphi \) (spójrz na rysunek). Wybierz równanie kierunkowe prostej \( p \) .}
{\( p\colon y=-\sqrt3x+3 \)}{\( p\colon y=\sqrt3x+3 \)}{\( p\colon y=1{,}7x+3 \)}{\( p\colon y=-1{,}7x+3 \)}{}{}}
\LANGcs{
\Question{
Přímka \( p \) je dána bodem \( A \) a směrovým úhlem \( \varphi \) (viz obrázek). Vyberte rovnici přímky \( p \) ve směrnicovém tvaru.}
{\( p\colon y=-\sqrt3x+3 \)}{\( p\colon y=\sqrt3x+3 \)}{\( p\colon y=1{,}7x+3 \)}{\( p\colon y=-1{,}7x+3 \)}{}{}}
\LANGsk{
\Question{
Priamka \( p \) je daná bodom \( A \) a smerovým uhlom \( \varphi \) (viď obrázok). Vyberte rovnicu priamky \( p \) v smernicovom tvare.}
{\( p\colon y=-\sqrt3x+3 \)}{\( p\colon y=\sqrt3x+3 \)}{\( p\colon y=1{,}7x+3 \)}{\( p\colon y=-1{,}7x+3 \)}{}{}}
\LANGes{
\Question{
La recta \( p \) viene dada por el punto \( A \) y el ángulo director \( \varphi \) (mira la imagen). Determina la ecuación de la recta \( p \) en forma explícita.}
{\( p\colon y=-\sqrt3x+3 \)}{\( p\colon y=\sqrt3x+3 \)}{\( p\colon y=1.7x+3 \)}{\( p\colon y=-1.7x+3 \)}{}{}}
\End

\Data\ID{1103061204}
\Updated{2023-07-04 09:07:11}
\Level{A}
\SubArea{Analytic geometry in a plane}
\IMG{1103061204.pdf}
\LANGen{
\Question{
From the following list choose the equation of a straight line that passes through the given point \( K \) and is not parallel to the given line \( m \) (see the picture).}
{\( g\colon y=-\frac32x+\frac{13}2 \)}{\( b\colon 2x-3y-13=0 \)}{\( f\colon y=\frac23x-\frac{13}3 \)}{$\begin{aligned}
q\colon x&=5+3t, \\
y&=-1+2t;\ t\in\mathbb{R}
\end{aligned}$}{}{}}
\LANGpl{
\Question{
Wybierz równanie prostej przechodzącej przez punkt \( K \) oraz  nie równoległej do prostej \( m \) (spójrz na rysunek).}
{\( g\colon y=-\frac32x+\frac{13}2 \)}{\( b\colon 2x-3y-13=0 \)}{\( f\colon y=\frac23x-\frac{13}3 \)}{$\begin{aligned}
q\colon x&=5+3t, \\
y&=-1+2t;\ t\in\mathbb{R}
\end{aligned}$}{}{}}
\LANGcs{
\Question{
Z uvedených možností vyberte rovnici přímky, která prochází bodem \( K \) a není rovnoběžná s danou přímkou \( m \) (viz obrázek).}
{\( g\colon y=-\frac32x+\frac{13}2 \)}{\( b\colon 2x-3y-13=0 \)}{\( f\colon y=\frac23x-\frac{13}3 \)}{$\begin{aligned}
q\colon x&=5+3t, \\
y&=-1+2t;\ t\in\mathbb{R}
\end{aligned}$}{}{}}
\LANGsk{
\Question{
Z uvedených možností vyberte rovnicu priamky, ktorá prechádza bodom \( K \) a nie je rovnobežná s danou priamkou \( m \) (viď obrázok).}
{\( g\colon y=-\frac32x+\frac{13}2 \)}{\( b\colon 2x-3y-13=0 \)}{\( f\colon y=\frac23x-\frac{13}3 \)}{$\begin{aligned}
q\colon x&=5+3t, \\
y&=-1+2t;\ t\in\mathbb{R}
\end{aligned}$}{}{}}
\LANGes{
\Question{
De la siguiente lista, elige la  ecuación de la recta que pasa por el punto \( K \) y  no es paralela   a la recta \( m \) (mira la imagen).}
{\( g\colon y=-\frac32x+\frac{13}2 \)}{\( b\colon 2x-3y-13=0 \)}{\( f\colon y=\frac23x-\frac{13}3 \)}{$\begin{aligned}
q\colon x&=5+3t, \\
y&=-1+2t;\ t\in\mathbb{R}
\end{aligned}$}{}{}}
\End

\Data\ID{1103061205}
\Updated{2023-07-04 09:07:46}
\Level{A}
\SubArea{Analytic geometry in a plane}
\IMG{1103061205.pdf}
\LANGen{
\Question{
From the following list choose the equation of a straight line that passes through the given point \( K \) and is not perpendicular to the given line \( m \) (see the picture).}
{\( r\colon y=\frac23x-\frac{13}3 \)}{\( p\colon 3x+2y-13=0 \)}{\( s\colon y=-\frac32x+\frac{13}2 \)}{$\begin{aligned}
q\colon x&=5+2t, \\
y&=-1-3t;\ t\in\mathbb{R}
\end{aligned}$}{}{}}
\LANGpl{
\Question{
Wybierz równanie prostej przechodzącej przez punkt  \( K \) oraz  nie prostopadłej do prostej \( m \) (spójrz na rysunek).}
{\( r\colon y=\frac23x-\frac{13}3 \)}{\( p\colon 3x+2y-13=0 \)}{\( s\colon y=-\frac32x+\frac{13}2 \)}{$\begin{aligned}
q\colon x&=5+2t, \\
y&=-1-3t;\ t\in\mathbb{R}
\end{aligned}$}{}{}}
\LANGcs{
\Question{
Z uvedených možností vyberte rovnici přímky, která prochází daným bodem \( K \) a není kolmá k dané přímce \( m \) (viz obrázek).}
{\( r\colon y=\frac23x-\frac{13}3 \)}{\( p\colon 3x+2y-13=0 \)}{\( s\colon y=-\frac32x+\frac{13}2 \)}{$\begin{aligned}
q\colon x&=5+2t, \\
y&=-1-3t;\ t\in\mathbb{R}
\end{aligned}$}{}{}}
\LANGsk{
\Question{
Z uvedených možností vyberte rovnicu priamky, ktorá prechádza daným bodom \( K \) a nie je kolmá k danej priamke \( m \) (viď obrázok).}
{\( r\colon y=\frac23x-\frac{13}3 \)}{\( p\colon 3x+2y-13=0 \)}{\( s\colon y=-\frac32x+\frac{13}2 \)}{$\begin{aligned}
q\colon x&=5+2t, \\
y&=-1-3t;\ t\in\mathbb{R}
\end{aligned}$}{}{}}
\LANGes{
\Question{
De la siguiente lista, elige la ecuación de la recta que pasa por el punto\( K \) y  no es perpendicular a la rectar \( m \) (mira la imagen).}
{\( r\colon y=\frac23x-\frac{13}3 \)}{\( p\colon 3x+2y-13=0 \)}{\( s\colon y=-\frac32x+\frac{13}2 \)}{$\begin{aligned}
q\colon x&=5+2t, \\
y&=-1-3t;\ t\in\mathbb{R}
\end{aligned}$}{}{}}
\End

\Data\ID{1103061207}
\Updated{2021-12-05 07:12:24}
\Level{A}
\SubArea{Analytic geometry in a plane}
\IMG{1103061207.pdf}
\LANGen{
\Question{
We are given the straight line \( m= \left\{[3-t;t]\text{, } t\in\mathbb{R} \right\} \) which intersects lines \( a \), \( b \), \( c \) in points \( A \), \( B \), \( C \) consecutively (see the picture). Find the values of a parameter \( t \) corresponding to these line intersections.}
{\( t_A=1; t_B=\frac32;\ t_C=2 \)}{\( t_A=-1; t_B=-2;\ t_C=-3 \)}{\( t_A=2; t_B=\frac32;\ t_C=1 \)}{\( t_A=2; t_B=\frac52;\ t_C=3 \)}{}{}}
\LANGpl{
\Question{
Dana jest prosta \( m= \left\{[3-t;t]\text{, } t\in\mathbb{R} \right\} \) przecinające proste \( a \), \( b \), \( c \) w punktach \( A \), \( B \), \( C \)  (spójrz na rysunek). Wyznacz wartość parametru \( t \) odpowiadające tym przecięciom.}
{\( t_A=1; t_B=\frac32;\ t_C=2 \)}{\( t_A=-1; t_B=-2;\ t_C=-3 \)}{\( t_A=2; t_B=\frac32;\ t_C=1 \)}{\( t_A=2; t_B=\frac52;\ t_C=3 \)}{}{}}
\LANGcs{
\Question{
Je dána přímka \( m= \left\{[3-t;t]\text{, } t\in\mathbb{R} \right\} \), která protíná přímky \( a \), \( b \), \( c \) po řadě v bodech \( A \), \( B \), \( C \) (viz obrázek). Určete hodnoty parametru \( t \) odpovídající těmto průsečíkům.}
{\( t_A=1; t_B=\frac32;\ t_C=2 \)}{\( t_A=-1; t_B=-2;\ t_C=-3 \)}{\( t_A=2; t_B=\frac32;\ t_C=1 \)}{\( t_A=2; t_B=\frac52;\ t_C=3 \)}{}{}}
\LANGsk{
\Question{
Je daná priamka \( m= \left\{[3-t;t]\text{, } t\in\mathbb{R} \right\} \), ktorá pretína priamky \( a \), \( b \), \( c \) po rade v bodoch \( A \), \( B \), \( C \) (viď obrázok). Určte hodnoty parametra \( t \) ktorý odpovedá týmto priesečníkom.}
{\( t_A=1; t_B=\frac32;\ t_C=2 \)}{\( t_A=-1; t_B=-2;\ t_C=-3 \)}{\( t_A=2; t_B=\frac32;\ t_C=1 \)}{\( t_A=2; t_B=\frac52;\ t_C=3 \)}{}{}}
\LANGes{
\Question{
Dada la recta  \( m= \left\{[3-t;t]\text{, } t\in\mathbb{R} \right\} \) que corta a las rectas \( a \), \( b \), \( c \) en los puntos \( A \), \( B \), \( C \) consecutivamente (mira la imagen). Halla los valores del parámetro \( t \) que corresponden a estas intersecciones.}
{\( t_A=1; t_B=\frac32;\ t_C=2 \)}{\( t_A=-1; t_B=-2;\ t_C=-3 \)}{\( t_A=2; t_B=\frac32;\ t_C=1 \)}{\( t_A=2; t_B=\frac52;\ t_C=3 \)}{}{}}
\End

\Data\ID{1103061302}
\Updated{2020-07-15 03:07:12}
\Level{A}
\SubArea{Analytic geometry in a plane}
\IMG{1103061302.pdf}
\LANGen{
\Question{
Let there be straight lines \( p\colon x+4y-16=0 \) and \( q\colon y= \frac18 x+b \), where \( b \) is a real parameter. Specify the value of \( b \), such that \( p \) and \( q \) do intersect on \( x \)-axis.}
{\( b=-2 \)}{\( b=-4 \)}{\( b=2 \)}{\( b=0 \)}{}{}}
\LANGpl{
\Question{
Dane są proste \( p\colon x+4y-16=0 \) i \( q\colon y= \frac18 x+b \), gdzie \( b \) to parametr rzeczywisty. Określ wartość \( b \) tak, aby \( p \) i \( q \) przecinały się na osi \( x \).}
{\( b=-2 \)}{\( b=-4 \)}{\( b=2 \)}{\( b=0 \)}{}{}}
\LANGcs{
\Question{
Jsou dány přímky \( p\colon x+4y-16=0 \) a \( q\colon y= \frac18 x+b \), kde \( b \) je reálný parametr. Určete hodnotu parametru \( b \) tak, aby se přímky \( p \) a \( q \) protínaly na ose \( x \).}
{\( b=-2 \)}{\( b=-4 \)}{\( b=2 \)}{\( b=0 \)}{}{}}
\LANGsk{
\Question{
Sú dané priamky \( p\colon x+4y-16=0 \) a \( q\colon y= \frac18 x+b \), kde \( b \) je reálny parameter. Určte hodnotu parametra \( b \) tak, aby sa priamky \( p \) a \( q \) pretínali na osi \( x \).}
{\( b=-2 \)}{\( b=-4 \)}{\( b=2 \)}{\( b=0 \)}{}{}}
\LANGes{
\Question{
Dadas las rectas \( p\colon x+4y-16=0 \) y \( q\colon y= \frac18 x+b \), donde \( b \) es un parámetro real. Halla el valor de \( b \), suponiendo que \( p \) y \( q \)  cortan  al eje \( x \).}
{\( b=-2 \)}{\( b=-4 \)}{\( b=2 \)}{\( b=0 \)}{}{}}
\End

\Data\ID{1103061303}
\Updated{2020-07-15 03:07:12}
\Level{A}
\SubArea{Analytic geometry in a plane}
\IMG{1103061303.pdf}
\LANGen{
\Question{
Let there be a straight line \( p\colon 5x-y-10=0 \). Choose the equation of a straight line \( q \) that passes through the point \( A=[-2;2] \) and intersects with \( p \) on \( y \)-axis.}
{\( q\colon y=-6x-10 \)}{\( q\colon y=-5x-10 \)}{\( q\colon y=-5x-8 \)}{\( q\colon y=-6x-8 \)}{}{}}
\LANGpl{
\Question{
Dana jest prosta \( p\colon 5x-y-10=0 \). Wybierz równanie prostej \( q \) przechodzącej przez punkt \( A=[-2;2] \) i przecinającą  \( p \) na osi \( y \).}
{\( q\colon y=-6x-10 \)}{\( q\colon y=-5x-10 \)}{\( q\colon y=-5x-8 \)}{\( q\colon y=-6x-8 \)}{}{}}
\LANGcs{
\Question{
Je dána přímka \( p\colon 5x-y-10=0 \). Vyberte rovnici přímky \( q \), která prochází bodem \( A=[-2;2] \) a s přímkou \( p \) se protíná na ose \( y \).}
{\( q\colon y=-6x-10 \)}{\( q\colon y=-5x-10 \)}{\( q\colon y=-5x-8 \)}{\( q\colon y=-6x-8 \)}{}{}}
\LANGsk{
\Question{
Je daná priamka \( p\colon 5x-y-10=0 \). Vyberte rovnicu priamky \( q \), ktorá prechádza bodom \( A=[-2;2] \) a s priamkou \( p \) sa pretína na osi \( y \).}
{\( q\colon y=-6x-10 \)}{\( q\colon y=-5x-10 \)}{\( q\colon y=-5x-8 \)}{\( q\colon y=-6x-8 \)}{}{}}
\LANGes{
\Question{
Dada la recta \( p\colon 5x-y-10=0 \). Elige la ecuación de una recta \( q \) que pasa por el punto \( A=[-2;2] \) y corta con \( p \) en el eje \( y \).}
{\( q\colon y=-6x-10 \)}{\( q\colon y=-5x-10 \)}{\( q\colon y=-5x-8 \)}{\( q\colon y=-6x-8 \)}{}{}}
\End

\Data\ID{1103090806}
\Updated{2020-07-15 03:07:10}
\Level{A}
\SubArea{Analytic geometry in a plane}
\LANGen{
\Question{
We are given the line segment \( AB \):
\begin{align*}
x&=2+2t, \\
y&=-1+t;\ t\in [0;1],
\end{align*}
and the points \( K=\left[\frac72;-\frac14\right] \), \( L=[-2;-3] \) and \( M=\left[5;\frac12\right] \). Choose a picture where the mutual position of the points \( A \), \( B \), \( K \), \( L \), and \( M \) is indicated correctly.}
{}{}{}{}{}{}}
\LANGpl{
\Question{
Dany jest odcinek \( AB \):
\begin{align*}
x&=2+2t, \\
y&=-1+t;\ t\in\langle0;1\rangle,
\end{align*}
oraz punkty \( K=\left[\frac72;-\frac14\right] \), \( L=[-2;-3] \) i \( M=\left[5;\frac12\right] \). Wybierz rysunek, na którym znajduje się poprawne wzajemne położenie punktów  \( A \), \( B \), \( K \), \( L \), i \( M \).}
{}{}{}{}{}{}}
\LANGcs{
\Question{
Je dána úsečka \( AB \):
\begin{align*}
x&=2+2t, \\
y&=-1+t;\ t\in\langle0;1\rangle,
\end{align*}
a body \( K=\left[\frac72;-\frac14\right] \), \( L=[-2;-3] \) a \( M=\left[5;\frac12\right] \). Vyberte obrázek, na kterém je správně vyznačena vzájemná poloha všech pěti bodů \( A \), \( B \), \( K \), \( L \) a \( M \).}
{}{}{}{}{}{}}
\LANGsk{
\Question{
Je daná úsečka \( AB \):
\begin{align*}
x&=2+2t, \\
y&=-1+t;\ t\in\langle0;1\rangle,
\end{align*}
a body \( K=\left[\frac72;-\frac14\right] \), \( L=[-2;-3] \) a \( M=\left[5;\frac12\right] \). Vyberte obrázok, na ktorom je správne vyznačená vzájomná poloha všetkých piatich bodov \( A \), \( B \), \( K \), \( L \) a \( M \).}
{}{}{}{}{}{}}
\LANGes{
\Question{
Dado el segmento  \( AB \):
\begin{align*}
x&=2+2t, \\
y&=-1+t;\ t\in [0;1],
\end{align*}
y los puntos \( K=\left[\frac72;-\frac14\right] \), \( L=[-2;-3] \) y \( M=\left[5;\frac12\right] \). Elige la imagen donde está correctamente marcada la posición recíproca de los puntos \( A \), \( B \), \( K \), \( L \), y \( M \).}
{}{}{}{}{}{}}

\IMGanswer{1103090806a.pdf}
\IMGanswer{1103090806b.pdf}
\IMGanswer{1103090806c.pdf}
\IMGanswer{1103090806d.pdf}\End

\Data\ID{2010014201}
\Updated{2022-05-23 02:05:48}
\Level{A}
\SubArea{Analytic geometry in a plane}
\LANGen{
\Question{
Find the value of a parameter \( a \), such that the line \( ax-4y-12=0 \) is parallel to the line \( y=-\frac32 x+4 \).}
{\( a=-6\)}{\( a=-\frac32\)}{\( a=4\)}{\( a=\frac23\)}{}{}}
\LANGpl{
\Question{
Znajdź wartość parametru \( a \), taką, że prosta \( ax-4y-12=0 \) jest równoległa do prostej \( y=-\frac32 x+4 \).}
{\( a=-6\)}{\( a=-\frac32\)}{\( a=4\)}{\( a=\frac23\)}{}{}}
\LANGcs{
\Question{
Určete hodnotu parametru \( a \) tak, aby byla přímka \( ax-4y-12=0 \) rovnoběžná s přímkou \( y=-\frac32 x+4 \).}
{\( a=-6\)}{\( a=-\frac32\)}{\( a=4\)}{\( a=\frac23\)}{}{}}
\LANGsk{
\Question{
Určte hodnotu parametra \( a \) tak, aby priamka \( ax-4y-12=0 \) bola rovnobežná s priamkou \( y=-\frac32 x+4 \).}
{\( a=-6\)}{\( a=-\frac32\)}{\( a=4\)}{\( a=\frac23\)}{}{}}
\LANGes{
\Question{
Halla el valor del parámetro \( a \) suponiendo que la recta \( ax-4y-12=0 \) es paralela a la recta \( y=-\frac32 x+4 \).}
{\( a=-6\)}{\( a=-\frac32\)}{\( a=4\)}{\( a=\frac23\)}{}{}}
\End

\Data\ID{2010014202}
\Updated{2022-05-16 11:05:13}
\Level{A}
\SubArea{Analytic geometry in a plane}
\LANGen{
\Question{
Determine the relative position of the lines \( p\colon 6x+4y+8=0 \) and \( q\colon y=-\frac32 x+2 \).}
{parallel different lines, \( p\parallel q;\ p\neq q \)}{intersecting lines, \( p\cap q=\left\{\left[0;-2\right]\right\} \)}{intersecting lines, \( p\cap q=\left\{\left[0;2\right]\right\} \)}{identical lines, \( p=q \)}{}{}}
\LANGpl{
\Question{
Określ wzajemne położenie prostych \( p\colon 6x+4y+8=0 \) i \( q\colon y=-\frac32 x+2 \).}
{równoległe różne proste, \( p\parallel q;\ p\neq q \)}{przecinające się proste, \( p\cap q=\left\{\left[0;-2\right]\right\} \)}{przecinające się proste \( p\cap q=\left\{\left[0;2\right]\right\} \)}{identyczne proste, \( p=q \)}{}{}}
\LANGcs{
\Question{
Vyšetřete vzájemnou polohu přímek \( p\colon 6x+4y+8=0 \) a \( q\colon y=-\frac32 x+2 \).}
{různé rovnoběžky, \( p\parallel q;\ p\neq q \)}{různoběžky, \( p\cap q=\left\{\left[0;-2\right]\right\} \)}{různoběžky, \( p\cap q=\left\{\left[0;2\right]\right\} \)}{totožné přímky, \( p=q \)}{}{}}
\LANGsk{
\Question{
Určte vzájomnú polohu priamok \( p\colon 6x+4y+8=0 \) a \( q\colon y=-\frac32 x+2 \).}
{rovnobežné rôzne, \( p\parallel q;\ p\neq q \)}{rôznobežné, \( p\cap q=\left\{\left[0;-2\right]\right\} \)}{rôznobežné, \( p\cap q=\left\{\left[0;2\right]\right\} \)}{totožné, \( p=q \)}{}{}}
\LANGes{
\Question{
Determina la posición relativa de las rectas \( p\colon 6x+4y+8=0 \) y \( q\colon y=-\frac32 x+2 \).}
{rectas secantes, \( p\parallel q;\ p\neq q \)}{rectas secantes, \( p\cap q=\left\{\left[0;-2\right]\right\} \)}{rectas secantes, \( p\cap q=\left\{\left[0;2\right]\right\} \)}{rectas coincidentes, \( p=q \)}{}{}}
\End

\Data\ID{2010014205}
\Updated{2022-07-20 08:07:08}
\Level{A}
\SubArea{Analytic geometry in a plane}
\IMG{2010014201-figure0.pdf}
\LANGen{
\Question{
A straight line \( p \) is given by the point \( A \) and the normal vector \( \vec{n} \) (see the picture). Find its equation in general form.}
{\( p\colon 2x+5y-4=0 \)}{\( p\colon 5x-2y=0 \)}{\( p\colon 5x-2y-10=0 \)}{\( p\colon 2x+5y+33=0 \)}{}{}}
\LANGpl{
\Question{
Prostą \( p \) wyznacza punkt \( A \) i wektor normalny \( \vec{n} \) (patrz rysunek). Znajdź jej równanie w postaci ogólnej.}
{\( p\colon 2x+5y-4=0 \)}{\( p\colon 5x-2y=0 \)}{\( p\colon 5x-2y-10=0 \)}{\( p\colon 2x+5y+33=0 \)}{}{}}
\LANGcs{
\Question{
Přímka \( p \) je dána bodem \( A \) a normálovým vektorem \( \vec{n} \) (viz obrázek). Určete její obecnou rovnici.}
{\( p\colon 2x+5y-4=0 \)}{\( p\colon 5x-2y=0 \)}{\( p\colon 5x-2y-10=0 \)}{\( p\colon 2x+5y+33=0 \)}{}{}}
\LANGsk{
\Question{
Priamka \( p \) je daná bodom \( A \) a normálovým vektorom \( \vec{n} \) (pozri obrázok). Určte jej všeobecnú rovnicu.}
{\( p\colon 2x+5y-4=0 \)}{\( p\colon 5x-2y=0 \)}{\( p\colon 5x-2y-10=0 \)}{\( p\colon 2x+5y+33=0 \)}{}{}}
\LANGes{
\Question{
La recta \( p\) viene dada por el punto \( A\) y el vector normal  \( \vec{n} \) (observa la imagen). Determina su ecuación general.}
{\( p\colon 2x+5y-4=0 \)}{\( p\colon 5x-2y=0 \)}{\( p\colon 5x-2y-10=0 \)}{\( p\colon 2x+5y+33=0 \)}{}{}}
\End

\Data\ID{2010014209}
\Updated{2022-05-16 11:05:40}
\Level{A}
\SubArea{Analytic geometry in a plane}
\LANGen{
\Question{
In the following list identify a vector having the same direction as the line passing through
the points \(A\) 
and \(B\).
\[ 
 A = \left [4;1\right ],\ \qquad B = \left [3;2\right ]
\]}
{\(\left (-1;1\right )\)}{\(\left (1;1\right )\)}{\(\left (7;3\right )\)}{\(\left (5;5\right )\)}{}{}}
\LANGpl{
\Question{
Z poniższej listy wybierz  wektor mający ten sam kierunek, co prostą rzechodząca przez punkty \(A\)
oraz \(B\).
\[ 
 A = \left [4;1\right ],\ \qquad B = \left [3;2\right ]
\]}
{\(\left (-1;1\right )\)}{\(\left (1;1\right )\)}{\(\left (7;3\right )\)}{\(\left (5;5\right )\)}{}{}}
\LANGcs{
\Question{
Z nabízených možností vyberte směrový vektor přímky, která prochází body \(
 A = \left [4;1\right ],\ B = \left [3;2\right ].
\)}
{\(\left (-1;1\right )\)}{\(\left (1;1\right )\)}{\(\left (7;3\right )\)}{\(\left (5;5\right )\)}{}{}}
\LANGsk{
\Question{
Z nasledujúcich možností vyberte smerový vektor priamky, ktorá prechádza bodmi \(A\) 
a \(B\).
\[ 
 A = \left [4;1\right ],\ \qquad B = \left [3;2\right ]
\]}
{\(\left (-1;1\right )\)}{\(\left (1;1\right )\)}{\(\left (7;3\right )\)}{\(\left (5;5\right )\)}{}{}}
\LANGes{
\Question{
En la siguiente lista identifica un vector que tiene la misma dirección que la recta que pasa por los puntos \(A\) 
y \(B\).
\[ 
 A = \left [4;1\right ],\ \qquad B = \left [3;2\right ]
\]}
{\(\left (-1;1\right )\)}{\(\left (1;1\right )\)}{\(\left (7;3\right )\)}{\(\left (5;5\right )\)}{}{}}
\End

\Data\ID{2010014210}
\Updated{2022-05-16 11:05:40}
\Level{A}
\SubArea{Analytic geometry in a plane}
\LANGen{
\Question{
In the following list find the vector having the same direction as the line
\(p\).
\[ 
 p\colon 2x +3y + 1 = 0
\]}
{\(\left (-3;2\right )\)}{\(\left (2;3\right )\)}{\(\left (2;-3\right )\)}{\(\left (3;1\right )\)}{}{}}
\LANGpl{
\Question{
Z poniższej listy wybierz wektor mający ten sam kierunek co prostą
\(p\).
\[ 
 p\colon 2x +3y + 1 = 0
\]}
{\(\left (-3;2\right )\)}{\(\left (2;3\right )\)}{\(\left (2;-3\right )\)}{\(\left (3;1\right )\)}{}{}}
\LANGcs{
\Question{
Z nabízených možností vyberte vektor, který je rovnoběžný se směrovým vektorem přímky \(p\).
\[ 
 p\colon 2x +3y + 1 = 0
\]}
{\(\left (-3;2\right )\)}{\(\left (2;3\right )\)}{\(\left (2;-3\right )\)}{\(\left (3;1\right )\)}{}{}}
\LANGsk{
\Question{
Z ponúkaných možností vyberte vektor, ktorý je rovnobežný so smerovým vektorom priamky \(p\).
\[ 
 p\colon 2x +3y + 1 = 0
\]}
{\(\left (-3;2\right )\)}{\(\left (2;3\right )\)}{\(\left (2;-3\right )\)}{\(\left (3;1\right )\)}{}{}}
\LANGes{
\Question{
En la siguiente lista identifica un vector que tiene la misma dirección que la recta
\(p\).
\[ 
 p\colon 2x +3y + 1 = 0
\]}
{\(\left (-3;2\right )\)}{\(\left (2;3\right )\)}{\(\left (2;-3\right )\)}{\(\left (3;1\right )\)}{}{}}
\End

\Data\ID{2010014401}
\Updated{2022-05-16 11:05:40}
\Level{A}
\SubArea{Analytic geometry in a plane}
\LANGen{
\Question{
Find the perpendicular line to the line \(3x-2y+1=0\).}
{\(2x+3y-4=0\)}{\(6x-4y+2=0\)}{\(2x-3y+1=0\)}{\(-2x+3y-3=0\)}{}{}}
\LANGpl{
\Question{
Znajdź prostą prostopadłą do prostej \(3x-2y+1=0\).}
{\(2x+3y-4=0\)}{\(6x-4y+2=0\)}{\(2x-3y+1=0\)}{\(-2x+3y-3=0\)}{}{}}
\LANGcs{
\Question{
Určete přímku kolmou k přímce \(3x-2y+1=0\).}
{\(2x+3y-4=0\)}{\(6x-4y+2=0\)}{\(2x-3y+1=0\)}{\(-2x+3y-3=0\)}{}{}}
\LANGsk{
\Question{
Určte priamku, ktorá je kolmá na priamku \(3x-2y+1=0\).}
{\(2x+3y-4=0\)}{\(6x-4y+2=0\)}{\(2x-3y+1=0\)}{\(-2x+3y-3=0\)}{}{}}
\LANGes{
\Question{
Determina la recta perpendicular a la recta \(3x-2y+1=0\).}
{\(2x+3y-4=0\)}{\(6x-4y+2=0\)}{\(2x-3y+1=0\)}{\(-2x+3y-3=0\)}{}{}}
\End

\Data\ID{2010014402}
\Updated{2022-05-16 11:05:40}
\Level{A}
\SubArea{Analytic geometry in a plane}
\LANGen{
\Question{
Find the perpendicular line to the line \(2x-3y+1=0\).}
{\(3x+2y-4=0\)}{\(4x-6y+2=0\)}{\(3x-2y+1=0\)}{\(-2x+3y-1=0\)}{}{}}
\LANGpl{
\Question{
Znajdź prostą prostopadłą do prostej \(2x-3y+1=0\).}
{\(3x+2y-4=0\)}{\(4x-6y+2=0\)}{\(3x-2y+1=0\)}{\(-2x+3y-1=0\)}{}{}}
\LANGcs{
\Question{
Určete přímku kolmou k přímce \(2x-3y+1=0\).}
{\(3x+2y-4=0\)}{\(4x-6y+2=0\)}{\(3x-2y+1=0\)}{\(-2x+3y-1=0\)}{}{}}
\LANGsk{
\Question{
Určte priamku, ktorá je kolmá na priamku \(2x-3y+1=0\).}
{\(3x+2y-4=0\)}{\(4x-6y+2=0\)}{\(3x-2y+1=0\)}{\(-2x+3y-1=0\)}{}{}}
\LANGes{
\Question{
Determina la recta perpendicular a la recta \(2x-3y+1=0\).}
{\(3x+2y-4=0\)}{\(4x-6y+2=0\)}{\(3x-2y+1=0\)}{\(-2x+3y-1=0\)}{}{}}
\End

\Data\ID{2010014403}
\Updated{2022-05-16 11:05:40}
\Level{A}
\SubArea{Analytic geometry in a plane}
\LANGen{
\Question{
Find the parallel line to the line \(p\colon -x+3y-5=0\).}
{\(2x-6y-5=0\)}{\(3x+y+2=0\)}{\(-2x-6y+1=0\)}{\(-x-3y-3=0\)}{}{}}
\LANGpl{
\Question{
Znajdź prostą równoległą do prostej \(p\colon -x+3y-5=0\).}
{\(2x-6y-5=0\)}{\(3x+y+2=0\)}{\(-2x-6y+1=0\)}{\(-x-3y-3=0\)}{}{}}
\LANGcs{
\Question{
Určete přímku rovnoběžnou s přímkou \(p\colon -x+3y-5=0\).}
{\(2x-6y-5=0\)}{\(3x+y+2=0\)}{\(-2x-6y+1=0\)}{\(-x-3y-3=0\)}{}{}}
\LANGsk{
\Question{
Určte priamku, ktorá je rovnobežná s priamkou \(p\colon -x+3y-5=0\).}
{\(2x-6y-5=0\)}{\(3x+y+2=0\)}{\(-2x-6y+1=0\)}{\(-x-3y-3=0\)}{}{}}
\LANGes{
\Question{
Determina la recta paralela a la recta \(p\colon -x+3y-5=0\).}
{\(2x-6y-5=0\)}{\(3x+y+2=0\)}{\(-2x-6y+1=0\)}{\(-x-3y-3=0\)}{}{}}
\End

\Data\ID{2010014404}
\Updated{2022-05-16 11:05:40}
\Level{A}
\SubArea{Analytic geometry in a plane}
\LANGen{
\Question{
Find the parallel line to the line \(p\colon -3x+y-5=0\).}
{\(6x-2y-5=0\)}{\(x+3y-5=0\)}{\(-6x-2y+1=0\)}{\(-x-3y-1=0\)}{}{}}
\LANGpl{
\Question{
Wyznacz prostą równoległą do prostej \(p\colon -3x+y-5=0\).}
{\(6x-2y-5=0\)}{\(x+3y-5=0\)}{\(-6x-2y+1=0\)}{\(-x-3y-1=0\)}{}{}}
\LANGcs{
\Question{
Určete přímku rovnoběžnou s přímkou  \(p\colon -3x+y-5=0\).}
{\(6x-2y-5=0\)}{\(x+3y-5=0\)}{\(-6x-2y+1=0\)}{\(-x-3y-1=0\)}{}{}}
\LANGsk{
\Question{
Určte priamku, ktorá je rovnobežná s priamkou \(p\colon -3x+y-5=0\).}
{\(6x-2y-5=0\)}{\(x+3y-5=0\)}{\(-6x-2y+1=0\)}{\(-x-3y-1=0\)}{}{}}
\LANGes{
\Question{
Determina la recta paralela a la recta \(p\colon -3x+y-5=0\).}
{\(6x-2y-5=0\)}{\(x+3y-5=0\)}{\(-6x-2y+1=0\)}{\(-x-3y-1=0\)}{}{}}
\End

\Data\ID{2010014405}
\Updated{2022-05-16 11:05:44}
\Level{A}
\SubArea{Analytic geometry in a plane}
\LANGen{
\Question{
Find the perpendicular line to the line \(p\colon x=1+5t\), \(y=-3-2t\), \(t \in \mathbb{R}\).}
{\(-5x+2y+1=0\)}{\(2x+5y-3=0\)}{\(-4x+10y-1=0\)}{\(2x-5y-3=0\)}{}{}}
\LANGpl{
\Question{
Wyznacz prostą prostopadłą do prostej \(p\colon x=1+5t\), \(y=-3-2t\), \(t \in \mathbb{R}\).}
{\(-5x+2y+1=0\)}{\(2x+5y-3=0\)}{\(-4x+10y-1=0\)}{\(2x-5y-3=0\)}{}{}}
\LANGcs{
\Question{
Určete přímku kolmou k přímce \(p\colon x=1+5t\), \(y=-3-2t\), \(t \in \mathbb{R}\).}
{\(-5x+2y+1=0\)}{\(2x+5y-3=0\)}{\(-4x+10y-1=0\)}{\(2x-5y-3=0\)}{}{}}
\LANGsk{
\Question{
Určte priamku, ktorá je kolmá na priamku \(p\colon x=1+5t\), \(y=-3-2t\), \(t \in \mathbb{R}\).}
{\(-5x+2y+1=0\)}{\(2x+5y-3=0\)}{\(-4x+10y-1=0\)}{\(2x-5y-3=0\)}{}{}}
\LANGes{
\Question{
Determina la recta perpendicular a la recta \(p\colon x=1+5t\), \(y=-3-2t\), \(t \in \mathbb{R}\).}
{\(-5x+2y+1=0\)}{\(2x+5y-3=0\)}{\(-4x+10y-1=0\)}{\(2x-5y-3=0\)}{}{}}
\End

\Data\ID{2010014406}
\Updated{2022-05-16 11:05:44}
\Level{A}
\SubArea{Analytic geometry in a plane}
\LANGen{
\Question{
Find the perpendicular line to the line \(p\colon x=1-2t\), \(y=-3+5t\), \(t \in \mathbb{R}\).}
{\(2x-5y+1=0\).}{\(5x+2y-1=0\).}{\(4x+10y+4=0\).}{\(x-3y+2=0\).}{}{}}
\LANGpl{
\Question{
Znajdź prostą prostopadłą do prostej \(p\colon x=1-2t\), \(y=-3+5t\), \(t \in \mathbb{R}\).}
{\(2x-5y+1=0\).}{\(5x+2y-1=0\).}{\(4x+10y+4=0\).}{\(x-3y+2=0\).}{}{}}
\LANGcs{
\Question{
Určete přímku kolmou k přímce \(p\colon x=1-2t\), \(y=-3+5t\), \(t \in \mathbb{R}\).}
{\(2x-5y+1=0\).}{\(5x+2y-1=0\).}{\(4x+10y+4=0\).}{\(x-3y+2=0\).}{}{}}
\LANGsk{
\Question{
Určte priamku, ktorá je kolmá na priamku \(p\colon x=1-2t\), \(y=-3+5t\), \(t \in \mathbb{R}\).}
{\(2x-5y+1=0\).}{\(5x+2y-1=0\).}{\(4x+10y+4=0\).}{\(x-3y+2=0\).}{}{}}
\LANGes{
\Question{
Determina la recta perpendicular a la recta \(p\colon x=1-2t\), \(y=-3+5t\), \(t \in \mathbb{R}\).}
{\(2x-5y+1=0\).}{\(5x+2y-1=0\).}{\(4x+10y+4=0\).}{\(x-3y+2=0\).}{}{}}
\End

\Data\ID{2010014407}
\Updated{2022-05-16 11:05:44}
\Level{A}
\SubArea{Analytic geometry in a plane}
\LANGen{
\Question{
Find the parallel line to the line \(p\colon x=-2-t\), \(y=1+3t\), \(t \in \mathbb{R}\).}
{\(-3x-y+2=0\)}{\(-2x+6y+1=0\)}{\(x-3y-2=0\)}{\(x+3y+5=0\)}{}{}}
\LANGpl{
\Question{
Znajdź prostą równoległą do prostej \(p\colon x=-2-t\), \(y=1+3t\), \(t \in \mathbb{R}\).}
{\(-3x-y+2=0\)}{\(-2x+6y+1=0\)}{\(x-3y-2=0\)}{\(x+3y+5=0\)}{}{}}
\LANGcs{
\Question{
Určete přímku rovnoběžnou s přímkou \(p\colon x=-2-t\), \(y=1+3t\), \(t \in \mathbb{R}\).}
{\(-3x-y+2=0\)}{\(-2x+6y+1=0\)}{\(x-3y-2=0\)}{\(x+3y+5=0\)}{}{}}
\LANGsk{
\Question{
Určte priamku, ktorá je rovnobežná s priamkou \(p\colon x=-2-t\), \(y=1+3t\), \(t \in \mathbb{R}\).}
{\(-3x-y+2=0\)}{\(-2x+6y+1=0\)}{\(x-3y-2=0\)}{\(x+3y+5=0\)}{}{}}
\LANGes{
\Question{
Determina la recta paralela a la recta \(p\colon x=-2-t\), \(y=1+3t\), \(t \in \mathbb{R}\).}
{\(-3x-y+2=0\)}{\(-2x+6y+1=0\)}{\(x-3y-2=0\)}{\(x+3y+5=0\)}{}{}}
\End

\Data\ID{2010014408}
\Updated{2022-05-16 11:05:44}
\Level{A}
\SubArea{Analytic geometry in a plane}
\LANGen{
\Question{
Find the parallel line to the line \(p\colon x=-2+3t\), \(y=1-t\), \(t \in \mathbb{R}\).}
{\(-x-3y+2=0\)}{\(6x-2y+1=0\)}{\(3x+y-2=0\)}{\(-2x+y-1=0\)}{}{}}
\LANGpl{
\Question{
Znajdź prostą równoległą do prostej \(p\colon x=-2+3t\), \(y=1-t\), \(t \in \mathbb{R}\).}
{\(-x-3y+2=0\)}{\(6x-2y+1=0\)}{\(3x+y-2=0\)}{\(-2x+y-1=0\)}{}{}}
\LANGcs{
\Question{
Určete přímku rovnoběžnou s přímkou \(p\colon x=-2+3t\), \(y=1-t\), \(t \in \mathbb{R}\).}
{\(-x-3y+2=0\)}{\(6x-2y+1=0\)}{\(3x+y-2=0\)}{\(-2x+y-1=0\)}{}{}}
\LANGsk{
\Question{
Určte priamku, ktorá je rovnobežná s priamkou \(p\colon x=-2+3t\), \(y=1-t\), \(t \in \mathbb{R}\).}
{\(-x-3y+2=0\)}{\(6x-2y+1=0\)}{\(3x+y-2=0\)}{\(-2x+y-1=0\)}{}{}}
\LANGes{
\Question{
Determina la recta paralela a la recta \(p\colon x=-2+3t\), \(y=1-t\), \(t \in \mathbb{R}\).}
{\(-x-3y+2=0\)}{\(6x-2y+1=0\)}{\(3x+y-2=0\)}{\(-2x+y-1=0\)}{}{}}
\End

\Data\ID{2010014409}
\Updated{2022-05-16 11:05:44}
\Level{A}
\SubArea{Analytic geometry in a plane}
\LANGen{
\Question{
Find the number \(c\) such that the point \(C=[c; 7]\) lies on the line \(p\colon 3x-4y+1= 0\).}
{\(9\)}{\(5\)}{\(7\)}{\(-1\)}{}{}}
\LANGpl{
\Question{
Znajdź liczbę \(c\) taką, że punkt \(C=[c; 7]\) leży na prostej \(p\colon 3x-4y+1= 0\).}
{\(9\)}{\(5\)}{\(7\)}{\(-1\)}{}{}}
\LANGcs{
\Question{
Určete reálné číslo \(c\) tak, aby bod \(C=[c; 7]\) ležel na přímce \(p\colon 3x-4y+1= 0\).}
{\(9\)}{\(5\)}{\(7\)}{\(-1\)}{}{}}
\LANGsk{
\Question{
Určte reálné číslo \(c\) tak, aby bod \(C=[c; 7]\) ležal na priamke \(p\colon 3x-4y+1= 0\).}
{\(9\)}{\(5\)}{\(7\)}{\(-1\)}{}{}}
\LANGes{
\Question{
Halla el número \(c\) suponiendo que el punto \(C=[c; 7]\) se encuentra en la recta \(p\colon 3x-4y+1= 0\).}
{\(9\)}{\(5\)}{\(7\)}{\(-1\)}{}{}}
\End

\Data\ID{2010014410}
\Updated{2022-05-16 11:05:44}
\Level{A}
\SubArea{Analytic geometry in a plane}
\LANGen{
\Question{
Find the number \(c\) such that the point \(C=[4;c]\) lies on the line \(p\colon 4x-3y-1 = 0\).}
{\(5\)}{\(9\)}{\(3\)}{\(-1\)}{}{}}
\LANGpl{
\Question{
Znajdź liczbę \(c\) taką, że punkt \(C=[4;c]\) leży na prostej \(p\colon 4x-3y-1 = 0\).}
{\(5\)}{\(9\)}{\(3\)}{\(-1\)}{}{}}
\LANGcs{
\Question{
Určete reálné číslo \(c\) tak, aby bod \(C=[4;c]\) ležel na přímce \(p\colon 4x-3y-1 = 0\).}
{\(5\)}{\(9\)}{\(3\)}{\(-1\)}{}{}}
\LANGsk{
\Question{
Určte reálne číslo \(c\) tak, aby bod \(C=[4;c]\) ležal na priamke \(p\colon 4x-3y-1 = 0\).}
{\(5\)}{\(9\)}{\(3\)}{\(-1\)}{}{}}
\LANGes{
\Question{
Halla el número \(c\) suponiendo que el punto \(C=[4;c]\) se encuentra en la recta \(p\colon 4x-3y-1 = 0\).}
{\(5\)}{\(9\)}{\(3\)}{\(-1\)}{}{}}
\End

\Data\ID{2010014411}
\Updated{2022-05-16 11:05:44}
\Level{A}
\SubArea{Analytic geometry in a plane}
\LANGen{
\Question{
Find the number \(c\) such that the point \(C=[c;9]\) lies on the line \(p\colon x=2+3t\), \(y=1+4t\), \(t \in \mathbb{R}\).}
{\(8\)}{\(3\)}{\(7\)}{\(-4\)}{}{}}
\LANGpl{
\Question{
Znajdź liczbę \(c\) taką, że punkt \(C=[c;9]\) leży na prostej \(p\colon x=2+3t\), \(y=1+4t\), \(t \in \mathbb{R}\).}
{\(8\)}{\(3\)}{\(7\)}{\(-4\)}{}{}}
\LANGcs{
\Question{
Určete reálné číslo \(c\) tak, aby bod \(C=[c;9]\) ležel na přímce \(p\colon x=2+3t\), \(y=1+4t\), \(t \in \mathbb{R}\).}
{\(8\)}{\(3\)}{\(7\)}{\(-4\)}{}{}}
\LANGsk{
\Question{
Určte reálné číslo \(c\) tak, aby bod \(C=[c;9]\) ležal na priamke \(p\colon x=2+3t\), \(y=1+4t\), \(t \in \mathbb{R}\).}
{\(8\)}{\(3\)}{\(7\)}{\(-4\)}{}{}}
\LANGes{
\Question{
Halla el número \(c\)  suponiendo que el punto \(C=[c;9]\) se encuentra en la recta \(p\colon x=2+3t\), \(y=1+4t\), \(t \in \mathbb{R}\).}
{\(8\)}{\(3\)}{\(7\)}{\(-4\)}{}{}}
\End

\Data\ID{2010014412}
\Updated{2022-05-16 11:05:44}
\Level{A}
\SubArea{Analytic geometry in a plane}
\LANGen{
\Question{
Find the number \(c\) such that the point \(C=[5;c]\) lies on the line \(p\colon x=2+3t\), \(y=1+4t\), \(t \in \mathbb{R}\).}
{\(5\)}{\(1\)}{\(-1\)}{\(2\)}{}{}}
\LANGpl{
\Question{
Znajdź liczbę \(c\) taką, że punkt \(C=[5;c]\) leży na prostej \(p\colon x=2+3t\), \(y=1+4t\), \(t \in \mathbb{R}\).}
{\(5\)}{\(1\)}{\(-1\)}{\(2\)}{}{}}
\LANGcs{
\Question{
Určete reálné číslo \(c\) tak, aby bod \(C=[5;c]\) ležel na přímce \(p\colon x=2+3t\), \(y=1+4t\), \(t \in \mathbb{R}\).}
{\(5\)}{\(1\)}{\(-1\)}{\(2\)}{}{}}
\LANGsk{
\Question{
Určte reálné číslo \(c\) tak, aby bod \(C=[5;c]\) ležal na priamke \(p\colon x=2+3t\), \(y=1+4t\), \(t \in \mathbb{R}\).}
{\(5\)}{\(1\)}{\(-1\)}{\(2\)}{}{}}
\LANGes{
\Question{
Halla el número \(c\)  suponiendo que el punto \(C=[5;c]\) se encuentra en la recta \(p\colon x=2+3t\), \(y=1+4t\), \(t \in \mathbb{R}\).}
{\(5\)}{\(1\)}{\(-1\)}{\(2\)}{}{}}
\End

\Data\ID{2010014413}
\Updated{2022-05-16 11:05:44}
\Level{A}
\SubArea{Analytic geometry in a plane}
\LANGen{
\Question{
Find the non-zero coordinate of the intersection point of the line \(p\colon -4x+3y-1=0\) and the \(x\)-axis.}
{\(-\frac14\)}{\(5\)}{\(-3\)}{\(\frac13\)}{}{}}
\LANGpl{
\Question{
Znajdź niezerową współrzędną punktu przecięcia linii \(p\colon -4x+3y-1=0\) i osi \(x\).}
{\(-\frac14\)}{\(5\)}{\(-3\)}{\(\frac13\)}{}{}}
\LANGcs{
\Question{
Určete nenulovou souřadnici průsečíku přímky \(p\colon -4x+3y-1=0\) a osy \(x\).}
{\(-\frac14\)}{\(5\)}{\(-3\)}{\(\frac13\)}{}{}}
\LANGsk{
\Question{
Určte nenulovú súradnicu priesečníku priamky \(p\colon -4x+3y-1=0\) a osi \(x\).}
{\(-\frac14\)}{\(5\)}{\(-3\)}{\(\frac13\)}{}{}}
\LANGes{
\Question{
Determina la coordenada distinta de cero del punto de intersección de la recta \(p\colon -4x+3y-1=0\) y el eje \(x\).}
{\(-\frac14\)}{\(5\)}{\(-3\)}{\(\frac13\)}{}{}}
\End

\Data\ID{2010014414}
\Updated{2022-05-16 11:05:44}
\Level{A}
\SubArea{Analytic geometry in a plane}
\LANGen{
\Question{
Find the non-zero coordinate of the intersection point of the line \(p\colon -x+y-1=0\) and the \(y\)-axis.}
{\(1\)}{\(-1\)}{\(0\)}{\(2\)}{}{}}
\LANGpl{
\Question{
Wyznacz niezerową współrzędną punktu przecięcia linii \(p\colon -x+y-1=0\) i osi \(y\).}
{\(1\)}{\(-1\)}{\(0\)}{\(2\)}{}{}}
\LANGcs{
\Question{
Určete nenulovou souřadnici průsečíku přímky \(p\colon -x+y-1=0\) a osy \(y\).}
{\(1\)}{\(-1\)}{\(0\)}{\(2\)}{}{}}
\LANGsk{
\Question{
Určte nenulovú súradnicu priesečníku priamky \(p\colon -x+y-1=0\) a osi \(y\).}
{\(1\)}{\(-1\)}{\(0\)}{\(2\)}{}{}}
\LANGes{
\Question{
Determina la coordenada distinta de cero del punto de intersección de la recta \(p\colon -x+y-1=0\) y el eje \(y\).}
{\(1\)}{\(-1\)}{\(0\)}{\(2\)}{}{}}
\End

\Data\ID{2010014415}
\Updated{2022-05-16 11:05:49}
\Level{A}
\SubArea{Analytic geometry in a plane}
\LANGen{
\Question{
Find the parallel line to the line \(p\colon x-2y-3=0\), which goes through the point \(M=[1;1]\).}
{\(x-2y+1=0\)}{\(2x-y-1=0\)}{\(2x+y-3=0\)}{\(2x-4y-3=0\)}{}{}}
\LANGpl{
\Question{
Znajdź prostą równoległą do prostej \(p\colon x-2y-3=0\), która przechodzi przez punkt \(M=[1;1]\).}
{\(x-2y+1=0\)}{\(2x-y-1=0\)}{\(2x+y-3=0\)}{\(2x-4y-3=0\)}{}{}}
\LANGcs{
\Question{
Určete přímku rovnoběžnou s přímkou \(p\colon x-2y-3=0\), která prochází bodem \(M=[1;1]\).}
{\(x-2y+1=0\)}{\(2x-y-1=0\)}{\(2x+y-3=0\)}{\(2x-4y-3=0\)}{}{}}
\LANGsk{
\Question{
Určte priamku rovnobežnú s priamkou \(p\colon x-2y-3=0\), ktorá prechádza bodom \(M=[1;1]\).}
{\(x-2y+1=0\)}{\(2x-y-1=0\)}{\(2x+y-3=0\)}{\(2x-4y-3=0\)}{}{}}
\LANGes{
\Question{
Determina la recta paralela a la recta \(p\colon x-2y-3=0\), que pasa por el punto \(M=[1;1]\).}
{\(x-2y+1=0\)}{\(2x-y-1=0\)}{\(2x+y-3=0\)}{\(2x-4y-3=0\)}{}{}}
\End

\Data\ID{2010014416}
\Updated{2022-05-16 11:05:49}
\Level{A}
\SubArea{Analytic geometry in a plane}
\LANGen{
\Question{
Find the perpendicular line to the line \(p\colon 3x-y+2=0\), which goes through the point \(M=[-1;1]\).}
{\(x+3y-2=0\)}{\(x+3y+2=0\)}{\(-x+3y-2=0\)}{\(x-3y+1=0\)}{}{}}
\LANGpl{
\Question{
Znajdź prostą prostopadłą do prostej \(p\colon 3x-y+2=0\), która przechodzi przez punkt \(M=[-1;1]\).}
{\(x+3y-2=0\)}{\(x+3y+2=0\)}{\(-x+3y-2=0\)}{\(x-3y+1=0\)}{}{}}
\LANGcs{
\Question{
Určete přímku kolmou k přímce \(p\colon 3x-y+2=0\), která prochází bodem \(M=[-1;1]\).}
{\(x+3y-2=0\)}{\(x+3y+2=0\)}{\(-x+3y-2=0\)}{\(x-3y+1=0\)}{}{}}
\LANGsk{
\Question{
Určte priamku kolmú na priamku \(p\colon 3x-y+2=0\), ktorá prechádza bodom \(M=[-1;1]\).}
{\(x+3y-2=0\)}{\(x+3y+2=0\)}{\(-x+3y-2=0\)}{\(x-3y+1=0\)}{}{}}
\LANGes{
\Question{
Determina la recta perpendicular a la recta \(p\colon 3x-y+2=0\), que pasa por el punto \(M=[-1;1]\).}
{\(x+3y-2=0\)}{\(x+3y+2=0\)}{\(-x+3y-2=0\)}{\(x-3y+1=0\)}{}{}}
\End

\Data\ID{2010014601}
\Updated{2022-05-16 11:05:49}
\Level{A}
\SubArea{Analytic geometry in a plane}
\LANGen{
\Question{
Find the normal vector of the line passing through the points
\(A = [1;3]\) and
\(B = [-2;5]\).}
{\((2;3)\)}{\((-3;2)\)}{\((3;-2)\)}{\((2;-3)\)}{}{}}
\LANGpl{
\Question{
Znajdź wektor normalny prostej przechodzącej przez punkty
\(A = [1;3]\) i
\(B = [-2;5]\).}
{\((2;3)\)}{\((-3;2)\)}{\((3;-2)\)}{\((2;-3)\)}{}{}}
\LANGcs{
\Question{
Z nabízených možností vyberte normálový vektor přímky, která prochází body
\(A = [1;3]\) a
\(B = [-2;5]\).}
{\((2;3)\)}{\((-3;2)\)}{\((3;-2)\)}{\((2;-3)\)}{}{}}
\LANGsk{
\Question{
Určte normálový vektor priamky, ktorá prechádza bodmi
\(A = [1;3]\) a
\(B = [-2;5]\).}
{\((2;3)\)}{\((-3;2)\)}{\((3;-2)\)}{\((2;-3)\)}{}{}}
\LANGes{
\Question{
Halla un vector normal de la recta que pasa por los puntos
\(A = [1;3]\) y
\(B = [-2;5]\).}
{\((2;3)\)}{\((-3;2)\)}{\((3;-2)\)}{\((2;-3)\)}{}{}}
\End

\Data\ID{2010014602}
\Updated{2022-05-16 11:05:49}
\Level{A}
\SubArea{Analytic geometry in a plane}
\LANGen{
\Question{
Find the normal vector of the following line. 
\[ 
p\colon \begin{aligned}[t] x =&1 +4t, &
\\y =& - 3 -2t;\ t\in \mathbb{R}
\\ \end{aligned}
\]}
{\((1;2)\)}{\((4;-2)\)}{\((1;-3)\)}{\((-2;1)\)}{}{}}
\LANGpl{
\Question{
Znajdź wektor normalny następującej linii.
\[ 
p\colon \begin{aligned}[t] x =&1 +4t, &
\\y =& - 3 -2t;\ t\in \mathbb{R}
\\ \end{aligned}
\]}
{\((1;2)\)}{\((4;-2)\)}{\((1;-3)\)}{\((-2;1)\)}{}{}}
\LANGcs{
\Question{
Z nabízených možností vyberte normálový vektor přímky, která je vyjádřena parametrickými rovnicemi:
\[ 
p\colon \begin{aligned}[t] x =&1 +4t, &
\\y =& - 3 -2t;\ t\in \mathbb{R}
\\ \end{aligned}
\]}
{\((1;2)\)}{\((4;-2)\)}{\((1;-3)\)}{\((-2;1)\)}{}{}}
\LANGsk{
\Question{
Určte normálový vektor priamky, ktorá je vyjadrená parametricky:
\[ 
p\colon \begin{aligned}[t] x =&1 +4t, &
\\y =& - 3 -2t;\ t\in \mathbb{R}
\\ \end{aligned}
\]}
{\((1;2)\)}{\((4;-2)\)}{\((1;-3)\)}{\((-2;1)\)}{}{}}
\LANGes{
\Question{
Halla un vector normal de la recta
\[ 
p\colon \begin{aligned}[t] x =&1 +4t, &
\\y =& - 3 -2t;\ t\in \mathbb{R}
\\ \end{aligned}
\]}
{\((1;2)\)}{\((4;-2)\)}{\((1;-3)\)}{\((-2;1)\)}{}{}}
\End

\Data\ID{2010014603}
\Updated{2022-05-16 11:05:49}
\Level{A}
\SubArea{Analytic geometry in a plane}
\LANGen{
\Question{
In the following list identify a line which is perpendicular to the line
\( 2x +3y -7= 0\).}
{\(\begin{aligned}[t]  x& = 2t, &
\\y & = -11+3t;\  t\in \mathbb{R}
\\ \end{aligned}\)}{\(\begin{aligned}[t]  x& = 1+3t, &
\\y & = 11 - 2t;\  t\in \mathbb{R}
\\ \end{aligned}\)}{\(\begin{aligned}[t]  x& = 2+t, &
\\y & = 3 - t;\  t\in \mathbb{R}
\\ \end{aligned}\)}{\(\begin{aligned}[t]  x& = 2t+7, &
\\y & = - 3t+1;\  t\in \mathbb{R}
\\ \end{aligned}\)}{}{}}
\LANGpl{
\Question{
Z poniższej listy wybierz  prostą prostopadłą do prostej
\( 2x +3y -7= 0\).}
{\(\begin{aligned}[t]  x& = 2t, &
\\y & = -11+3t;\  t\in \mathbb{R}
\\ \end{aligned}\)}{\(\begin{aligned}[t]  x& = 1+3t, &
\\y & = 11 - 2t;\  t\in \mathbb{R}
\\ \end{aligned}\)}{\(\begin{aligned}[t]  x& = 2+t, &
\\y & = 3 - t;\  t\in \mathbb{R}
\\ \end{aligned}\)}{\(\begin{aligned}[t]  x& = 2t+7, &
\\y & = - 3t+1;\  t\in \mathbb{R}
\\ \end{aligned}\)}{}{}}
\LANGcs{
\Question{
Z následujících přímek zadaných parametricky vyberte tu, která je kolmá k přímce
\( 2x +3y -7= 0\).}
{\(\begin{aligned}[t]  x& = 2t, &
\\y & = -11+3t;\  t\in \mathbb{R}
\\ \end{aligned}\)}{\(\begin{aligned}[t]  x& = 1+3t, &
\\y & = 11 - 2t;\  t\in \mathbb{R}
\\ \end{aligned}\)}{\(\begin{aligned}[t]  x& = 2+t, &
\\y & = 3 - t;\  t\in \mathbb{R}
\\ \end{aligned}\)}{\(\begin{aligned}[t]  x& = 2t+7, &
\\y & = - 3t+1;\  t\in \mathbb{R}
\\ \end{aligned}\)}{}{}}
\LANGsk{
\Question{
Z následujúcich priamok zadaných parametricky vyberte tú, ktorá je kolmá na priamku
\( 2x +3y -7= 0\).}
{\(\begin{aligned}[t]  x& = 2t, &
\\y & = -11+3t;\  t\in \mathbb{R}
\\ \end{aligned}\)}{\(\begin{aligned}[t]  x& = 1+3t, &
\\y & = 11 - 2t;\  t\in \mathbb{R}
\\ \end{aligned}\)}{\(\begin{aligned}[t]  x& = 2+t, &
\\y & = 3 - t;\  t\in \mathbb{R}
\\ \end{aligned}\)}{\(\begin{aligned}[t]  x& = 2t+7, &
\\y & = - 3t+1;\  t\in \mathbb{R}
\\ \end{aligned}\)}{}{}}
\LANGes{
\Question{
En la siguiente lista identifica una recta que es perpendicular a la recta 
\( 2x +3y -7= 0\).}
{\(\begin{aligned}[t]  x& = 2t, &
\\y & = -11+3t;\  t\in \mathbb{R}
\\ \end{aligned}\)}{\(\begin{aligned}[t]  x& = 1+3t, &
\\y & = 11 - 2t;\  t\in \mathbb{R}
\\ \end{aligned}\)}{\(\begin{aligned}[t]  x& = 2+t, &
\\y & = 3 - t;\  t\in \mathbb{R}
\\ \end{aligned}\)}{\(\begin{aligned}[t]  x& = 2t+7, &
\\y & = - 3t+1;\  t\in \mathbb{R}
\\ \end{aligned}\)}{}{}}
\End

\Data\ID{2010014604}
\Updated{2023-07-04 09:07:23}
\Level{A}
\SubArea{Analytic geometry in a plane}
\LANGen{
\Question{
Among the lines in the following list (slope-intercept form) identify a line perpendicular
to the line
\[ 
 y = \frac{2}{3}x - 1.
\]}
{\(y = -\frac{3}
{2}x +1\)}{\(y = \frac{2}
{3}x +1\)}{\(y = \frac{3}
{2}x - 1\)}{\(y = -\frac{1}
{2}x + 1\)}{}{}}
\LANGpl{
\Question{
Wśród prostych na poniższej liście (postać kierunkowa) wskaż prostą prostopadłą do prostej
\[
 y = \frac{2}{3}x - 1.
\]}
{\(y = -\frac{3}
{2}x +1\)}{\(y = \frac{2}
{3}x +1\)}{\(y = \frac{3}
{2}x - 1\)}{\(y = -\frac{1}
{2}x + 1\)}{}{}}
\LANGcs{
\Question{
Z následujících přímek zadaných rovnicí ve směrnicovém tvaru vyberte tu, která je kolmá k přímce
\[ 
 y = \frac{2}{3}x - 1.
\]}
{\(y = -\frac{3}
{2}x +1\)}{\(y = \frac{2}
{3}x +1\)}{\(y = \frac{3}
{2}x - 1\)}{\(y = -\frac{1}
{2}x + 1\)}{}{}}
\LANGsk{
\Question{
Z nasledujúcich priamok zadaných rovnicou v smernicovom tvare vyberte tú, ktorá je kolmá na priamku
\[ 
 y = \frac{2}{3}x - 1.
\]}
{\(y = -\frac{3}
{2}x +1\)}{\(y = \frac{2}
{3}x +1\)}{\(y = \frac{3}
{2}x - 1\)}{\(y = -\frac{1}
{2}x + 1\)}{}{}}
\LANGes{
\Question{
Entre las rectas de la siguiente lista (en forma de ecuación explícita) identifica una recta perpendicular a la recta
\[ 
 y = \frac{2}{3}x - 1.
\]}
{\(y = -\frac{3}
{2}x +1\)}{\(y = \frac{2}
{3}x +1\)}{\(y = \frac{3}
{2}x - 1\)}{\(y = -\frac{1}
{2}x + 1\)}{}{}}
\End

\Data\ID{9000106001}
\Updated{2020-07-17 07:07:28}
\Level{A}
\SubArea{Analytic geometry in a plane}
\LANGen{
\Question{
In the following list identify a vector in the direction of the following parametric
line.
\[\begin{aligned}
 x =\ &1 + t, & &
 \\y =\ &3 + 2t;\  t\in \mathbb{R} & &
\end{aligned}\]}
{\(\left (1;2\right )\)}{\(\left (1;3\right )\)}{\(\left (0;2\right )\)}{\(\left (3;1\right )\)}{}{}}
\LANGpl{
\Question{
Wyznacz wektor w kierunku prostej wyrażonej równaniem parametrycznym
\[\begin{aligned}
 x =\ &1 + t, & &
 \\y =\ &3 + 2t;\  t\in \mathbb{R} & &
\end{aligned}\]}
{\(\left (1;2\right )\)}{\(\left (1;3\right )\)}{\(\left (0;2\right )\)}{\(\left (3;1\right )\)}{}{}}
\LANGcs{
\Question{
Z nabízených možností vyberte směrový vektor přímky, která je
vyjádřena parametrickými rovnicemi:
\[\begin{aligned}
 x =\ &1 + t, & &
 \\y =\ &3 + 2t;\  t\in \mathbb{R} & &
\end{aligned}\]}
{\(\left (1;2\right )\)}{\(\left (1;3\right )\)}{\(\left (0;2\right )\)}{\(\left (3;1\right )\)}{}{}}
\LANGsk{
\Question{
Z ponúknutých možností vyberte smerový vektor priamky, ktorá je
vyjadrená parametrickými rovnicami:
\[\begin{aligned}
 x =\ &1 + t, & &
 \\y =\ &3 + 2t;\  t\in \mathbb{R} & &
\end{aligned}\]}
{\(\left (1;2\right )\)}{\(\left (1;3\right )\)}{\(\left (0;2\right )\)}{\(\left (3;1\right )\)}{}{}}
\LANGes{
\Question{
En la siguiente lista, identifica un vector director de la recta expresada por ecuaciones paramétricas.
\[\begin{aligned}
 x =\ &1 + t, & &
 \\y =\ &3 + 2t;\  t\in \mathbb{R} & &
\end{aligned}\]}
{\(\left (1;2\right )\)}{\(\left (1;3\right )\)}{\(\left (0;2\right )\)}{\(\left (3;1\right )\)}{}{}}
\End

\Data\ID{9000106002}
\Updated{2020-07-17 07:07:31}
\Level{A}
\SubArea{Analytic geometry in a plane}
\LANGen{
\Question{
In the following list identify a vector in the direction of the following parametric
line.
\[\begin{aligned}
 x =\ &t - 1, & &
 \\y =\ &t - 2;\ t\in \mathbb{R}\text{.} & &
\end{aligned}\]}
{\(\left (1;1\right )\)}{\(\left (1;2\right )\)}{\(\left (-1;-2\right )\)}{\(\left (1;-1\right )\)}{}{}}
\LANGpl{
\Question{
Wyznacz wektor w kierunku prostej wyrażonej równaniem parametrycznym
\[\begin{aligned}
 x =\ &t - 1, & &
 \\y =\ &t - 2;\ t\in \mathbb{R}\text{.} & &
\end{aligned}\]}
{\(\left (1;1\right )\)}{\(\left (1;2\right )\)}{\(\left (-1;-2\right )\)}{\(\left (1;-1\right )\)}{}{}}
\LANGcs{
\Question{
Z nabízených možností vyberte směrový vektor přímky, která je
vyjádřena parametrickými rovnicemi:
\[\begin{aligned}
 x =\ &t - 1, & &
 \\y =\ &t - 2;\ t\in \mathbb{R}\text{.} & &
\end{aligned}\]}
{\(\left (1;1\right )\)}{\(\left (1;2\right )\)}{\(\left (-1;-2\right )\)}{\(\left (1;-1\right )\)}{}{}}
\LANGsk{
\Question{
Z ponúknutých možností vyberte smerový vektor priamky, ktorá je
vyjadrená parametrickými rovnicami:
\[\begin{aligned}
 x =\ &t - 1, & &
 \\y =\ &t - 2;\ t\in \mathbb{R}\text{.} & &
\end{aligned}\]}
{\(\left (1;1\right )\)}{\(\left (1;2\right )\)}{\(\left (-1;-2\right )\)}{\(\left (1;-1\right )\)}{}{}}
\LANGes{
\Question{
En la siguiente lista, identifica un vector director de la recta expresada por ecuaciones paramétricas.
\[\begin{aligned}
 x =\ &t - 1, & &
 \\y =\ &t - 2;\ t\in \mathbb{R}\text{.} & &
\end{aligned}\]}
{\(\left (1;1\right )\)}{\(\left (1;2\right )\)}{\(\left (-1;-2\right )\)}{\(\left (1;-1\right )\)}{}{}}
\End

\Data\ID{9000106003}
\Updated{2020-07-17 07:07:40}
\Level{A}
\SubArea{Analytic geometry in a plane}
\LANGen{
\Question{
In the following list identify a vector in the direction of the following parametric
line.
\[\begin{aligned}
 x =\ &2, & &
 \\y =\ &t;\ t\in \mathbb{R} & &
\end{aligned}\]}
{\(\left (0;1\right )\)}{\(\left (2;1\right )\)}{\(\left (2;0\right )\)}{\(\left (1;0\right )\)}{}{}}
\LANGpl{
\Question{
Wyznacz wektor w kierunku prostej wyrażonej równaniem parametrycznym.
\[\begin{aligned}
 x =\ &2, & &
 \\y =\ &t;\ t\in \mathbb{R} & &
\end{aligned}\]}
{\(\left (0;1\right )\)}{\(\left (2;1\right )\)}{\(\left (2;0\right )\)}{\(\left (1;0\right )\)}{}{}}
\LANGcs{
\Question{
Z nabízených možností vyberte směrový vektor přímky, která je
vyjádřena parametrickými rovnicemi:
\[\begin{aligned}
 x =\ &2, & &
 \\y =\ &t;\ t\in \mathbb{R} & &
\end{aligned}\]}
{\(\left (0;1\right )\)}{\(\left (2;1\right )\)}{\(\left (2;0\right )\)}{\(\left (1;0\right )\)}{}{}}
\LANGsk{
\Question{
Z ponúknutých možností vyberte smerový vektor priamky, ktorá je
vyjadrená parametrickými rovnicami:
\[\begin{aligned}
 x =\ &2, & &
 \\y =\ &t;\ t\in \mathbb{R} & &
\end{aligned}\]}
{\(\left (0;1\right )\)}{\(\left (2;1\right )\)}{\(\left (2;0\right )\)}{\(\left (1;0\right )\)}{}{}}
\LANGes{
\Question{
En la siguiente lista, identifica un vector director de la recta expresada por ecuaciones paramétricas.
\[\begin{aligned}
 x =\ &2, & &
 \\y =\ &t;\ t\in \mathbb{R} & &
\end{aligned}\]}
{\(\left (0;1\right )\)}{\(\left (2;1\right )\)}{\(\left (2;0\right )\)}{\(\left (1;0\right )\)}{}{}}
\End

\Data\ID{9000106004}
\Updated{2020-07-17 07:07:47}
\Level{A}
\SubArea{Analytic geometry in a plane}
\LANGen{
\Question{
In the following list identify a vector in the direction of the following parametric
line.

\[\begin{aligned}
 x =\ &t, & &
 \\y =\ &1;\ t\in \mathbb{R}. & &
\end{aligned}\]}
{\(\left (1;0\right )\)}{\(\left (0;0\right )\)}{\(\left (0;1\right )\)}{\(\left (1;1\right )\)}{}{}}
\LANGpl{
\Question{
Wyznacz wektor w kierunku prostej wyrażonej równaniem parametrycznym.

\[\begin{aligned}
 x =\ &t, & &
 \\y =\ &1;\ t\in \mathbb{R}. & &
\end{aligned}\]}
{\(\left (1;0\right )\)}{\(\left (0;0\right )\)}{\(\left (0;1\right )\)}{\(\left (1;1\right )\)}{}{}}
\LANGcs{
\Question{
Z nabízených možností vyberte směrový vektor přímky, která je
vyjádřena parametrickými rovnicemi:

\[\begin{aligned}
 x =\ &t, & &
 \\y =\ &1;\ t\in \mathbb{R}. & &
\end{aligned}\]}
{\(\left (1;0\right )\)}{\(\left (0;0\right )\)}{\(\left (0;1\right )\)}{\(\left (1;1\right )\)}{}{}}
\LANGsk{
\Question{
Z ponúknutých možností vyberte smerový vektor priamky, ktorá je
vyjadrená parametrickými rovnicami:
\[\begin{aligned}
 x =\ &t, & &
 \\y =\ &1;\ t\in \mathbb{R}. & &
\end{aligned}\]}
{\(\left (1;0\right )\)}{\(\left (0;0\right )\)}{\(\left (0;1\right )\)}{\(\left (1;1\right )\)}{}{}}
\LANGes{
\Question{
En la siguiente lista, identifica un vector director de la recta expresada por ecuaciones paramétricas.

\[\begin{aligned}
 x =\ &t, & &
 \\y =\ &1;\ t\in \mathbb{R}. & &
\end{aligned}\]}
{\(\left (1;0\right )\)}{\(\left (0;0\right )\)}{\(\left (0;1\right )\)}{\(\left (1;1\right )\)}{}{}}
\End

\Data\ID{9000106005}
\Updated{2022-04-23 05:04:00}
\Level{A}
\SubArea{Analytic geometry in a plane}
\LANGen{
\Question{
In the following list identify a vector having the same direction as the line passing through
the points \(A\) 
and \(B\).
\[ 
 A = \left [2;1\right ]\text{, }\qquad B = \left [3;2\right ]
\]}
{\(\left (1;1\right )\)}{\(\left (-1;1\right )\)}{\(\left (5;3\right )\)}{\(\left (3;5\right )\)}{}{}}
\LANGpl{
\Question{
Wyznacz wektor o tym samym kierunku co prosta przechodząca przez punkty 
 \(A\) 
i \(B\).
\[ 
 A = \left [2;1\right ]\text{, }\qquad B = \left [3;2\right ]
\]}
{\(\left (1;1\right )\)}{\(\left (-1;1\right )\)}{\(\left (5;3\right )\)}{\(\left (3;5\right )\)}{}{}}
\LANGcs{
\Question{
Z nabízených možností vyberte směrový vektor přímky, která prochází
body \(A\) 
a \(B\), kde
\[ 
 A = \left [2;1\right ]\text{, }B = \left [3;2\right ]\text{.}
\]}
{\(\left (1;1\right )\)}{\(\left (-1;1\right )\)}{\(\left (5;3\right )\)}{\(\left (3;5\right )\)}{}{}}
\LANGsk{
\Question{
Z ponúknutých možností vyberte smerový vektor priamky, ktorá prechádza
bodmi \(A\) 
a \(B\), kde
\[ 
 A = \left [2;1\right ]\text{, }B = \left [3;2\right ]\text{.}
\]}
{\(\left (1;1\right )\)}{\(\left (-1;1\right )\)}{\(\left (5;3\right )\)}{\(\left (3;5\right )\)}{}{}}
\LANGes{
\Question{
En la siguiente lista, identifica un vector que tiene la misma dirección que la recta que pasa por los puntos \(A\) 
y \(B\).
\[ 
 A = \left [2;1\right ]\text{, }\qquad B = \left [3;2\right ]
\]}
{\(\left (1;1\right )\)}{\(\left (-1;1\right )\)}{\(\left (5;3\right )\)}{\(\left (3;5\right )\)}{}{}}
\End

\Data\ID{9000106006}
\Updated{2021-12-05 10:12:32}
\Level{A}
\SubArea{Analytic geometry in a plane}
\LANGen{
\Question{
In the following list identify a vector having the same direction as the line passing through
the points \(A\) 
and \(B\).
\[ 
 A = \left [-3;-1\right ]\text{, }\qquad B = \left [-1;-2\right ]
\]}
{\(\left (2;-1\right )\)}{\(\left (-4;-3\right )\)}{\(\left (1;2\right )\)}{\(\left (2;1\right )\)}{}{}}
\LANGpl{
\Question{
Wyznacz współrzędne wektora o tym samym kierunku co prosta przechodząca przez punkty 
\(A\) 
i \(B\).
\[ 
 A = \left [-3;-1\right ]\text{, }\qquad B = \left [-1;-2\right ]
\]}
{\(\left (2;-1\right )\)}{\(\left (-4;-3\right )\)}{\(\left (1;2\right )\)}{\(\left (2;1\right )\)}{}{}}
\LANGcs{
\Question{
Z nabízených možností vyberte směrový vektor přímky, která prochází
body \(A\) 
a \(B\), kde
\[ 
 A = \left [-3;-1\right ]\text{, }B = \left [-1;-2\right ]\text{.}
\]}
{\(\left (2;-1\right )\)}{\(\left (-4;-3\right )\)}{\(\left (1;2\right )\)}{\(\left (2;1\right )\)}{}{}}
\LANGsk{
\Question{
Z ponúknutých možností vyberte smerový vektor priamky, ktorá prechádza
bodmi \(A\) 
a \(B\), kde
\[ 
 A = \left [-3;-1\right ]\text{, }B = \left [-1;-2\right ]\text{.}
\]}
{\(\left (2;-1\right )\)}{\(\left (-4;-3\right )\)}{\(\left (1;2\right )\)}{\(\left (2;1\right )\)}{}{}}
\LANGes{
\Question{
En la siguiente lista, identifica un vector que tiene la misma dirección que la recta que pasa por los puntos \(A\) 
y \(B\).
\[ 
 A = \left [-3;-1\right ]\text{, }\qquad B = \left [-1;-2\right ]
\]}
{\(\left (2;-1\right )\)}{\(\left (-4;-3\right )\)}{\(\left (1;2\right )\)}{\(\left (2;1\right )\)}{}{}}
\End

\Data\ID{9000106007}
\Updated{2020-07-17 07:07:49}
\Level{A}
\SubArea{Analytic geometry in a plane}
\LANGen{
\Question{
In the following list identify a normal vector of the following line. 
\[ 
 2x + 3y + 4 = 0
\]}
{\(\left (2;3\right )\)}{\(\left (3;4\right )\)}{\(\left (2;4\right )\)}{\(\left (2;-3\right )\)}{}{}}
\LANGpl{
\Question{
Wyznacz wektor prostopadły wyrażony  równaniem.
\[ 
 2x + 3y + 4 = 0
\]}
{\(\left (2;3\right )\)}{\(\left (3;4\right )\)}{\(\left (2;4\right )\)}{\(\left (2;-3\right )\)}{}{}}
\LANGcs{
\Question{
Z nabízených možností vyberte normálový vektor přímky, která je
vyjádřena obecnou rovnicí 
\[ 
 2x + 3y + 4 = 0\text{.}
\]}
{\(\left (2;3\right )\)}{\(\left (3;4\right )\)}{\(\left (2;4\right )\)}{\(\left (2;-3\right )\)}{}{}}
\LANGsk{
\Question{
Z ponúknutých možností vyberte normálový vektor priamky, ktorá je
vyjadrená všeobecnou rovnicou 
\[ 
 2x + 3y + 4 = 0\text{.}
\]}
{\(\left (2;3\right )\)}{\(\left (3;4\right )\)}{\(\left (2;4\right )\)}{\(\left (2;-3\right )\)}{}{}}
\LANGes{
\Question{
En la siguiente lista, identifica un vector normal de la recta
\[ 
 2x + 3y + 4 = 0
\].}
{\(\left (2;3\right )\)}{\(\left (3;4\right )\)}{\(\left (2;4\right )\)}{\(\left (2;-3\right )\)}{}{}}
\End

\Data\ID{9000106008}
\Updated{2020-07-17 07:07:11}
\Level{A}
\SubArea{Analytic geometry in a plane}
\LANGen{
\Question{
In the following list identify a normal vector of the following line. 
\[ 
 x - 2y - 10 = 0
\]}
{\(\left (1;-2\right )\)}{\(\left (1;2\right )\)}{\(\left (2;10\right )\)}{\(\left (-2;-10\right )\)}{}{}}
\LANGpl{
\Question{
Wyznacz wektor prostopadły wyrażony równaniem.
\[ 
 x - 2y - 10 = 0
\]}
{\(\left (1;-2\right )\)}{\(\left (1;2\right )\)}{\(\left (2;10\right )\)}{\(\left (-2;-10\right )\)}{}{}}
\LANGcs{
\Question{
Z nabízených možností vyberte normálový vektor přímky, která je
vyjádřena obecnou rovnicí 
\[ 
 x - 2y - 10 = 0\text{.}
\]}
{\(\left (1;-2\right )\)}{\(\left (1;2\right )\)}{\(\left (2;10\right )\)}{\(\left (-2;-10\right )\)}{}{}}
\LANGsk{
\Question{
Z ponúknutých možností vyberte normálový vektor priamky, ktorá je
vyjadrená všeobecnou rovnicou 
\[ 
 x - 2y - 10 = 0\text{.}
\]}
{\(\left (1;-2\right )\)}{\(\left (1;2\right )\)}{\(\left (2;10\right )\)}{\(\left (-2;-10\right )\)}{}{}}
\LANGes{
\Question{
En la siguiente lista, identifica un vector normal de la recta.
\[ 
 x - 2y - 10 = 0
\]}
{\(\left (1;-2\right )\)}{\(\left (1;2\right )\)}{\(\left (2;10\right )\)}{\(\left (-2;-10\right )\)}{}{}}
\End

\Data\ID{9000106009}
\Updated{2020-07-17 07:07:49}
\Level{A}
\SubArea{Analytic geometry in a plane}
\LANGen{
\Question{
In the following list identify a normal vector of the following line. 
\[ 
 x - 2 = 0
\]}
{\(\left (1;0\right )\)}{\(\left (1;2\right )\)}{\(\left (1;-2\right )\)}{\(\left (0;-2\right )\)}{}{}}
\LANGpl{
\Question{
Wyznacz wektor prostopadły wyrażony równaniem.
\[ 
 x - 2 = 0
\]}
{\(\left (1;0\right )\)}{\(\left (1;2\right )\)}{\(\left (1;-2\right )\)}{\(\left (0;-2\right )\)}{}{}}
\LANGcs{
\Question{
Z nabízených možností vyberte normálový vektor přímky, která je
vyjádřena obecnou rovnicí 
\[ 
 x - 2 = 0\text{.}
\]}
{\(\left (1;0\right )\)}{\(\left (1;2\right )\)}{\(\left (1;-2\right )\)}{\(\left (0;-2\right )\)}{}{}}
\LANGsk{
\Question{
Z ponúknutých možností vyberte normálový vektor priamky, ktorá je
vyjadrená všeobecnou rovnicou 
\[ 
 x - 2 = 0\text{.}
\]}
{\(\left (1;0\right )\)}{\(\left (1;2\right )\)}{\(\left (1;-2\right )\)}{\(\left (0;-2\right )\)}{}{}}
\LANGes{
\Question{
En la siguiente lista, identifica un vector normal de la recta.
\[ 
 x - 2 = 0
\]}
{\(\left (1;0\right )\)}{\(\left (1;2\right )\)}{\(\left (1;-2\right )\)}{\(\left (0;-2\right )\)}{}{}}
\End

\Data\ID{9000106010}
\Updated{2020-07-17 07:07:38}
\Level{A}
\SubArea{Analytic geometry in a plane}
\LANGen{
\Question{
In the following list identify a normal vector of the following line. 
\[ 
 3y - 1 = 0
\]}
{\(\left (0;3\right )\)}{\(\left (3;-1\right )\)}{\(\left (-1;0\right )\)}{\(\left (1;-3\right )\)}{}{}}
\LANGpl{
\Question{
Wyznacz wektor prostopadły wyrażony równaniem.
\[ 
 3y - 1 = 0
\]}
{\(\left (0;3\right )\)}{\(\left (3;-1\right )\)}{\(\left (-1;0\right )\)}{\(\left (1;-3\right )\)}{}{}}
\LANGcs{
\Question{
Z nabízených možností vyberte normálový vektor přímky, která je
vyjádřena obecnou rovnicí 
\[ 
 3y - 1 = 0\text{.}
\]}
{\(\left (0;3\right )\)}{\(\left (3;-1\right )\)}{\(\left (-1;0\right )\)}{\(\left (1;-3\right )\)}{}{}}
\LANGsk{
\Question{
Z ponúknutých možností vyberte normálový vektor priamky, ktorá je
vyjadrená všeobecnou rovnicou
\[ 
 3y - 1 = 0\text{.}
\]}
{\(\left (0;3\right )\)}{\(\left (3;-1\right )\)}{\(\left (-1;0\right )\)}{\(\left (1;-3\right )\)}{}{}}
\LANGes{
\Question{
En la siguiente lista, identifica un vector normal de la recta.
\[ 
 3y - 1 = 0
\]}
{\(\left (0;3\right )\)}{\(\left (3;-1\right )\)}{\(\left (-1;0\right )\)}{\(\left (1;-3\right )\)}{}{}}
\End

\Data\ID{9000106201}
\Updated{2020-07-17 08:07:41}
\Level{A}
\SubArea{Analytic geometry in a plane}
\LANGen{
\Question{
In the following list identify a vector having the same direction as the parametric
line \(p\).
\[ \begin{alignedat}{80}
 p\colon x & = 1 + 2t, & &\phantom{t\in \mathbb{R}} & & & &
 \\y & = 3 - 4t;\ & &t\in \mathbb{R} & & & &
\\\end{alignedat}\]}
{\((1;-2)\)}{\((1;3)\)}{\((3;1)\)}{\((2;3)\)}{}{}}
\LANGpl{
\Question{
Wyznacz wektor o tym samym kierunku co prosta
 \(p\) wyrażona równaniem parametrycznym.
\[ \begin{alignedat}{80}
 p\colon x & = 1 + 2t, & &\phantom{t\in \mathbb{R}} & & & &
 \\y & = 3 - 4t;\ & &t\in \mathbb{R} & & & &
\\\end{alignedat}\]}
{\((1;-2)\)}{\((1;3)\)}{\((3;1)\)}{\((2;3)\)}{}{}}
\LANGcs{
\Question{
Z nabízených možností vyberte směrový vektor přímky, která je
vyjádřena parametrickými rovnicemi:
\[ \begin{alignedat}{80}
 p\colon x & = 1 + 2t, & &\phantom{t\in \mathbb{R}} & & & &
 \\y & = 3 - 4t;\ & &t\in \mathbb{R}. & & & &
\\\end{alignedat}\]}
{\((1;-2)\)}{\((1;3)\)}{\((3;1)\)}{\((2;3)\)}{}{}}
\LANGsk{
\Question{
Z ponúknutých možností vyberte smerový vektor priamky, ktorá je
vyjadrená parametrickými rovnicami:
\[ \begin{alignedat}{80}
 p\colon x & = 1 + 2t, & &\phantom{t\in \mathbb{R}} & & & &
 \\y & = 3 - 4t;\ & &t\in \mathbb{R}. & & & &
\\\end{alignedat}\]}
{\((1;-2)\)}{\((1;3)\)}{\((3;1)\)}{\((2;3)\)}{}{}}
\LANGes{
\Question{
En la siguiente lista, identifica un vector que tiene la misma dirección que la recta paramétrica \(p\).
\[ \begin{alignedat}{80}
 p\colon x & = 1 + 2t, & &\phantom{t\in \mathbb{R}} & & & &
 \\y & = 3 - 4t;\ & &t\in \mathbb{R} & & & &
\\\end{alignedat}\]}
{\((1;-2)\)}{\((1;3)\)}{\((3;1)\)}{\((2;3)\)}{}{}}
\End

\Data\ID{9000106202}
\Updated{2020-07-17 08:07:52}
\Level{A}
\SubArea{Analytic geometry in a plane}
\LANGen{
\Question{
In the following list identify a vector having the same direction as the parametric
line \(p\).
\[ \begin{aligned}
 x & = 1 - t, \\
y & = t;\ t\in \mathbb{R}
\\\end{aligned}\]}
{\((1;-1)\)}{\((0;0)\)}{\((1;0)\)}{\((1;1)\)}{}{}}
\LANGpl{
\Question{
Wyznacz wektor o tym samym kierunku co prosta
 \(p\) wyrażona równaniem parametrycznym.
\[ \begin{aligned}
 x & = 1 - t, \\
y & = t;\ t\in \mathbb{R}
\\\end{aligned}\]}
{\((1;-1)\)}{\((0;0)\)}{\((1;0)\)}{\((1;1)\)}{}{}}
\LANGcs{
\Question{
Z nabízených možností vyberte směrový vektor přímky, která je
vyjádřena parametrickými rovnicemi:
\[ \begin{aligned}
 x & = 1 - t, \\
y & = t;\ t\in \mathbb{R}.
\\\end{aligned}\]}
{\((1;-1)\)}{\((0;0)\)}{\((1;0)\)}{\((1;1)\)}{}{}}
\LANGsk{
\Question{
Z ponúknutých možností vyberte smerový vektor priamky, ktorá je
vyjadrená parametrickými rovnicami:
\[ \begin{aligned}
 x & = 1 - t, \\
y & = t;\ t\in \mathbb{R}.
\\\end{aligned}\]}
{\((1;-1)\)}{\((0;0)\)}{\((1;0)\)}{\((1;1)\)}{}{}}
\LANGes{
\Question{
En la siguiente lista, identifica un vector que tiene la misma dirección que la recta paramétrica \(p\).
\[ \begin{aligned}
 x & = 1 - t, \\
y & = t;\ t\in \mathbb{R}
\\\end{aligned}\]}
{\((1;-1)\)}{\((0;0)\)}{\((1;0)\)}{\((1;1)\)}{}{}}
\End

\Data\ID{9000106203}
\Updated{2020-07-17 08:07:00}
\Level{A}
\SubArea{Analytic geometry in a plane}
\LANGen{
\Question{
In the following list identify a vector having the same direction as the parametric
line \(p\).
\[ \begin{aligned}
 p\colon x & = -5, \\
y & = 5t;\ t\in \mathbb{R}
\end{aligned}\]}
{\((0;1)\)}{\((5;5)\)}{\((5;0)\)}{\((-5;5)\)}{}{}}
\LANGpl{
\Question{
Wyznacz wektor o tym samym kierunku co prosta
 \(p\) wyrażona równaniem parametrycznym.

\[ \begin{aligned}
 p\colon x & = -5, \\
y & = 5t;\ t\in \mathbb{R}
\end{aligned}\]}
{\((0;1)\)}{\((5;5)\)}{\((5;0)\)}{\((-5;5)\)}{}{}}
\LANGcs{
\Question{
Z nabízených možností vyberte směrový vektor přímky, která je
vyjádřena parametrickými rovnicemi:

\[ \begin{aligned}
 p\colon x & = -5, \\
y & = 5t;\ t\in \mathbb{R}.
\end{aligned}\]}
{\((0;1)\)}{\((5;5)\)}{\((5;0)\)}{\((-5;5)\)}{}{}}
\LANGsk{
\Question{
Z ponúknutých možností vyberte smerový vektor priamky, ktorá je
vyjadrená parametrickými rovnicami:

\[ \begin{aligned}
 p\colon x & = -5, \\
y & = 5t;\ t\in \mathbb{R}.
\end{aligned}\]}
{\((0;1)\)}{\((5;5)\)}{\((5;0)\)}{\((-5;5)\)}{}{}}
\LANGes{
\Question{
En la siguiente lista, identifica un vector que tiene la misma dirección que la recta paramétrica \(p\).
\[ \begin{aligned}
 p\colon x & = -5, \\
y & = 5t;\ t\in \mathbb{R}
\end{aligned}\]}
{\((0;1)\)}{\((5;5)\)}{\((5;0)\)}{\((-5;5)\)}{}{}}
\End

\Data\ID{9000106204}
\Updated{2020-07-17 08:07:13}
\Level{A}
\SubArea{Analytic geometry in a plane}
\LANGen{
\Question{
In the following list identify a vector having the same direction as the parametric
line \(p\).
\[ \begin{aligned}
 x & = 2t, \\
y & = 0;\ t\in \mathbb{R}.
\\\end{aligned}\]}
{\((1;0)\)}{\((0;0)\)}{\((0;1)\)}{\((0;2)\)}{}{}}
\LANGpl{
\Question{
Wyznacz wektor o tym samym kierunku co prosta
 \(p\) wyrażona równaniem parametrycznym.
\[ \begin{aligned}
 x & = 2t, \\
y & = 0;\ t\in \mathbb{R}.
\\\end{aligned}\]}
{\((1;0)\)}{\((0;0)\)}{\((0;1)\)}{\((0;2)\)}{}{}}
\LANGcs{
\Question{
Z nabízených možností vyberte směrový vektor přímky, která je
vyjádřena parametrickými rovnicemi:
\[ \begin{aligned}
 x & = 2t, \\
y & = 0;\ t\in \mathbb{R}.
\\\end{aligned}\]}
{\((1;0)\)}{\((0;0)\)}{\((0;1)\)}{\((0;2)\)}{}{}}
\LANGsk{
\Question{
Z ponúknutých možností vyberte smerový vektor priamky, ktorá je
vyjadrená parametrickými rovnicami:
\[ \begin{aligned}
 x & = 2t, \\
y & = 0;\ t\in \mathbb{R}.
\\\end{aligned}\]}
{\((1;0)\)}{\((0;0)\)}{\((0;1)\)}{\((0;2)\)}{}{}}
\LANGes{
\Question{
En la siguiente lista, identifica un vector que tiene la misma dirección que la recta paramétrica \(p\).
\[ \begin{aligned}
 x & = 2t, \\
y & = 0;\ t\in \mathbb{R}.
\\\end{aligned}\]}
{\((1;0)\)}{\((0;0)\)}{\((0;1)\)}{\((0;2)\)}{}{}}
\End

\Data\ID{9000106205}
\Updated{2021-12-05 07:12:25}
\Level{A}
\SubArea{Analytic geometry in a plane}
\LANGen{
\Question{
In the following list find the vector having the same direction as the line passing through
the points \(A = [3;-3]\) 
and \(B = [-1;-9]\).}
{\((2;3)\)}{\((2;-12)\)}{\((4;-6)\)}{\((-4;6)\)}{}{}}
\LANGpl{
\Question{
Wyznacz wektor o tym samym kierunku co prosta przechodząca przez punkty \(A = [3;-3]\) 
i \(B = [-1;-9]\).}
{\((2;3)\)}{\((2;-12)\)}{\((4;-6)\)}{\((-4;6)\)}{}{}}
\LANGcs{
\Question{
Z nabízených možností vyberte směrový vektor přímky, která prochází
body \(A\),
\(B\), kde
\(A = [3;-3]\),
\(B = [-1;-9]\).}
{\((2;3)\)}{\((2;-12)\)}{\((4;-6)\)}{\((-4;6)\)}{}{}}
\LANGsk{
\Question{
Z ponúkaných možností vyberte smerový vektor priamky, ktorá prechádza
bodmi \(A\),
\(B\), kde
\(A = [3;-3]\),
\(B = [-1;-9]\).}
{\((2;3)\)}{\((2;-12)\)}{\((4;-6)\)}{\((-4;6)\)}{}{}}
\LANGes{
\Question{
En la siguiente lista, identifica un vector que tiene la misma dirección que la recta que pasa por los puntos \(A = [3;-3]\) 
y \(B = [-1;-9]\).}
{\((2;3)\)}{\((2;-12)\)}{\((4;-6)\)}{\((-4;6)\)}{}{}}
\End

\Data\ID{9000106206}
\Updated{2021-12-05 07:12:41}
\Level{A}
\SubArea{Analytic geometry in a plane}
\LANGen{
\Question{
In the following list find the vector having the same direction as the line passing through
the points \(A = [3;4]\) 
and \(B = [5;8]\).}
{\((1;2)\)}{\((4;2)\)}{\((1;3)\)}{\((2;-4)\)}{}{}}
\LANGpl{
\Question{
Wyznacz wektor o tym samym kierunku co prosta przechodząca przez punkty \(A = [3;4]\) 
i \(B = [5;8]\).}
{\((1;2)\)}{\((4;2)\)}{\((1;3)\)}{\((2;-4)\)}{}{}}
\LANGcs{
\Question{
Z nabízených možností vyberte směrový vektor přímky, která prochází
body \(A\),
\(B\), kde
\(A = [3;4]\),
\(B = [5;8]\).}
{\((1;2)\)}{\((4;2)\)}{\((1;3)\)}{\((2;-4)\)}{}{}}
\LANGsk{
\Question{
Z ponúknutých možností vyberte smerový vektor priamky, ktorá prechádza
bodmi \(A\),
\(B\), kde
\(A = [3;4]\),
\(B = [5;8]\).}
{\((1;2)\)}{\((4;2)\)}{\((1;3)\)}{\((2;-4)\)}{}{}}
\LANGes{
\Question{
En la siguiente lista, identifica un vector que tiene la misma dirección que la recta que pasa por los puntos  \(A = [3;4]\) 
y \(B = [5;8]\).}
{\((1;2)\)}{\((4;2)\)}{\((1;3)\)}{\((2;-4)\)}{}{}}
\End

\Data\ID{9000106207}
\Updated{2020-07-17 08:07:27}
\Level{A}
\SubArea{Analytic geometry in a plane}
\LANGen{
\Question{
In the following list identify a normal vector to the line
\(p\).
\[ 
 p\colon - x + 2y + 3 = 0
\]}
{\((1;-2)\)}{\((2;3)\)}{\((-1;3)\)}{\((2;-1)\)}{}{}}
\LANGpl{
\Question{
Wskaż wektor prostopadły do prostej
\(p\).
\[ 
 p\colon - x + 2y + 3 = 0
\]}
{\((1;-2)\)}{\((2;3)\)}{\((-1;3)\)}{\((2;-1)\)}{}{}}
\LANGcs{
\Question{
Z nabízených možností vyberte normálový vektor přímky, která je
vyjádřena obecnou rovnicí: 
\[ 
 -x + 2y + 3 = 0
\]}
{\((1;-2)\)}{\((2;3)\)}{\((-1;3)\)}{\((2;-1)\)}{}{}}
\LANGsk{
\Question{
Z ponúkaných možností vyberte normálový vektor priamky, ktorá je
vyjadrená všeobecnou rovnicou: 
\[ 
 -x + 2y + 3 = 0
\]}
{\((1;-2)\)}{\((2;3)\)}{\((-1;3)\)}{\((2;-1)\)}{}{}}
\LANGes{
\Question{
En la siguiente lista, identifica un vector normal de la recta
\(p\).
\[ 
 p\colon - x + 2y + 3 = 0
\]}
{\((1;-2)\)}{\((2;3)\)}{\((-1;3)\)}{\((2;-1)\)}{}{}}
\End

\Data\ID{9000106208}
\Updated{2020-07-17 08:07:28}
\Level{A}
\SubArea{Analytic geometry in a plane}
\LANGen{
\Question{
In the following list identify a normal vector to the line
\(p\).
\[ 
 p\colon 2y - 1 = 0
\]}
{\((0;1)\)}{\((2;0)\)}{\((2;-1)\)}{$(2;1)$}{}{}}
\LANGpl{
\Question{
Wskaż wektor prostopadły do prostej 
\(p\).
\[ 
 p\colon 2y - 1 = 0
\]}
{\((0;1)\)}{\((2;0)\)}{\((2;-1)\)}{$(2;1)$}{}{}}
\LANGcs{
\Question{
Z nabízených možností vyberte normálový vektor přímky, která je
vyjádřena obecnou rovnicí: 
\[ 
 2y - 1 = 0\text{.}
\]}
{\((0;1)\)}{\((2;0)\)}{\((2;-1)\)}{$(2;1)$}{}{}}
\LANGsk{
\Question{
Z ponúknutých možností vyberte normálový vektor priamky, ktorá je
vyjadrená všeobecnou rovnicou: 
\[ 
 2y - 1 = 0\text{.}
\]}
{\((0;1)\)}{\((2;0)\)}{\((2;-1)\)}{$(2;1)$}{}{}}
\LANGes{
\Question{
En la siguiente lista, identifica un vector normal de la recta
\(p\).
\[ 
 p\colon 2y - 1 = 0
\]}
{\((0;1)\)}{\((2;0)\)}{\((2;-1)\)}{$(2;1)$}{}{}}
\End

\Data\ID{9000106209}
\Updated{2020-07-17 08:07:55}
\Level{A}
\SubArea{Analytic geometry in a plane}
\LANGen{
\Question{
In the following list find the vector having the same direction as the line
\(p\).
\[ 
 p\colon y = 2x + 1
\]}
{\((1;2)\)}{\((2;-1)\)}{\((2;1)\)}{\((1;-2)\)}{}{}}
\LANGpl{
\Question{
Wskaż wektor o tym samym kierunku co prosta
\(p\).
\[ 
 p\colon y = 2x + 1
\]}
{\((1;2)\)}{\((2;-1)\)}{\((2;1)\)}{\((1;-2)\)}{}{}}
\LANGcs{
\Question{
Z nabízených možností vyberte směrový vektor přímky, která je
vyjádřena ve směrnicovém tvaru rovnicí: 
\[ 
 y = 2x + 1\text{.}
\]}
{\((1;2)\)}{\((2;-1)\)}{\((2;1)\)}{\((1;-2)\)}{}{}}
\LANGsk{
\Question{
Z ponúknutých možností vyberte smerový vektor priamky, ktorá je
vyjadrená v smernicovom tvare rovnicou: 
\[ 
 y = 2x + 1\text{.}
\]}
{\((1;2)\)}{\((2;-1)\)}{\((2;1)\)}{\((1;-2)\)}{}{}}
\LANGes{
\Question{
En la siguiente lista, identifica un vector que tiene la misma dirección que la recta 
\(p\).
\[ 
 p\colon y = 2x + 1
\]}
{\((1;2)\)}{\((2;-1)\)}{\((2;1)\)}{\((1;-2)\)}{}{}}
\End

\Data\ID{9000106210}
\Updated{2022-04-23 05:04:00}
\Level{A}
\SubArea{Analytic geometry in a plane}
\LANGen{
\Question{
In the following list find the vector having the same direction as the line
\(p\).
\[ 
 p\colon 3x - 2y + 1 = 0
\]}
{\((2;3)\)}{\((3;-2)\)}{\((-2;1)\)}{\((3;2)\)}{}{}}
\LANGpl{
\Question{
Wskaż wektor o tym samym kierunku co prosta
\(p\).
\[ 
 p\colon 3x - 2y + 1 = 0
\]}
{\((2;3)\)}{\((3;-2)\)}{\((-2;1)\)}{\((3;2)\)}{}{}}
\LANGcs{
\Question{
Z nabízených možností vyberte směrový vektor přímky, která je
vyjádřena obecnou rovnicí: 
\[ 
 3x - 2y + 1 = 0\text{.}
\]}
{\((2;3)\)}{\((3;-2)\)}{\((-2;1)\)}{\((3;2)\)}{}{}}
\LANGsk{
\Question{
Z ponúknutých možností vyberte smerový vektor priamky, ktorá je
vyjadrená všeobecnou rovnicou: 
\[ 
 3x - 2y + 1 = 0\text{.}
\]}
{\((2;3)\)}{\((3;-2)\)}{\((-2;1)\)}{\((3;2)\)}{}{}}
\LANGes{
\Question{
En la siguiente lista, identifica un vector que tiene la misma dirección que la recta 
\(p\).
\[ 
 p\colon 3x - 2y + 1 = 0
\]}
{\((2;3)\)}{\((3;-2)\)}{\((-2;1)\)}{\((3;2)\)}{}{}}
\End

\Data\ID{9000106801}
\Updated{2020-07-17 08:07:07}
\Level{A}
\SubArea{Analytic geometry in a plane}
\LANGen{
\Question{
Find the vector in the direction of the following line. 
\[ 
 2x + 1 = 3y - 2
\]}
{\((3;2)\)}{\((-3;2)\)}{\((2;-3)\)}{\((-3;3)\)}{}{}}
\LANGpl{
\Question{
Wskaż wektor o tym samym kierunku co prosta wyrażona równaniem.
\[ 
 2x + 1 = 3y - 2
\]}
{\((3;2)\)}{\((-3;2)\)}{\((2;-3)\)}{\((-3;3)\)}{}{}}
\LANGcs{
\Question{
Z nabízených možností vyberte směrový vektor přímky, která je
vyjádřena rovnicí: 
\[ 
 2x + 1 = 3y - 2
\]}
{\((3;2)\)}{\((-3;2)\)}{\((2;-3)\)}{\((-3;3)\)}{}{}}
\LANGsk{
\Question{
Z ponúknutých možností vyberte smerový vektor priamky, ktorá je
vyjadrená rovnicou: 
\[ 
 2x + 1 = 3y - 2
\]}
{\((3;2)\)}{\((-3;2)\)}{\((2;-3)\)}{\((-3;3)\)}{}{}}
\LANGes{
\Question{
Halla un vector director de la recta. 
\[ 
 2x + 1 = 3y - 2
\]}
{\((3;2)\)}{\((-3;2)\)}{\((2;-3)\)}{\((-3;3)\)}{}{}}
\End

\Data\ID{9000106802}
\Updated{2022-04-23 05:04:00}
\Level{A}
\SubArea{Analytic geometry in a plane}
\LANGen{
\Question{
Find the normal vector of the line passing through the point
\(A = [3;-1]\) and
\(B = [2;2]\).}
{\((3;1)\)}{\((-1;3)\)}{\((1;-3)\)}{\((1;3)\)}{}{}}
\LANGpl{
\Question{
Wskaż wektor prostopadły do prostej przechodzącej przez punkty
\(A = [3;-1]\) i
\(B = [2;2]\).}
{\((3;1)\)}{\((-1;3)\)}{\((1;-3)\)}{\((1;3)\)}{}{}}
\LANGcs{
\Question{
Z nabízených možností vyberte normálový vektor přímky, která prochází
body \(A\),
\(B\), kde
\(A = [3;-1]\),
\(B = [2;2]\).}
{\((3;1)\)}{\((-1;3)\)}{\((1;-3)\)}{\((1;3)\)}{}{}}
\LANGsk{
\Question{
Z ponúknutých možností vyberte normálový vektor priamky, ktorá prechádza
bodmi \(A\),
\(B\), kde
\(A = [3;-1]\),
\(B = [2;2]\).}
{\((3;1)\)}{\((-1;3)\)}{\((1;-3)\)}{\((1;3)\)}{}{}}
\LANGes{
\Question{
Halla un vector normal de la recta que pasa por los puntos \(A = [3;-1]\) y
\(B = [2;2]\).}
{\((3;1)\)}{\((-1;3)\)}{\((1;-3)\)}{\((1;3)\)}{}{}}
\End

\Data\ID{9000106803}
\Updated{2020-07-17 08:07:53}
\Level{A}
\SubArea{Analytic geometry in a plane}
\LANGen{
\Question{
Find the vector in the direction of the following line. 
\[ 
 p\colon y = \frac{2} 
{3}x - 3
\]}
{\((3;2)\)}{\(\left (\frac{2}
{3};-1\right )\)}{\((3;-1)\)}{\(\left (\frac{2}
{3};1\right )\)}{}{}}
\LANGpl{
\Question{
Wskaż wektor o tym samym kierunku co prosta wyrażona równaniem.
\[ 
 p\colon y = \frac{2} 
{3}x - 3
\]}
{\((3;2)\)}{\(\left (\frac{2}
{3};-1\right )\)}{\((3;-1)\)}{\(\left (\frac{2}
{3};1\right )\)}{}{}}
\LANGcs{
\Question{
Z nabízených možností vyberte směrový vektor přímky, která je
vyjádřena rovnicí ve směrnicovém tvaru: 
\[ 
 y = \frac{2} 
{3}x - 3
\]}
{\((3;2)\)}{\(\left (\frac{2}
{3};-1\right )\)}{\((3;-1)\)}{\(\left (\frac{2}
{3};1\right )\)}{}{}}
\LANGsk{
\Question{
Z ponúknutých možností vyberte smerový vektor priamky, ktorá je
vyjadrená rovnicou v smernicovom tvare: 
\[ 
 y = \frac{2} 
{3}x - 3
\]}
{\((3;2)\)}{\(\left (\frac{2}
{3};-1\right )\)}{\((3;-1)\)}{\(\left (\frac{2}
{3};1\right )\)}{}{}}
\LANGes{
\Question{
Halla un vector normal de la siguiente recta
\[ 
 p\colon y = \frac{2} 
{3}x - 3
\]}
{\((3;2)\)}{\(\left (\frac{2}
{3};-1\right )\)}{\((3;-1)\)}{\(\left (\frac{2}
{3};1\right )\)}{}{}}
\End

\Data\ID{9000106804}
\Updated{2022-04-23 05:04:00}
\Level{A}
\SubArea{Analytic geometry in a plane}
\LANGen{
\Question{
Find the normal vector of the following line. 
\[ 
p\colon \begin{aligned}[t] x =&1 - 6t, &
\\y =& - 2 + 3t;\ t\in \mathbb{R}
\\ \end{aligned}
\]}
{\((1;2)\)}{\((-6;3)\)}{\((1;-2)\)}{\((2;1)\)}{}{}}
\LANGpl{
\Question{
Wskaż wektor prostopadły do prostej wyrażonej równaniem.

\[ 
p\colon \begin{aligned}[t] x =&1 - 6t, &
\\y =& - 2 + 3t;\ t\in \mathbb{R}
\\ \end{aligned}
\]}
{\((1;2)\)}{\((-6;3)\)}{\((1;-2)\)}{\((2;1)\)}{}{}}
\LANGcs{
\Question{
Z nabízených možností vyberte normálový vektor přímky, která je
vyjádřena parametrickými rovnicemi:

\[ 
p\colon \begin{aligned}[t] x =&1 - 6t, &
\\y =& - 2 + 3t;\ t\in \mathbb{R}
\\ \end{aligned}
\]}
{\((1;2)\)}{\((-6;3)\)}{\((1;-2)\)}{\((2;1)\)}{}{}}
\LANGsk{
\Question{
Z ponúknutých možností vyberte normálový vektor priamky, ktorá je
vyjadrená parametrickými rovnicami:

\[ 
p\colon \begin{aligned}[t] x =&1 - 6t, &
\\y =& - 2 + 3t;\ t\in \mathbb{R}
\\ \end{aligned}
\]}
{\((1;2)\)}{\((-6;3)\)}{\((1;-2)\)}{\((2;1)\)}{}{}}
\LANGes{
\Question{
Halla un vector normal de la recta
\[ 
p\colon \begin{aligned}[t] x =&1 - 6t, &
\\y =& - 2 + 3t;\ t\in \mathbb{R}
\\ \end{aligned}
\]}
{\((1;2)\)}{\((-6;3)\)}{\((1;-2)\)}{\((2;1)\)}{}{}}
\End

\Data\ID{9000107501}
\Updated{2022-04-23 05:04:00}
\Level{A}
\SubArea{Analytic geometry in a plane}
\LANGen{
\Question{
In the following list identify a line which is perpendicular to the line
\( 3x - 2y + 11 = 0\).}
{\(\begin{aligned}[t]  x& = 3t, &
\\y & = 1 - 2t;\  t\in \mathbb{R}
\\ \end{aligned}\)}{\(\begin{aligned}[t]  x& = 1 + 2t, &
\\y & = 2 - 3t;\ t\in \mathbb{R}
\\ \end{aligned}\)}{\(\begin{aligned}[t]  x& = 2 - t, &
\\y & = 3 + t;\  t\in \mathbb{R}
\\ \end{aligned}\)}{\(\begin{aligned}[t] x& = 2 + 3t, &
\\y & = 1 + 2t;\  t\in \mathbb{R}
\\ \end{aligned}\)}{}{}}
\LANGpl{
\Question{
Wskaż prostą, która jest prostopadła do prostej
\(q\colon 3x - 2y + 11 = 0\).}
{\(\begin{aligned}[t] p\colon x& = 3t, &
\\y & = 1 - 2t;\ t\in \mathbb{R}
\\ \end{aligned}\)}{\(\begin{aligned}[t] p\colon x& = 1 + 2t, &
\\y & = 2 - 3t;\ t\in \mathbb{R}
\\ \end{aligned}\)}{\(\begin{aligned}[t] p\colon x& = 2 - t, &
\\y & = 3 + t;\  t\in \mathbb{R}
\\ \end{aligned}\)}{\(\begin{aligned}[t] p\colon x& = 2 + 3t, &
\\y & = 1 + 2t;\ t\in \mathbb{R}
\\ \end{aligned}\)}{}{}}
\LANGcs{
\Question{
Z následujících přímek zadaných parametricky vyberte tu, která je kolmá k
přímce \(q\colon 3x - 2y + 11 = 0\):}
{\(p\colon x = 3t,\ y = 1 - 2t;\ t\in \mathbb{R}\)}{\(p\colon x = 1 + 2t,\ y = 2 - 3t;\ t\in \mathbb{R}\)}{\(p\colon x = 2 - t,\ y = 3 + t;\ t\in \mathbb{R}\)}{\(p\colon x = 2 + 3t,\ y = 1 + 2t;\ t\in \mathbb{R}\)}{}{}}
\LANGsk{
\Question{
Z nasledujúcich priamok zadaných parametricky vyberte tú, ktorá je kolmá k
priamke \(q\colon 3x - 2y + 11 = 0\):}
{\(p\colon x = 3t,\ y = 1 - 2t;\ t\in \mathbb{R}\)}{\(p\colon x = 1 + 2t,\ y = 2 - 3t;\ t\in \mathbb{R}\)}{\(p\colon x = 2 - t,\ y = 3 + t;\ t\in \mathbb{R}\)}{\(p\colon x = 2 + 3t,\ y = 1 + 2t;\ t\in \mathbb{R}\)}{}{}}
\LANGes{
\Question{
En la siguiente lista identifica una recta que es perpendicular a la recta 
\( 3x - 2y + 11 = 0\).}
{\(\begin{aligned}[t] x& = 3t, &
\\y & = 1 - 2t;\  t\in \mathbb{R}
\\ \end{aligned}\)}{\(\begin{aligned}[t]  x& = 1 + 2t, &
\\y & = 2 - 3t;\ t\in \mathbb{R}
\\ \end{aligned}\)}{\(\begin{aligned}[t]  x& = 2 - t, &
\\y & = 3 + t;\  t\in \mathbb{R}
\\ \end{aligned}\)}{\(\begin{aligned}[t] x& = 2 + 3t, &
\\y & = 1 + 2t;\  t\in \mathbb{R}
\\ \end{aligned}\)}{}{}}
\End

\Data\ID{9000107502}
\Updated{2020-07-17 08:07:56}
\Level{A}
\SubArea{Analytic geometry in a plane}
\LANGen{
\Question{
In the following list identify a line which is perpendicular to the line
\(q\).
\[ 
\begin{aligned}q\colon x& = 5 - t,&
\\y & = 3t;\ t\in \mathbb{R}
\\ \end{aligned}
\]}
{\(p\colon x - 3y - 7 = 0\)}{\(p\colon - x - 3y + 11 = 0\)}{\(p\colon 3x - y = 0\)}{\(p\colon y = 5\)}{}{}}
\LANGpl{
\Question{
Wskaż prostą, która jest prostopadła do prostej
\(q\).
\[ 
\begin{aligned}q\colon x& = 5 - t,&
\\y & = 3t;\ t\in \mathbb{R}
\\ \end{aligned}
\]}
{\(p\colon x - 3y - 7 = 0\)}{\(p\colon - x - 3y + 11 = 0\)}{\(p\colon 3x - y = 0\)}{\(p\colon y = 5\)}{}{}}
\LANGcs{
\Question{
Z následujících přímek zadaných obecnými rovnicemi vyberte tu, která
je kolmá k přímce \(q\colon x = 5 - t;\ y = 3t;\ t\in \mathbb{R}\).}
{\(p\colon x - 3y - 7 = 0\)}{\(p\colon - x - 3y + 11 = 0\)}{\(p\colon 3x - y = 0\)}{\(p\colon y = 5\)}{}{}}
\LANGsk{
\Question{
Z nasledujúcich priamok zadaných všeobecnými rovnicami vyberte tú, ktorá
je kolmá k priamke \(q\colon x = 5 - t;\ y = 3t;\ t\in \mathbb{R}\).}
{\(p\colon x - 3y - 7 = 0\)}{\(p\colon - x - 3y + 11 = 0\)}{\(p\colon 3x - y = 0\)}{\(p\colon y = 5\)}{}{}}
\LANGes{
\Question{
En la siguiente lista identifica una recta que es perpendicular a la recta 
\(q\).
\[ 
\begin{aligned}q\colon x& = 5 - t,&
\\y & = 3t;\ t\in \mathbb{R}
\\ \end{aligned}
\]}
{\(p\colon x - 3y - 7 = 0\)}{\(p\colon - x - 3y + 11 = 0\)}{\(p\colon 3x - y = 0\)}{\(p\colon y = 5\)}{}{}}
\End

\Data\ID{9000107503}
\Updated{2022-04-23 05:04:00}
\Level{A}
\SubArea{Analytic geometry in a plane}
\LANGen{
\Question{
Among the lines in the following list (slope-intercept form) identify a line perpendicular
to the line
\[ 
 y = \frac{3}{4}x + 1.
\]}
{\(y = -\frac{4}
{3}x - 2\)}{\(y = -\frac{3}
{4}x - 1\)}{\(y = \frac{4} 
{3}x - 5\)}{\(y = 3\)}{}{}}
\LANGpl{
\Question{
Spośród prostych ( proste w postaci kierunkowej)  wybierz prostą prostopadłą do prostej
 \(q\).
\[ 
 q\colon y = \frac{3} 
{4}x + 1
\]}
{\(p\colon y = -\frac{4}
{3}x - 2\)}{\(p\colon y = -\frac{3}
{4}x - 1\)}{\(p\colon y = \frac{4} 
{3}x - 5\)}{\(p\colon y = 3\)}{}{}}
\LANGcs{
\Question{
Z následujících přímek zadaných rovnicí ve směrnicovém tvaru vyberte tu, která je
kolmá k přímce $q$. \[q\colon y = \frac{3} 
{4}x + 1\]}
{\(p\colon y = -\frac{4}
{3}x - 2\)}{\(p\colon y = -\frac{3}
{4}x - 1\)}{\(p\colon y = \frac{4} 
{3}x - 5\)}{\(p\colon y = 3\)}{}{}}
\LANGsk{
\Question{
Z nasledujúcich priamok zadaných rovnicou v smernicovom tvare vyberte tú, ktorá je
kolmá k priamke $q$. \[q\colon y = \frac{3} 
{4}x + 1\]}
{\(p\colon y = -\frac{4}
{3}x - 2\)}{\(p\colon y = -\frac{3}
{4}x - 1\)}{\(p\colon y = \frac{4} 
{3}x - 5\)}{\(p\colon y = 3\)}{}{}}
\LANGes{
\Question{
Entre las rectas de la siguiente lista (rectas en forma de explícita) identifica una recta perpendicular
a la recta \(q\).
\[ 
 q\colon y = \frac{3} 
{4}x + 1
\]}
{\(p\colon y = -\frac{4}
{3}x - 2\)}{\(p\colon y = -\frac{3}
{4}x - 1\)}{\(p\colon y = \frac{4} 
{3}x - 5\)}{\(p\colon y = 3\)}{}{}}
\End

\Data\ID{9000107510}
\Updated{2020-07-17 04:07:05}
\Level{A}
\SubArea{Analytic geometry in a plane}
\LANGen{
\Question{
In the following list identify a line parallel to the line
\(q\).
\[ 
\begin{aligned}q\colon x& = t, &
\\y & = 1 + 5t;\ t\in \mathbb{R}
\\ \end{aligned}
\]}
{\(p\colon - 5x + y - 13 = 0\)}{\(p\colon x + 5y - 1 = 0\)}{\(p\colon x - 5 = 0\)}{\(p\colon 10x + 2y - 1 = 0\)}{}{}}
\LANGpl{
\Question{
Wskaż prostą równoległą do prostej
\(q\).
\[ 
\begin{aligned}q\colon x& = t, &
\\y & = 1 + 5t;\ t\in \mathbb{R}
\\ \end{aligned}
\]}
{\(p\colon - 5x + y - 13 = 0\)}{\(p\colon x + 5y - 1 = 0\)}{\(p\colon x - 5 = 0\)}{\(p\colon 10x + 2y - 1 = 0\)}{}{}}
\LANGcs{
\Question{
Z následujících přímek zadaných obecnými rovnicemi vyberte tu, která je rovnoběžná
s přímkou \(q\colon x = t,\ y = 1 + 5t;\ t\in \mathbb{R}\):}
{\(p\colon - 5x + y - 13 = 0\)}{\(p\colon x + 5y - 1 = 0\)}{\(p\colon x - 5 = 0\)}{\(p\colon 10x + 2y - 1 = 0\)}{}{}}
\LANGsk{
\Question{
Z následujúcich priamok zadaných všeobecnými rovnicami vyberte tú, ktorá je rovnobežná
s priamkou \(q\colon x = t,\ y = 1 + 5t;\ t\in \mathbb{R}\):}
{\(p\colon - 5x + y - 13 = 0\)}{\(p\colon x + 5y - 1 = 0\)}{\(p\colon x - 5 = 0\)}{\(p\colon 10x + 2y - 1 = 0\)}{}{}}
\LANGes{
\Question{
En la siguiente lista, identifica una recta paralela a la recta
\(q\).
\[ 
\begin{aligned}q\colon x& = t, &
\\y & = 1 + 5t;\ t\in \mathbb{R}
\\ \end{aligned}
\]}
{\(p\colon - 5x + y - 13 = 0\)}{\(p\colon x + 5y - 1 = 0\)}{\(p\colon x - 5 = 0\)}{\(p\colon 10x + 2y - 1 = 0\)}{}{}}
\End

\Data\ID{9000151310}
\Updated{2021-12-05 07:12:38}
\Level{A}
\SubArea{Analytic geometry in a plane}
\LANGen{
\Question{
Find the value of the real parameter
\(a\) which ensures, that
the following two lines \(p\) 
and \(q\) 
are perpendicular. 
\[ 
 p\colon ax + y - 4 = 0,\qquad q\colon x + 2y + 4 = 0.
\]}
{\(- 2\)}{\(2\)}{\(1\)}{\(- 1\)}{}{}}
\LANGpl{
\Question{
Wyznacz wartość rzeczywistą parametru 
\(a\) tak, aby proste  \(p\) 
i \(q\) 
były prostopadłe. 
\[ 
 p\colon ax + y - 4 = 0,\qquad q\colon x + 2y + 4 = 0.
\]}
{\(- 2\)}{\(2\)}{\(1\)}{\(- 1\)}{}{}}
\LANGcs{
\Question{
Jsou dány dvě přímky \(p\),
\(q\) 
zadané obecnými rovnicemi takto: 
\[ 
 p\colon ax + y - 4 = 0,\qquad q\colon x + 2y + 4 = 0.
\]
Určete hodnotu parametru \(a\in \mathbb{R}\) 
tak, aby přímky \(p\),
\(q\) byly
navzájem kolmé.}
{\(- 2\)}{\(2\)}{\(1\)}{\(- 1\)}{}{}}
\LANGsk{
\Question{
Sú dané dve priamky \(p\),
\(q\) 
zadané všeobecnými rovnicami takto: 
\[ 
 p\colon ax + y - 4 = 0,\qquad q\colon x + 2y + 4 = 0.
\]
Určte hodnotu parametra \(a\in \mathbb{R}\) 
tak, aby priamky \(p\),
\(q\) boli navzájom kolmé.}
{\(- 2\)}{\(2\)}{\(1\)}{\(- 1\)}{}{}}
\LANGes{
\Question{
Halla el valor del parámetro real
\(a\) para que las rectas \(p\) 
y \(q\) 
sean perpendiculares. 
\[ 
 p\colon ax + y - 4 = 0,\qquad q\colon x + 2y + 4 = 0.
\]}
{\(- 2\)}{\(2\)}{\(1\)}{\(- 1\)}{}{}}
\End

\Data\ID{1003090802}
\Updated{2022-04-23 05:04:00}
\Level{B}
\SubArea{Analytic geometry in a plane}
\LANGen{
\Question{
Find the distance between parallel lines \( p \) and \( q \), if they are given by their general form equations, where \( p \) is \( 2x-4y+5=0 \) and \( q \) is \( x-2y+3=0 \).}
{\( \frac{\sqrt5}{10} \)}{\( \frac{11\sqrt5}{10} \)}{\( \frac{3}{2\sqrt5} \)}{\( \frac{3\sqrt5}{10} \)}{}{}}
\LANGpl{
\Question{
Oblicz odległość pomiędzy prostymi równoległymi \( p \) i \( q \), jeśli proste są określone równaniami, gdzie \( p \) to \( 2x-4y+5=0 \) a \( q \) to \( x-2y+3=0 \).}
{\( \frac{\sqrt5}{10} \)}{\( \frac{11\sqrt5}{10} \)}{\( \frac{3}{2\sqrt5} \)}{\( \frac{3\sqrt5}{10} \)}{}{}}
\LANGcs{
\Question{
Vypočtěte vzdálenost rovnoběžek \( p \), \( q \), jsou-li zadány jejich obecné rovnice: \( p \) :  \( 2x-4y+5=0 \), \( q \) :  \( x-2y+3=0 \).}
{\( \frac{\sqrt5}{10} \)}{\( \frac{11\sqrt5}{10} \)}{\( \frac{3}{2\sqrt5} \)}{\( \frac{3\sqrt5}{10} \)}{}{}}
\LANGsk{
\Question{
Vypočítajte vzdialenosť rovnobežiek \( p \), \( q \), ak sú zadané ich všeobecné rovnice: \( p \) :  \( 2x-4y+5=0 \), \( q \) :  \( x-2y+3=0 \).}
{\( \frac{\sqrt5}{10} \)}{\( \frac{11\sqrt5}{10} \)}{\( \frac{3}{2\sqrt5} \)}{\( \frac{3\sqrt5}{10} \)}{}{}}
\LANGes{
\Question{
Halla la distancia entre las rectas paralelas \( p \) y \( q \), suponiendo que la ecuación general  de la recta \( p \)  es \( 2x-4y+5=0 \)  y de la recta  \( q \) es \( x-2y+3=0 \).}
{\( \frac{\sqrt5}{10} \)}{\( \frac{11\sqrt5}{10} \)}{\( \frac{3}{2\sqrt5} \)}{\( \frac{3\sqrt5}{10} \)}{}{}}
\End

\Data\ID{1003090803}
\Updated{2020-07-15 03:07:21}
\Level{B}
\SubArea{Analytic geometry in a plane}
\LANGen{
\Question{
Find the distance between parallel lines \( p \) and \( q \), if they are given by slope-intercept form equations, where \( p \) is \( y=-3x+5 \) and \( q \) is \( y=-3x-1 \).}
{\( \frac{3\sqrt{10}}5 \)}{\( \frac{2\sqrt{10}}5 \)}{\( \frac{4\sqrt{10}}5 \)}{\( \frac{\sqrt{10}}5 \)}{}{}}
\LANGpl{
\Question{
Oblicz odległość pomiędzy prostymi równoległymi \( p \) i \( q \), określonymi równaniami kierunkowymi, gdzie \( p \) to \( y=-3x+5 \) a \( q \) to \( y=-3x-1 \).}
{\( \frac{3\sqrt{10}}5 \)}{\( \frac{2\sqrt{10}}5 \)}{\( \frac{4\sqrt{10}}5 \)}{\( \frac{\sqrt{10}}5 \)}{}{}}
\LANGcs{
\Question{
Vypočtěte vzdálenost rovnoběžek \( p \), \( q \), jsou-li zadány jejich rovnice ve směrnicovém tvaru: \( p \) : \( y=-3x+5 \),  \( q \) : \( y=-3x-1 \).}
{\( \frac{3\sqrt{10}}5 \)}{\( \frac{2\sqrt{10}}5 \)}{\( \frac{4\sqrt{10}}5 \)}{\( \frac{\sqrt{10}}5 \)}{}{}}
\LANGsk{
\Question{
Vypočítajte vzdialenosť rovnobežiek \( p \), \( q \), ak sú zadané ich rovnice v smernicovom tvare: \( p \) : \( y=-3x+5 \),  \( q \) : \( y=-3x-1 \).}
{\( \frac{3\sqrt{10}}5 \)}{\( \frac{2\sqrt{10}}5 \)}{\( \frac{4\sqrt{10}}5 \)}{\( \frac{\sqrt{10}}5 \)}{}{}}
\LANGes{
\Question{
Halla la distancia entre las rectas paralelas \( p \) y \( q \), suponiendo que la ecuación general de la recta \( p \) es \( y=-3x+5 \) y de la recta \( q \) es \( y=-3x-1 \).}
{\( \frac{3\sqrt{10}}5 \)}{\( \frac{2\sqrt{10}}5 \)}{\( \frac{4\sqrt{10}}5 \)}{\( \frac{\sqrt{10}}5 \)}{}{}}
\End

\Data\ID{1003090804}
\Updated{2022-04-23 05:04:00}
\Level{B}
\SubArea{Analytic geometry in a plane}
\LANGen{
\Question{
Find the distance between parallel lines \( p \) and \( q \) given by their parametric equations.
\begin{align*}
 p\colon x&=3+3t, & q\colon x&=2-3s, \\
y&=-1+t;\ t\in\mathbb{R}; & y&=1-s;\ s\in\mathbb{R}.
\end{align*}}
{\( \frac{7\sqrt{10}}{10} \)}{\( \frac{\sqrt{10}}{2} \)}{\( \frac{\sqrt{10}}{5} \)}{\( \frac{5\sqrt{10}}{2} \)}{}{}}
\LANGpl{
\Question{
Oblicz odległość pomiędzy prostymi równoległymi  \( p \) i \( q \) określonymi równaniami parametrycznymi.
\begin{align*}
 p\colon x&=3+3t, & q\colon x&=2-3s, \\
y&=-1+t;\ t\in\mathbb{R}; & y&=1-s;\ s\in\mathbb{R}.
\end{align*}}
{\( \frac{7\sqrt{10}}{10} \)}{\( \frac{\sqrt{10}}{2} \)}{\( \frac{\sqrt{10}}{5} \)}{\( \frac{5\sqrt{10}}{2} \)}{}{}}
\LANGcs{
\Question{
Vypočtěte vzdálenost rovnoběžek \( p \), \( q \) jsou-li zadána jejich parametrická vyjádření:
\begin{align*}
 p\colon x&=3+3t, & q\colon x&=2-3s, \\
y&=-1+t;\ t\in\mathbb{R}; & y&=1-s;\ s\in\mathbb{R}.
\end{align*}}
{\( \frac{7\sqrt{10}}{10} \)}{\( \frac{\sqrt{10}}{2} \)}{\( \frac{\sqrt{10}}{5} \)}{\( \frac{5\sqrt{10}}{2} \)}{}{}}
\LANGsk{
\Question{
Nájdi vzdialenosť dvoch priamok \( p \) a \( q \), ktoré sú dané parametrickými rovnicami.
\begin{align*}
 p\colon x&=3+3t, & q\colon x&=2-3s, \\
y&=-1+t;\ t\in\mathbb{R}; & y&=1-s;\ s\in\mathbb{R}.
\end{align*}}
{\( \frac{7\sqrt{10}}{10} \)}{\( \frac{\sqrt{10}}{2} \)}{\( \frac{\sqrt{10}}{5} \)}{\( \frac{5\sqrt{10}}{2} \)}{}{}}
\LANGes{
\Question{
Halla la distancia entre las rectas paralelas \( p \) y \( q \) que vienen dadas por sus ecuaciones parámetricas.
\begin{align*}
 p\colon x&=3+3t, & q\colon x&=2-3s, \\
y&=-1+t;\ t\in\mathbb{R}; & y&=1-s;\ s\in\mathbb{R}.
\end{align*}}
{\( \frac{7\sqrt{10}}{10} \)}{\( \frac{\sqrt{10}}{2} \)}{\( \frac{\sqrt{10}}{5} \)}{\( \frac{5\sqrt{10}}{2} \)}{}{}}
\End

\Data\ID{1103061301}
\Updated{2021-12-05 07:12:03}
\Level{B}
\SubArea{Analytic geometry in a plane}
\IMG{1103061301.pdf}
\LANGen{
\Question{
Let \( ABC \) be a triangle (see the picture). Find the standard form equations of the lines \( t \), \( v \) and \( o \), where \( t \) contains the median to \( AB \), \( v \) contains the altitude to \( AB \) and \( o \) is the line of symmetry of \( AB \). 
Choose the option with all equations correct.}
{\( t\colon 2x+y-10=0	;\ v\colon 4x+y-16=0;\ o\colon 4x+y-20=0 \)}{\( t\colon 2x+y-10=0;\ v\colon x-4y+13=0;\ o\colon x-4y-5=0 \)}{\( t\colon x-2y-5=0;\ v\colon 4x+y-16=0;\ o\colon 4x+y-20=0 \)}{\( t\colon x-2y-5=0;\ v\colon x-4y+13=0;\ o\colon x-4y-5=0 \)}{}{}}
\LANGpl{
\Question{
Dany jest trójkąt \( ABC \)  (spójrz na rysunek). Wskaż równania prostych  \( t \), \( v \) i \( o \), jeśłi \( t \) to środkowa \( AB \), \( v \)  wysokość do \( AB \) a \( o \) to symetralna \( AB \). 
Wybierz opcję, w której wszystkie równania są poprawne.}
{\( t\colon 2x+y-10=0	;\ v\colon 4x+y-16=0;\ o\colon 4x+y-20=0 \)}{\( t\colon 2x+y-10=0;\ v\colon x-4y+13=0;\ o\colon x-4y-5=0 \)}{\( t\colon x-2y-5=0;\ v\colon 4x+y-16=0;\ o\colon 4x+y-20=0 \)}{\( t\colon x-2y-5=0;\ v\colon x-4y+13=0;\ o\colon x-4y-5=0 \)}{}{}}
\LANGcs{
\Question{
Je dán trojúhelník \( ABC \) (viz obrázek). Určete obecné rovnice přímek \( t \), \( v \), \( o \), kde \( t \) je těžnice na stranu \( AB \), \( v \) je přímka, na které leží výška na stranu \( AB \) a přímka \( o \) je osa strany \( AB \). 
Vyberte možnost, kde jsou všechny tři rovnice správně.}
{\( t\colon 2x+y-10=0	;\ v\colon 4x+y-16=0;\ o\colon 4x+y-20=0 \)}{\( t\colon 2x+y-10=0;\ v\colon x-4y+13=0;\ o\colon x-4y-5=0 \)}{\( t\colon x-2y-5=0;\ v\colon 4x+y-16=0;\ o\colon 4x+y-20=0 \)}{\( t\colon x-2y-5=0;\ v\colon x-4y+13=0;\ o\colon x-4y-5=0 \)}{}{}}
\LANGsk{
\Question{
Je daný trojuholník \( ABC \) (viď obrázok). Určte všeobecné rovnice priamok \( t \), \( v \), \( o \), kde \( t \) je ťažnica na stranu \( AB \), \( v \) je priamka, na ktorej leží výška na stranu \( AB \) a priamka \( o \) je os strany \( AB \). 
Vyberte možnosť, kde sú všetky tri rovnice správne.}
{\( t\colon 2x+y-10=0	;\ v\colon 4x+y-16=0;\ o\colon 4x+y-20=0 \)}{\( t\colon 2x+y-10=0;\ v\colon x-4y+13=0;\ o\colon x-4y-5=0 \)}{\( t\colon x-2y-5=0;\ v\colon 4x+y-16=0;\ o\colon 4x+y-20=0 \)}{\( t\colon x-2y-5=0;\ v\colon x-4y+13=0;\ o\colon x-4y-5=0 \)}{}{}}
\LANGes{
\Question{
Dado el triángulo \( ABC \)  (observa la imagen). Determina las ecuaciones generales de las rectas \( t \), \( v \) y \( o \), donde \( t \) es la mediana del lado \( AB \), \( v \) es la altura del lado \( AB \) y \( o \) es el eje de simetría del lado \( AB \). 
Elige la opción donde todas las ecuaciones son correctas.}
{\( t\colon 2x+y-10=0	;\ v\colon 4x+y-16=0;\ o\colon 4x+y-20=0 \)}{\( t\colon 2x+y-10=0;\ v\colon x-4y+13=0;\ o\colon x-4y-5=0 \)}{\( t\colon x-2y-5=0;\ v\colon 4x+y-16=0;\ o\colon 4x+y-20=0 \)}{\( t\colon x-2y-5=0;\ v\colon x-4y+13=0;\ o\colon x-4y-5=0 \)}{}{}}
\End

\Data\ID{1103090801}
\Updated{2022-04-23 05:04:00}
\Level{B}
\SubArea{Analytic geometry in a plane}
\IMG{1103090801.pdf}
\LANGen{
\Question{
Find a general form equation of the straight line that passes through the point \( M=[2;3] \) and is parallel with the line of symmetry of the line segment \( AB \), where \( A=[-1;4] \), and \( B=\left[\frac52;-3\right] \) (see the picture).}
{\( x-2y+4=0 \)}{\( 2x+y-7=0 \)}{\( 3x+2y-12=0 \)}{\( 2x-3y+5=0 \)}{}{}}
\LANGpl{
\Question{
Wskaż równanie prostej przechodzącej przez punkt \( M=[2;3] \) oraz równoległej do symetralnej odcinka \( AB \), gdzie \( A=[-1;4] \), i \( B=\left[\frac52;-3\right] \) (spójrz na rysunek).}
{\( x-2y+4=0 \)}{\( 2x+y-7=0 \)}{\( 3x+2y-12=0 \)}{\( 2x-3y+5=0 \)}{}{}}
\LANGcs{
\Question{
Určete obecnou rovnici přímky, která prochází bodem \( M=[2;3] \) a je rovnoběžná s osou úsečky \( AB \), přičemž \( A=[-1;4] \) a \( B=\left[\frac52;-3\right] \) (viz obrázek).}
{\( x-2y+4=0 \)}{\( 2x+y-7=0 \)}{\( 3x+2y-12=0 \)}{\( 2x-3y+5=0 \)}{}{}}
\LANGsk{
\Question{
Určte všeobecnú rovnicu priamky, ktorá prechádza bodom \( M=[2;3] \) a je rovnobežná s osou úsečky \( AB \), pričom \( A=[-1;4] \) a \( B=\left[\frac52;-3\right] \) (viď obrázok).}
{\( x-2y+4=0 \)}{\( 2x+y-7=0 \)}{\( 3x+2y-12=0 \)}{\( 2x-3y+5=0 \)}{}{}}
\LANGes{
\Question{
Determina la ecuación general de la recta que pasa por el punto \( M=[2;3] \) y es paralela al eje de siemtría del segmento \( AB  \), donde \( A=[-1;4] \) y \( B=\left[\frac52;-3\right] \) (mira la imagen).}
{\( x-2y+4=0 \)}{\( 2x+y-7=0 \)}{\( 3x+2y-12=0 \)}{\( 2x-3y+5=0 \)}{}{}}
\End

\Data\ID{1103090805}
\Updated{2021-12-05 07:12:37}
\Level{B}
\SubArea{Analytic geometry in a plane}
\IMG{1103090805.pdf}
\LANGen{
\Question{
Find the straight lines passing through the coordinate origin at the distance of \( 2 \) from the point \( M=[0;4] \). Express their equations in the slope-intercept form.}
{\( y=\pm\sqrt3 x \)}{\( y=\pm4 x \)}{\( y=\pm\frac32 x \)}{\( y=\pm2\sqrt3 x \)}{}{}}
\LANGpl{
\Question{
Wyznacz proste przechodzące przez początek układu współrzędnych w odległości \( 2 \) od punktu \( M=[0;4] \). Wskaż równanie w postaci kierunkowej.}
{\( y=\pm\sqrt3 x \)}{\( y=\pm4 x \)}{\( y=\pm\frac32 x \)}{\( y=\pm2\sqrt3 x \)}{}{}}
\LANGcs{
\Question{
Najděte přímky, které procházejí počátkem soustavy souřadnic a mají od bodu \( M=[0;4] \) vzdálenost \( 2 \). Rovnice přímek vyjádřete ve směrnicovém tvaru.}
{\( y=\pm\sqrt3 x \)}{\( y=\pm4 x \)}{\( y=\pm\frac32 x \)}{\( y=\pm2\sqrt3 x \)}{}{}}
\LANGsk{
\Question{
Nájdite priamky, ktoré prechádzajú začiatkom sústavy súradníc a majú od bodu \( M=[0;4] \) vzdialenosť \( 2 \). Rovnice priamok vyjadrite v smernicovom tvare.}
{\( y=\pm\sqrt3 x \)}{\( y=\pm4 x \)}{\( y=\pm\frac32 x \)}{\( y=\pm2\sqrt3 x \)}{}{}}
\LANGes{
\Question{
Halla las rectas que pasan por el origen de coordenadas  y distan \( 2 \) del punto \( M=[0;4] \). Expresa sus ecuaciones en la forma explícita.}
{\( y=\pm\sqrt3 x \)}{\( y=\pm4 x \)}{\( y=\pm\frac32 x \)}{\( y=\pm2\sqrt3 x \)}{}{}}
\End

\Data\ID{1103109001}
\Updated{2021-12-05 07:12:36}
\Level{B}
\SubArea{Analytic geometry in a plane}
\IMG{1103109001.pdf}
\LANGen{
\Question{
Let \( A \) be the point \( [4;3] \) and let the line \( p \) has the equation \( x-y+3=0 \). Find the coordinates of the point \( A' \) which is a mirror reflection of \( A \) about the line of symmetry \( p \) (see the picture).}
{\( A'=[0;7] \)}{\( A'=[1;8] \)}{\( A'=[-1;8] \)}{\( A'=[-1;7] \)}{}{}}
\LANGpl{
\Question{
Dany jest punkt \( A \) o współrzędnych \( [4;3] \) oraz prosta \( p \) o równaniu \( x-y+3=0 \). Wyznacz współrzędne punktu \( A' \), który jest obrazem punktu \( A \) w symetrii osiowej  względem prostej \( p \) (spójrz na rysunek).}
{\( A'=[0;7] \)}{\( A'=[1;8] \)}{\( A'=[-1;8] \)}{\( A'=[-1;7] \)}{}{}}
\LANGcs{
\Question{
Je dán bod \( A \) =\( [4;3] \) a přímka \( p \): \( x-y+3=0 \). Určete souřadnice bodu \( A' \), který je obrazem bodu \( A \) v osové souměrnosti s osou \( p \) (viz obrázek).}
{\( A'=[0;7] \)}{\( A'=[1;8] \)}{\( A'=[-1;8] \)}{\( A'=[-1;7] \)}{}{}}
\LANGsk{
\Question{
Je daný bod \( A \) =\( [4;3] \) a priamka \( p \): \( x-y+3=0 \). Určte súradnice bodu \( A' \), ktorý je obrazom bodu \( A \) v osovej súmernosti s osou \( p \) (viď obrázok).}
{\( A'=[0;7] \)}{\( A'=[1;8] \)}{\( A'=[-1;8] \)}{\( A'=[-1;7] \)}{}{}}
\LANGes{
\Question{
Sea \( A \) el punto \( [4;3] \) y sea \( p \) la recta que tiene por ecuación \( x-y+3=0 \). Halla las coordenadas del punto \( A' \) que es el simétrico de \( A \) respecto el eje de simetría   \( p \) (mira la imagen).}
{\( A'=[0;7] \)}{\( A'=[1;8] \)}{\( A'=[-1;8] \)}{\( A'=[-1;7] \)}{}{}}
\End

\Data\ID{1103109002}
\Updated{2021-12-05 07:12:01}
\Level{B}
\SubArea{Analytic geometry in a plane}
\IMG{1103109002.pdf}
\LANGen{
\Question{
Let \( A=[0;1] \), \( B=[4;-2] \) and \( S=[4;3] \) be the points (see the picture). Find the coordinates of the points \( C \) and \( D \) so that \( ABCD \) is a parallelogram with the centre \( S \).}
{\( C=[8;5]\text{, } D=[4;8] \)}{\( C=[7;5]\text{, } D=[4;8] \)}{\( C=[8;5]\text{, } D=[4;7] \)}{\( C=[4;8]\text{, } D=[8;5] \)}{}{}}
\LANGpl{
\Question{
Dane są punkty  \( A=[0;1] \), \( B=[4;-2] \) and \( S=[4;3] \)  (spójrz na rysunek). Wyznacz współrzędne punktu  \( C \) i \( D \) tak, aby powstał równoległobok \( ABCD \) o środku  \( S \).}
{\( C=[8;5]\text{, } D=[4;8] \)}{\( C=[7;5]\text{, } D=[4;8] \)}{\( C=[8;5]\text{, } D=[4;7] \)}{\( C=[4;8]\text{, } D=[8;5] \)}{}{}}
\LANGcs{
\Question{
Jsou dány body \( A=[0;1] \), \( B=[4;-2] \) a \( S=[4;3] \) (viz obrázek). Určete souřadnice bodů \( C \) a \( D \) tak, aby \( ABCD \) byl rovnoběžník se středem \( S \).}
{\( C=[8;5]\text{, } D=[4;8] \)}{\( C=[7;5]\text{, } D=[4;8] \)}{\( C=[8;5]\text{, } D=[4;7] \)}{\( C=[4;8]\text{, } D=[8;5] \)}{}{}}
\LANGsk{
\Question{
Sú dané body \( A=[0;1] \), \( B=[4;-2] \) a \( S=[4;3] \) (viď obrázok). Určte súradnice bodov \( C \) a \( D \) tak, aby \( ABCD \) bol rovnobežník so stredom \( S \).}
{\( C=[8;5]\text{, } D=[4;8] \)}{\( C=[7;5]\text{, } D=[4;8] \)}{\( C=[8;5]\text{, } D=[4;7] \)}{\( C=[4;8]\text{, } D=[8;5] \)}{}{}}
\LANGes{
\Question{
Sean \( A=[0;1] \), \( B=[4;-2] \) y \( S=[4;3] \) los puntos (mira la imagen). Determina las coordenadas de los puntos \( C \) y \( D \) para que \( ABCD \) sea un paralelogramo con el centro en \( S \).}
{\( C=[8;5]\text{, } D=[4;8] \)}{\( C=[7;5]\text{, } D=[4;8] \)}{\( C=[8;5]\text{, } D=[4;7] \)}{\( C=[4;8]\text{, } D=[8;5] \)}{}{}}
\End

\Data\ID{1103109003}
\Updated{2021-12-05 07:12:47}
\Level{B}
\SubArea{Analytic geometry in a plane}
\IMG{1103109003.pdf}
\LANGen{
\Question{
Let \( 2x+6y-5=0 \) be the line \( p \) and \( x+3y-4=0 \) be the line \( o \), where \( p \) and \( o \) are parallel (see the picture). Find the general form equation of a line \( p' \) which is the reflection of the line \( p \) about the line of symmetry \( o \).}
{\( p'\colon 2x+6y-11=0 \)}{\( p'\colon 2x+6y-2=0 \)}{\( p'\colon 2x+6y+5=0 \)}{\( p'\colon -2x-6y-11=0 \)}{}{}}
\LANGpl{
\Question{
Niech \( 2x+6y-5=0 \) będzie prostą \( p \) a \( x+3y-4=0 \) prostą \( o \), \( p \) i  \( o \) są równoległe (spójrz na rysunek). Wyznacz równanie prostej \( p' \), która jest obrazem prostej  \( p \) w symetrii osiowej względem prostej  \( o \).}
{\( p'\colon 2x+6y-11=0 \)}{\( p'\colon 2x+6y-2=0 \)}{\( p'\colon 2x+6y+5=0 \)}{\( p'\colon -2x-6y-11=0 \)}{}{}}
\LANGcs{
\Question{
Jsou dány rovnoběžky \( p \): \( 2x+6y-5=0 \)  a  \( o \): \( x+3y-4=0 \) (viz obrázek). Určete obecnou rovnici přímky \( p' \), která je obrazem přímky \( p \) v osové souměrnosti s osou \( o \).}
{\( p'\colon 2x+6y-11=0 \)}{\( p'\colon 2x+6y-2=0 \)}{\( p'\colon 2x+6y+5=0 \)}{\( p'\colon -2x-6y-11=0 \)}{}{}}
\LANGsk{
\Question{
Sú dané rovnobežky \( p \): \( 2x+6y-5=0 \)  a  \( o \): \( x+3y-4=0 \) (viď obrázok). Určte všeobecnú rovnicu priamky \( p' \), ktorá je obrazom priamky \( p \) v osovej súmernosti s osou \( o \).}
{\( p'\colon 2x+6y-11=0 \)}{\( p'\colon 2x+6y-2=0 \)}{\( p'\colon 2x+6y+5=0 \)}{\( p'\colon -2x-6y-11=0 \)}{}{}}
\LANGes{
\Question{
Sean \( 2x+6y-5=0 \) la recta \( p \) y \( x+3y-4=0 \) la recta \( o \), donde \( p \) y \( o \) son paralelas (mira la imagen). Determina la ecuación general de una recta \( p' \) que es simétrica de la recta \( p \) respecto al eje de simetría \( o \).}
{\( p'\colon 2x+6y-11=0 \)}{\( p'\colon 2x+6y-2=0 \)}{\( p'\colon 2x+6y+5=0 \)}{\( p'\colon -2x-6y-11=0 \)}{}{}}
\End

\Data\ID{1103109004}
\Updated{2021-12-05 07:12:38}
\Level{B}
\SubArea{Analytic geometry in a plane}
\IMG{1103109004.pdf}
\LANGen{
\Question{
Let \( p \) be the line with the equation \( x-2y-1=0 \) and let \( S \) be the point with coordinates \( [2;2] \) (see the picture). Find the general form equation of the line \( p' \) which is the image of the line \( p \) in the point symmetry with the centre in \( S \).}
{\( p'\colon x-2y+5=0 \)}{\( p'\colon 2x-4y+9=0 \)}{\( p'\colon x-2y+4=0 \)}{\( p'\colon x-2y+6=0 \)}{}{}}
\LANGpl{
\Question{
Dana jest prosta  \( p \) określona równaniem \( x-2y-1=0 \) oraz punkt \( S \) o współrzędnych \( [2;2] \) (spójrz na rysunek). Wyznacz równanie prostej  \( p' \), która jest obrazem prostej  \( p \) w symetrii środkowej względem punktu  \( S \).}
{\( p'\colon x-2y+5=0 \)}{\( p'\colon 2x-4y+9=0 \)}{\( p'\colon x-2y+4=0 \)}{\( p'\colon x-2y+6=0 \)}{}{}}
\LANGcs{
\Question{
Je dána přímka \( p \): \( x-2y-1=0 \) a bod \( S \) =\( [2;2] \) (viz obrázek). Určete obecnou rovnici přímky \( p' \), která je obrazem přímky \( p \) ve středové souměrnosti se středem \( S \).}
{\( p'\colon x-2y+5=0 \)}{\( p'\colon 2x-4y+9=0 \)}{\( p'\colon x-2y+4=0 \)}{\( p'\colon x-2y+6=0 \)}{}{}}
\LANGsk{
\Question{
Je daná priamka \( p \): \( x-2y-1=0 \) a bod \( S \) =\( [2;2] \) (viď obrázok). Určte všeobecnú rovnicu priamky \( p' \), ktorá je obrazom priamky \( p \) v stredovej súmernosti so stredom \( S \).}
{\( p'\colon x-2y+5=0 \)}{\( p'\colon 2x-4y+9=0 \)}{\( p'\colon x-2y+4=0 \)}{\( p'\colon x-2y+6=0 \)}{}{}}
\LANGes{
\Question{
Sea \( p \) la recta con ecuación \( x-2y-1=0 \) y sea \( S \) un punto con coordenadas \( [2;2] \) (mira la imagen). Determina la ecuación general de la recta \( p' \)  que es simétrica de la recta \( p \) respecto al punto \( S \).}
{\( p'\colon x-2y+5=0 \)}{\( p'\colon 2x-4y+9=0 \)}{\( p'\colon x-2y+4=0 \)}{\( p'\colon x-2y+6=0 \)}{}{}}
\End

\Data\ID{1103109005}
\Updated{2021-12-05 10:12:11}
\Level{B}
\SubArea{Analytic geometry in a plane}
\IMG{1103109005.pdf}
\LANGen{
\Question{
Let \( p \) be the line with the equation \( x-2y+5=0 \) and let \( \vec{v} \) be the vector \( (3;-2) \) (see the picture). Find the general form equation of the line \( p' \) which is the image of the line \( p \) translated by the vector \( \vec{v} \).}
{\( p'\colon x-2y-2=0 \)}{\( p'\colon 2x-4y-3=0 \)}{\( p'\colon x-2y-1=0 \)}{\( p'\colon 2x-4y+3=0 \)}{}{}}
\LANGpl{
\Question{
Dana jest prosta \( p \) określona równaniem  \( x-2y+5=0 \) oraz wektor \( \vec{v} \) o współrzędnych  \( (3;-2) \) (spójrz na rysunek). Wyznacz równanie prostej  \( p' \), która jest obrazem prostej \( p \) w przesunięciu o wektor \( \vec{v} \).}
{\( p'\colon x-2y-2=0 \)}{\( p'\colon 2x-4y-3=0 \)}{\( p'\colon x-2y-1=0 \)}{\( p'\colon 2x-4y+3=0 \)}{}{}}
\LANGcs{
\Question{
Je dána přímka \( p \): \( x-2y+5=0 \) a vektor \( \vec{v} \) = \( (3;-2) \) (viz obrázek). Určete obecnou rovnici přímky \( p' \), která je obrazem přímky \( p \) v posunutí daném vektorem \( \vec{v} \).}
{\( p'\colon x-2y-2=0 \)}{\( p'\colon 2x-4y-3=0 \)}{\( p'\colon x-2y-1=0 \)}{\( p'\colon 2x-4y+3=0 \)}{}{}}
\LANGsk{
\Question{
Je daná priamka \( p \): \( x-2y+5=0 \) a vektor \( \vec{v} \) = \( (3;-2) \) (viď obrázok). Určte všeobecnú rovnicu priamky \( p' \), ktorá je obrazom priamky \( p \) v posunutí danom vektorom \( \vec{v} \).}
{\( p'\colon x-2y-2=0 \)}{\( p'\colon 2x-4y-3=0 \)}{\( p'\colon x-2y-1=0 \)}{\( p'\colon 2x-4y+3=0 \)}{}{}}
\LANGes{
\Question{
Sea \( p \) la recta con ecuación \( x-2y+5=0 \) y sea \( \vec{v} \) el vector \( (3;-2) \) (mira la imagen). Determina la ecuación general de la recta \( p' \) que es una traslación de la recta \( p \) mediante el vector \( \vec{v} \).}
{\( p'\colon x-2y-2=0 \)}{\( p'\colon 2x-4y-3=0 \)}{\( p'\colon x-2y-1=0 \)}{\( p'\colon 2x-4y+3=0 \)}{}{}}
\End

\Data\ID{1103109006}
\Updated{2022-04-23 05:04:00}
\Level{B}
\SubArea{Analytic geometry in a plane}
\IMG{1103109006.pdf}
\LANGen{
\Question{
Let \( p \) be the line with the equation \( x-2y-1=0 \). Find the general form equations of all lines parallel to \( p \) such that their distance from \( p \) equals to \( \sqrt5 \).}
{\( x-2y+4=0;\ x-2y-6=0 \)}{\( x-2y+\sqrt5=0;\ x-2y-\sqrt5=0 \)}{\( x-2y-1+\sqrt5=0;\ x-2y-1-\sqrt5=0 \)}{\( x-2y+6=0;\ x-2y-4=0 \)}{}{}}
\LANGpl{
\Question{
Dana jest prosta \( p \) określona równaniem \( x-2y-1=0 \). Wyznacz równania wszystkich prostych równoległych do \( p \) tak, aby odległość od \( p \) była równa \( \sqrt5 \).}
{\( x-2y+4=0;\ x-2y-6=0 \)}{\( x-2y+\sqrt5=0;\ x-2y-\sqrt5=0 \)}{\( x-2y-1+\sqrt5=0;\ x-2y-1-\sqrt5=0 \)}{\( x-2y+6=0;\ x-2y-4=0 \)}{}{}}
\LANGcs{
\Question{
Je dána přímka \( p \): \( x-2y-1=0 \). Určete rovnice všech přímek rovnoběžných s přímkou \( p \), které mají od ní vzdálenost \( \sqrt5 \)  (viz obrázek).}
{\( x-2y+4=0;\ x-2y-6=0 \)}{\( x-2y+\sqrt5=0;\ x-2y-\sqrt5=0 \)}{\( x-2y-1+\sqrt5=0;\ x-2y-1-\sqrt5=0 \)}{\( x-2y+6=0;\ x-2y-4=0 \)}{}{}}
\LANGsk{
\Question{
Je daná priamka \( p \): \( x-2y-1=0 \). Určte rovnicu všetkých priamok rovnobežných s priamkou \( p \), ktoré majú od nej vzdialenosť \( \sqrt5 \)  (viď obrázok).}
{\( x-2y+4=0;\ x-2y-6=0 \)}{\( x-2y+\sqrt5=0;\ x-2y-\sqrt5=0 \)}{\( x-2y-1+\sqrt5=0;\ x-2y-1-\sqrt5=0 \)}{\( x-2y+6=0;\ x-2y-4=0 \)}{}{}}
\LANGes{
\Question{
Sea \( p \)  la recta con ecuación \( x-2y-1=0 \). Determina las ecuaciones generales de todas las rectas paralelas a la recta \( p \) suponiendo que la distancia a \( p \) equivale a \( \sqrt5 \).}
{\( x-2y+4=0;\ x-2y-6=0 \)}{\( x-2y+\sqrt5=0;\ x-2y-\sqrt5=0 \)}{\( x-2y-1+\sqrt5=0;\ x-2y-1-\sqrt5=0 \)}{\( x-2y+6=0;\ x-2y-4=0 \)}{}{}}
\End

\Data\ID{1103109007}
\Updated{2021-12-05 10:12:01}
\Level{B}
\SubArea{Analytic geometry in a plane}
\IMG{1103109007.pdf}
\LANGen{
\Question{
Let \( p \) be the line with the equation \( x-2y-1=0 \). Find the coordinates of all points lying on the line \( p \) such that their distance from the line \( x=4 \) equals to \( 2 \).}
{\( X_1 = \left[2;\frac12\right]\text{, }X_2 = \left[6;\frac52\right] \)}{\( X_1 = \left[2;1\right]\text{, }X_2 = \left[6;5\right] \)}{\( X_1 = \left[2;\frac14\right]\text{, }X_2 = \left[6;\frac54\right] \)}{\( X_1 = \left[2;\frac32\right]\text{, }X_2 = \left[6;\frac72\right] \)}{}{}}
\LANGpl{
\Question{
Dana jest prosta \( p \) o równaniu  \( x-2y-1=0 \).  Wyznacz współrzędne punktów leżących na prostej \( p \) tak, aby odległość od prostej \( x=4 \) była  równa  \( 2 \).}
{\( X_1 = \left[2;\frac12\right]\text{, }X_2 = \left[6;\frac52\right] \)}{\( X_1 = \left[2;1\right]\text{, }X_2 = \left[6;5\right] \)}{\( X_1 = \left[2;\frac14\right]\text{, }X_2 = \left[6;\frac54\right] \)}{\( X_1 = \left[2;\frac32\right]\text{, }X_2 = \left[6;\frac72\right] \)}{}{}}
\LANGcs{
\Question{
Je dána přímka \( p \): \( x-2y-1=0 \). Určete souřadnice všech takových bodů ležících na \( p \), které mají od přímky \( x=4 \) vzdálenost \( 2 \).}
{\( X_1 = \left[2;\frac12\right]\text{, }X_2 = \left[6;\frac52\right] \)}{\( X_1 = \left[2;1\right]\text{, }X_2 = \left[6;5\right] \)}{\( X_1 = \left[2;\frac14\right]\text{, }X_2 = \left[6;\frac54\right] \)}{\( X_1 = \left[2;\frac32\right]\text{, }X_2 = \left[6;\frac72\right] \)}{}{}}
\LANGsk{
\Question{
Je daná priamka \( p \): \( x-2y-1=0 \). Určte súradnice všetkých takých bodov ležiacich na \( p \), ktoré majú od priamky \( x=4 \) vzdialenosť \( 2 \).}
{\( X_1 = \left[2;\frac12\right]\text{, }X_2 = \left[6;\frac52\right] \)}{\( X_1 = \left[2;1\right]\text{, }X_2 = \left[6;5\right] \)}{\( X_1 = \left[2;\frac14\right]\text{, }X_2 = \left[6;\frac54\right] \)}{\( X_1 = \left[2;\frac32\right]\text{, }X_2 = \left[6;\frac72\right] \)}{}{}}
\LANGes{
\Question{
Sea \( p \)  la recta con ecuación \( x-2y-1=0 \). Halla las coordenadas de todos los puntos que se encuentran en la recta \( p \) y cuya distancia a la recta \( x=4 \) equivale a \( 2 \).}
{\( X_1 = \left[2;\frac12\right]\text{, }X_2 = \left[6;\frac52\right] \)}{\( X_1 = \left[2;1\right]\text{, }X_2 = \left[6;5\right] \)}{\( X_1 = \left[2;\frac14\right]\text{, }X_2 = \left[6;\frac54\right] \)}{\( X_1 = \left[2;\frac32\right]\text{, }X_2 = \left[6;\frac72\right] \)}{}{}}
\End

\Data\ID{1103109008}
\Updated{2021-12-05 10:12:12}
\Level{B}
\SubArea{Analytic geometry in a plane}
\IMG{1103109008.pdf}
\LANGen{
\Question{
Let \( p \) be the line with the equation \( x-2y-1=0 \). Find the coordinates of all points lying on the line \( p \) such that their distance from the line \( y=3 \) equals to \( 1 \).}
{\( X_1 = \left[5;2\right]\text{, }X_2 = \left[9;4\right] \)}{\( X_1 = \left[4;2\right]\text{, }X_2 = \left[8;4\right] \)}{\( X_1 = \left[2;4\right]\text{, }X_2 = \left[6;4\right] \)}{\( X_1 = \left[2;5\right]\text{, }X_2 = \left[4;9\right] \)}{}{}}
\LANGpl{
\Question{
Dana jest prosta \( p \) określona równaniem \( x-2y-1=0 \). Wyznacz współrzędne wszystkich punktów leżących na prostej \( p \) tak, aby ich odległość od prostej \( y=3 \) była równa \( 1 \).}
{\( X_1 = \left[5;2\right]\text{, }X_2 = \left[9;4\right] \)}{\( X_1 = \left[4;2\right]\text{, }X_2 = \left[8;4\right] \)}{\( X_1 = \left[2;4\right]\text{, }X_2 = \left[6;4\right] \)}{\( X_1 = \left[2;5\right]\text{, }X_2 = \left[4;9\right] \)}{}{}}
\LANGcs{
\Question{
Je dána přímka \( p \): \( x-2y-1=0 \). Určete souřadnice všech bodů ležících na \( p \), které mají od přímky \( y=3 \) vzdálenost \( 1 \).}
{\( X_1 = \left[5;2\right]\text{, }X_2 = \left[9;4\right] \)}{\( X_1 = \left[4;2\right]\text{, }X_2 = \left[8;4\right] \)}{\( X_1 = \left[2;4\right]\text{, }X_2 = \left[6;4\right] \)}{\( X_1 = \left[2;5\right]\text{, }X_2 = \left[4;9\right] \)}{}{}}
\LANGsk{
\Question{
Je daná priamka \( p \): \( x-2y-1=0 \). Určte súradnice všetkých bodov ležiacich na \( p \), ktoré majú od priamky \( y=3 \) vzdialenosť \( 1 \).}
{\( X_1 = \left[5;2\right]\text{, }X_2 = \left[9;4\right] \)}{\( X_1 = \left[4;2\right]\text{, }X_2 = \left[8;4\right] \)}{\( X_1 = \left[2;4\right]\text{, }X_2 = \left[6;4\right] \)}{\( X_1 = \left[2;5\right]\text{, }X_2 = \left[4;9\right] \)}{}{}}
\LANGes{
\Question{
Sea \( p \)  la recta con ecuación \( x-2y-1=0 \). Halla las coordenadas de todos los puntos que se encuentran en la recta \( p \) y cuya distancia a la recta \( y=3 \) equivale a \( 1 \).}
{\( X_1 = \left[5;2\right]\text{, }X_2 = \left[9;4\right] \)}{\( X_1 = \left[4;2\right]\text{, }X_2 = \left[8;4\right] \)}{\( X_1 = \left[2;4\right]\text{, }X_2 = \left[6;4\right] \)}{\( X_1 = \left[2;5\right]\text{, }X_2 = \left[4;9\right] \)}{}{}}
\End

\Data\ID{2010014203}
\Updated{2022-07-20 07:07:05}
\Level{B}
\SubArea{Analytic geometry in a plane}
\LANGen{
\Question{
Find the distance between parallel lines \(p\) and \(q\), if they are given by their general form equations, where \(p\) is  \(-2x-4y+8=0\) and \(q\) is  \(-x-2y+3=0\).}
{\(\frac{\sqrt{5}}5\)}{\(-\frac1{\sqrt{5}}\)}{\(\frac{2\sqrt{5}}5\)}{\(\frac13\)}{}{}}
\LANGpl{
\Question{
Znajdź odległość między równoległymi prostymi \(p\) i \(q\), jeśli są podane przez ich równania w postaci ogólnej, gdzie \(p\) to \(-2x-4y+8=0\) i \( q\) to \(-x-2y+3=0\).}
{\(\frac{\sqrt{5}}5\)}{\(-\frac1{\sqrt{5}}\)}{\(\frac{2\sqrt{5}}5\)}{\(\frac13\)}{}{}}
\LANGcs{
\Question{
Vypočtěte vzdálenost rovnoběžek \(p\) a \(q\), jsou-li zadány jejich obecné rovnice: \(p:\)  \(-2x-4y+8=0\) a \(q:\)  \(-x-2y+3=0\).}
{\(\frac{\sqrt{5}}5\)}{\(-\frac1{\sqrt{5}}\)}{\(\frac{2\sqrt{5}}5\)}{\(\frac13\)}{}{}}
\LANGsk{
\Question{
Určte vzdialenosť rovnobežných priamok \(p\) a \(q\), ak sú dané všeobecnými rovnicami kde \(p\) je \(-2x-4y+8=0\) a \(q\) je \(-x-2y+3=0\).}
{\(\frac{\sqrt{5}}5\)}{\(-\frac1{\sqrt{5}}\)}{\(\frac{2\sqrt{5}}5\)}{\(\frac13\)}{}{}}
\LANGes{
\Question{
Halla la distancia entre las rectas paralelas \(p\) y \(q\), suponiendo que la ecuación general de la recta \(p\) es  \(-2x-4y+8=0\) y la de \(q\) es  \(-x-2y+3=0\).}
{\(\frac{\sqrt{5}}5\)}{\(-\frac1{\sqrt{5}}\)}{\(\frac{2\sqrt{5}}5\)}{\(\frac13\)}{}{}}
\End

\Data\ID{2010014204}
\Updated{2022-05-16 11:05:13}
\Level{B}
\SubArea{Analytic geometry in a plane}
\LANGen{
\Question{
Find the distance between parallel lines \( p \) and \( q \) given by their parametric equations.
\begin{align*}
 p\colon x&=3-2t, & q\colon x&=2+2s, \\
y&=-1+t;\ t\in\mathbb{R}; & y&=1-s;\ s\in\mathbb{R}.
\end{align*}}
{\(\frac{3\sqrt{5}}5\)}{\(-\frac{3\sqrt{5}}5\)}{\(\sqrt{5}\)}{\(\frac{\sqrt{5}}3\)}{}{}}
\LANGpl{
\Question{
Znajdź odległość między równoległymi prostymi \( p \) i \( q \) określonymi przez ich równania parametryczne.
\begin{align*}
 p\colon x&=3-2t, & q\colon x&=2+2s, \\
y&=-1+t;\ t\in\mathbb{R}; & y&=1-s;\ s\in\mathbb{R}.
\end{align*}}
{\(\frac{3\sqrt{5}}5\)}{\(-\frac{3\sqrt{5}}5\)}{\(\sqrt{5}\)}{\(\frac{\sqrt{5}}3\)}{}{}}
\LANGcs{
\Question{
Vypočtěte vzdálenost rovnoběžek \( p \) a \( q \), jsou-li zadána jejich parametrická vyjádření:
\begin{align*}
 p\colon x&=3-2t, & q\colon x&=2+2s, \\
y&=-1+t;\ t\in\mathbb{R}; & y&=1-s;\ s\in\mathbb{R}.
\end{align*}}
{\(\frac{3\sqrt{5}}5\)}{\(-\frac{3\sqrt{5}}5\)}{\(\sqrt{5}\)}{\(\frac{\sqrt{5}}3\)}{}{}}
\LANGsk{
\Question{
Určte vzdialenosť rovnobežných priamok \( p \) a \( q \) daných parametrickými rovnicami:
\begin{align*}
 p\colon x&=3-2t, & q\colon x&=2+2s, \\
y&=-1+t;\ t\in\mathbb{R}; & y&=1-s;\ s\in\mathbb{R}.
\end{align*}}
{\(\frac{3\sqrt{5}}5\)}{\(-\frac{3\sqrt{5}}5\)}{\(\sqrt{5}\)}{\(\frac{\sqrt{5}}3\)}{}{}}
\LANGes{
\Question{
Halla la distancia entre las rectas paralelas \( p \) y \( q \) que vienen dadas por sus ecuaciones parámetricas.
\begin{align*}
 p\colon x&=3-2t, & q\colon x&=2+2s, \\
y&=-1+t;\ t\in\mathbb{R}; & y&=1-s;\ s\in\mathbb{R}.
\end{align*}}
{\(\frac{3\sqrt{5}}5\)}{\(-\frac{3\sqrt{5}}5\)}{\(\sqrt{5}\)}{\(\frac{\sqrt{5}}3\)}{}{}}
\End

\Data\ID{2010014206}
\Updated{2022-05-16 11:05:40}
\Level{B}
\SubArea{Analytic geometry in a plane}
\IMG{2010014201-figure1.pdf}
\LANGen{
\Question{
Let \( p \) be the line with the equation \( x+2y-1=0 \). Find the general form equations of all lines parallel to \( p \) such that their distance from \( p \) equals to \( \sqrt5 \).}
{\( x+2y-6=0;\ x+2y+4=0 \)}{\( x+2y-1=0;\ x+2y+1=0 \)}{\( 2x-y-6=0;\ 2x-y+4=0 \)}{\( 2x-y-1=0;\ 2x-y+1=0 \)}{}{}}
\LANGpl{
\Question{
Niech \( p \) będzie prostą o równaniu \( x+2y-1=0 \). Znajdź równania wszystkich prostych równoległych do \( p \) tak, że ich odległość od \( p \) jest równa \( \sqrt5 \).}
{\( x+2y-6=0;\ x+2y+4=0 \)}{\( x+2y-1=0;\ x+2y+1=0 \)}{\( 2x-y-6=0;\ 2x-y+4=0 \)}{\( 2x-y-1=0;\ 2x-y+1=0 \)}{}{}}
\LANGcs{
\Question{
Je dána přímka \( p: \) \( x+2y-1=0 \). Určete rovnice všech přímek rovnoběžných s přímkou \( p \), které od ní mají vzdálenost \( \sqrt5 \).}
{\( x+2y-6=0;\ x+2y+4=0 \)}{\( x+2y-1=0;\ x+2y+1=0 \)}{\( 2x-y-6=0;\ 2x-y+4=0 \)}{\( 2x-y-1=0;\ 2x-y+1=0 \)}{}{}}
\LANGsk{
\Question{
Daná je priamka \( p: \) \( x+2y-1=0 \). 
Určte rovnice všetkých priamok rovnobežných s priamkou \( p \), ktoré od nej majú vzdialenosť \( \sqrt5 \).}
{\( x+2y-6=0;\ x+2y+4=0 \)}{\( x+2y-1=0;\ x+2y+1=0 \)}{\( 2x-y-6=0;\ 2x-y+4=0 \)}{\( 2x-y-1=0;\ 2x-y+1=0 \)}{}{}}
\LANGes{
\Question{
Sea \( p \) la recta con ecuación \( x+2y-1=0 \). Determina las ecuaciones generales de todas las rectas paralelas a la recta \( p \) suponiendo que la distancia a \( p \) equivale a \( \sqrt5 \).}
{\( x+2y-6=0;\ x+2y+4=0 \)}{\( x+2y-1=0;\ x+2y+1=0 \)}{\( 2x-y-6=0;\ 2x-y+4=0 \)}{\( 2x-y-1=0;\ 2x-y+1=0 \)}{}{}}
\End

\Data\ID{2010014605}
\Updated{2022-05-16 11:05:49}
\Level{B}
\SubArea{Analytic geometry in a plane}
\LANGen{
\Question{
Find the distance from the point \(P = [2;4]\) 
to the line \(4x - 3y - 5 = 0\).}
{\(\frac95\)}{\(3\)}{\(\frac45\)}{The line contains \(P\).}{}{}}
\LANGpl{
\Question{
Wyznacz odległość punktu \(P = [2;4]\)
od prostej \(4x - 3y - 5 = 0\).}
{\(\frac95\)}{\(3\)}{\(\frac45\)}{Punkt \(P\) należy do prostej.}{}{}}
\LANGcs{
\Question{
Určete vzdálenost bodu \(P = [2;4]\) 
od přímky \(4x - 3y - 5 = 0\).}
{\(\frac95\)}{\(3\)}{\(\frac45\)}{\(P\) leží na dané přímce.}{}{}}
\LANGsk{
\Question{
Určte vzdialenosť bodu \(P = [2;4]\) 
od priamky \(4x - 3y - 5 = 0\).}
{\(\frac95\)}{\(3\)}{\(\frac45\)}{Bod \(P\) leží na priamke.}{}{}}
\LANGes{
\Question{
Halla la distancia del punto \(P = [2;4]\) 
a la recta \(4x - 3y - 5 = 0\).}
{\(\frac95\)}{\(3\)}{\(\frac45\)}{El punto está en la recta \(P\).}{}{}}
\End

\Data\ID{2010014606}
\Updated{2022-05-16 11:05:49}
\Level{B}
\SubArea{Analytic geometry in a plane}
\LANGen{
\Question{
Find all the values of the parameter \(c\) so that the distance from the point \(M = [1;-2]\) 
to the line \(-4x + 3y + c = 0\) equals \(5\).}
{\(c\in \{ - 15;35\}\)}{\(c\in \{ 15\}\)}{\(c\in \{ 15;25\}\)}{\(c\in \{ -5;5\}\)}{}{}}
\LANGpl{
\Question{
Znajdź wszystkie wartości parametru \(c\) tak, aby odległość punktu \(M = [1;-2]\)
od prostej \(-4x + 3y + c = 0\) była równa \(5\).}
{\(c\in \{ - 15;35\}\)}{\(c\in \{ 15\}\)}{\(c\in \{ 15;25\}\)}{\(c\in \{ -5;5\}\)}{}{}}
\LANGcs{
\Question{
Určete všechny hodnoty parametru \(c\) tak, aby měl bod \(M = [1;-2]\) 
od přímky \(-4x + 3y + c = 0\) vzdálenost \(5\).}
{\(c\in \{ - 15;35\}\)}{\(c\in \{ 15\}\)}{\(c\in \{ 15;25\}\)}{\(c\in \{ -5;5\}\)}{}{}}
\LANGsk{
\Question{
Určte hodnotu parametra \(c\) tak, aby bola vzdialenosť bodu \(M = [1;-2]\) 
od priamky \(-4x + 3y + c = 0\) rovná \(5\).}
{\(c\in \{ - 15;35\}\)}{\(c\in \{ 15\}\)}{\(c\in \{ 15;25\}\)}{\(c\in \{ -5;5\}\)}{}{}}
\LANGes{
\Question{
Halla el valor (valores) del parámetro \(c\) suponiendo que la distancia del punto \(M = [1;-2]\) 
a la recta \(-4x + 3y + c = 0\) es \(5\).}
{\(c\in \{ - 15;35\}\)}{\(c\in \{ 15\}\)}{\(c\in \{ 15;25\}\)}{\(c\in \{ -5;5\}\)}{}{}}
\End

\Data\ID{2010014607}
\Updated{2023-07-04 09:07:22}
\Level{B}
\SubArea{Analytic geometry in a plane}
\LANGen{
\Question{
Given points \(A = [3;3]\),
\(B = [-5;3]\) and 
\(C = [-1;-1]\), find the length of the altitude of the triangle \(ABC\) through the point \(C\). Hint: The altitude through the point \(C\) of a triangle \(ABC\) is the perpendicular line segment drawn from the vertex \(C\) to the line containing the side \(AB\).}
{\(4\)}{\(\frac43\)}{\(6\)}{\(\frac23\)}{}{}}
\LANGpl{
\Question{
Podane są  punkty \(A = [3;3]\),
\(B = [-5;3]\) i
\(C = [-1;-1]\), znajdź długość wysokości trójkąta \(ABC\) przechodzącej przez punkt \(C\). Podpowiedź: Wysokość przechodząca przez punkt \(C\) trójkąta \(ABC\) to prostopadły odcinek linii narysowany od wierzchołka \(C\) do linii zawierającej bok \(AB\).}
{\(4\)}{\(\frac43\)}{\(6\)}{\(\frac23\)}{}{}}
\LANGcs{
\Question{
V trojúhelníku \(ABC\), kde \(A = [3;3]\), \(B = [-5;3]\) a \(C = [-1;-1]\), určete velikost výšky procházející bodem \(C\). Nápověda: Výška procházející bodem \(C\) v trojúhelníku \(ABC\) je úsečka procházející vrcholem \(C\), která je kolmá k přímce obsahující stranu \(AB\).}
{\(4\)}{\(\frac43\)}{\(6\)}{\(\frac23\)}{}{}}
\LANGsk{
\Question{
Danú sú body \(A = [3;3]\), \(B = [-5;3]\) a 
\(C = [-1;-1]\), určte veľkosť výšky v trojuholníku \(ABC\), ktorá prechádza bodom \(C\). Nápoveda: Výška cez vrchol \(C\) v trojuholníku \(ABC\) je priamka prechádzajúca vrcholom \(C\), ktorá je kolmá na priamku obsahujúcu stranu \(AB\).}
{\(4\)}{\(\frac43\)}{\(6\)}{\(\frac23\)}{}{}}
\LANGes{
\Question{
Dados los puntos \(A = [3;3]\),
\(B = [-5;3]\) y 
\(C = [-1;-1]\), halla la longitud de la altura al punto \(C\) del triángulo \(ABC\). Pista: En geometría, la altura al punto \(C\) del triángulo \(ABC\) es un segmento que une el vértice \(C\) con un punto de su lado opuesto y es perpendicular este lado \(AB\) del triángulo.}
{\(4\)}{\(\frac43\)}{\(6\)}{\(\frac23\)}{}{}}
\End

\Data\ID{2010014608}
\Updated{2022-05-16 11:05:49}
\Level{B}
\SubArea{Analytic geometry in a plane}
\IMG{2010014601-figure0.pdf}
\LANGen{
\Question{
Find a general form equation of the straight line that passes through the point \( M=[2;-3] \) and is parallel with the line of symmetry of the line segment \( AB \), where \( A=[4;-1] \), and \( B=\left[-3;\frac32\right] \) (see the picture).}
{\( 14x-5y-43=0 \)}{\( 5x-14y-52=0 \)}{\( 14x+5y-13=0 \)}{\( 5x+14+32=0 \)}{}{}}
\LANGpl{
\Question{
Znajdź ogólne równanie prostej przechodzącej przez punkt \( M=[2;-3] \) i równoległej do prostej \( AB \), gdzie \( A=[ 4;-1] \) oraz \( B=\left[-3;\frac32\right] \) (patrz rysunek).}
{\( 14x-5y-43=0 \)}{\( 5x-14y-52=0 \)}{\( 14x+5y-13=0 \)}{\( 5x+14+32=0 \)}{}{}}
\LANGcs{
\Question{
Určete obecnou rovnici přímky, která prochází bodem \( M=[2;-3] \) a je rovnoběžná s osou úsečky \( AB \), přičemž \( A=[4;-1] \) a \( B=\left[-3;\frac32\right] \) (viz obrázek).}
{\( 14x-5y-43=0 \)}{\( 5x-14y-52=0 \)}{\( 14x+5y-13=0 \)}{\( 5x+14+32=0 \)}{}{}}
\LANGsk{
\Question{
Určte všeobecnú rovnicu priamky, ktorá prechádza bodom \( M=[2;-3] \) a je rovnobežná s osou úsečky \( AB \), ak platí \( A=[4;-1] \) a \( B=\left[-3;\frac32\right] \) (pozri obrázok).}
{\( 14x-5y-43=0 \)}{\( 5x-14y-52=0 \)}{\( 14x+5y-13=0 \)}{\( 5x+14+32=0 \)}{}{}}
\LANGes{
\Question{
Determina la ecuación general de la recta que pasa por el punto \( M=[2;-3] \) y es paralela al eje de simetría del segmento \( AB \), donde \( A=[4;-1] \), y \( B=\left[-3;\frac32\right] \) (mira la imagen).}
{\( 14x-5y-43=0 \)}{\( 5x-14y-52=0 \)}{\( 14x+5y-13=0 \)}{\( 5x+14+32=0 \)}{}{}}
\End

\Data\ID{2010014609}
\Updated{2022-05-16 11:05:19}
\Level{B}
\SubArea{Analytic geometry in a plane}
\LANGen{
\Question{
Find the angle \(\varphi \) 
between the lines \(2x +1 = 0\) 
and \(x - y + 7 = 0\).}
{\(45^{\circ }\)}{\(60^{\circ }\)}{\(90^{\circ }\)}{\(30^{\circ }\)}{}{}}
\LANGpl{
\Question{
Wyznacz kąt \(\varphi \)
między prostymi \(2x +1 = 0\)
oraz \(x - y + 7 = 0\).}
{\(45^{\circ }\)}{\(60^{\circ }\)}{\(90^{\circ }\)}{\(30^{\circ }\)}{}{}}
\LANGcs{
\Question{
Určete odchylku \(\varphi \) 
přímek zadaných obecnými rovnicemi \(2x +1 = 0\) 
a \(x - y + 7 = 0\).}
{\(45^{\circ }\)}{\(60^{\circ }\)}{\(90^{\circ }\)}{\(30^{\circ }\)}{}{}}
\LANGsk{
\Question{
Určte uhol \(\varphi \) 
priamok \(2x +1 = 0\) 
a \(x - y + 7 = 0\).}
{\(45^{\circ }\)}{\(60^{\circ }\)}{\(90^{\circ }\)}{\(30^{\circ }\)}{}{}}
\LANGes{
\Question{
Determina el ángulo \(\varphi \) 
entre las rectas \(2x +1 = 0\) 
y \(x - y + 7 = 0\).}
{\(45^{\circ }\)}{\(60^{\circ }\)}{\(90^{\circ }\)}{\(30^{\circ }\)}{}{}}
\End

\Data\ID{2010014610}
\Updated{2022-05-16 11:05:19}
\Level{B}
\SubArea{Analytic geometry in a plane}
\LANGen{
\Question{
Given points \(A = [4;-1]\),
\(B = [2,-3]\) and 
\(C = [5,5]\), find the angle
\(\beta \) (the interior angle
at the vertex \(B\))
in the triangle \(ABC\).}
{\(24^{\circ }27'\)}{\(144^{\circ }46'\)}{\(155^{\circ }33'\)}{\(11^{\circ }05'\)}{}{}}
\LANGpl{
\Question{
Podane są punkty \(A = [4;-1]\),
\(B = [2,-3]\) i
\(C = [5,5]\). Wyznacz kąt
\(\beta \) (kąt wewnętrzny
przy wierzchołku \(B\))
w trójkącie \(ABC\).}
{\(24^{\circ }27'\)}{\(144^{\circ }46'\)}{\(155^{\circ }33'\)}{\(11^{\circ }05'\)}{}{}}
\LANGcs{
\Question{
Je dán trojúhelník \(ABC\) , kde \(A = [4;-1]\), \(B = [2,-3]\) a \(C = [5,5]\). Vypočítejte velikost vnitřního úhlu 
\(\beta \) u vrcholu \(B\) trojúhelníku \(ABC\).}
{\(24^{\circ }27'\)}{\(144^{\circ }46'\)}{\(155^{\circ }33'\)}{\(11^{\circ }05'\)}{}{}}
\LANGsk{
\Question{
Dané sú body \(A = [4;-1]\), \(B = [2,-3]\) a \(C = [5,5]\) v trojuholníku \(ABC\). Určte uhol \(\beta \) (vnútorný uhol pro vrchole \(B\))
v tomto trojuholníku.}
{\(24^{\circ }27'\)}{\(144^{\circ }46'\)}{\(155^{\circ }33'\)}{\(11^{\circ }05'\)}{}{}}
\LANGes{
\Question{
Dados los puntos \(A = [4;-1]\),
\(B = [2,-3]\) y 
\(C = [5,5]\), halla el ángulo
\(\beta \) (el ángulo interior del vértice \(B\))
en el triángulo \(ABC\).}
{\(24^{\circ }27'\)}{\(144^{\circ }46'\)}{\(155^{\circ }33'\)}{\(11^{\circ }05'\)}{}{}}
\End

\Data\ID{9000107504}
\Updated{2020-07-17 09:07:43}
\Level{B}
\SubArea{Analytic geometry in a plane}
\LANGen{
\Question{
Find the angle between the lines 
\[ 
 p\colon 2x - 3y + 1 = 0;\quad q\colon 3x + 2y - 3 = 0.
\]}
{\(90^{\circ }\)}{\(60^{\circ }\)}{\(0^{\circ }\)}{\(30^{\circ }\)}{}{}}
\LANGpl{
\Question{
Wyznacz kąt pomiędzy prostymi
\[ 
 p\colon 2x - 3y + 1 = 0;\quad q\colon 3x + 2y - 3 = 0.
\]}
{\(90^{\circ }\)}{\(60^{\circ }\)}{\(0^{\circ }\)}{\(30^{\circ }\)}{}{}}
\LANGcs{
\Question{
Odchylka přímek \(p\colon 2x - 3y + 1 = 0;\ q\colon 3x + 2y - 3 = 0\) 
je rovna:}
{\(90^{\circ }\)}{\(60^{\circ }\)}{\(0^{\circ }\)}{\(30^{\circ }\)}{}{}}
\LANGsk{
\Question{
Odchýlka priamok \(p\colon 2x - 3y + 1 = 0;\ q\colon 3x + 2y - 3 = 0\) 
je rovná:}
{\(90^{\circ }\)}{\(60^{\circ }\)}{\(0^{\circ }\)}{\(30^{\circ }\)}{}{}}
\LANGes{
\Question{
Halla el ángulo  que forman las siguientes rectas
\[ 
 p\colon 2x - 3y + 1 = 0;\quad q\colon 3x + 2y - 3 = 0.
\]}
{\(90^{\circ }\)}{\(60^{\circ }\)}{\(0^{\circ }\)}{\(30^{\circ }\)}{}{}}
\End

\Data\ID{9000107505}
\Updated{2020-07-17 09:07:09}
\Level{B}
\SubArea{Analytic geometry in a plane}
\LANGen{
\Question{
Find \(\cos \varphi \) where
\(\varphi \) is the angle
between the lines \(p\) 
and \(q\).
\[ 
\begin{aligned}[t] p\colon x& = 1 + 4t, &
\\y & = 3 - 3t;\ t\in \mathbb{R};
\\ \end{aligned} \quad q\colon x + y - 3 = 0
\]}
{\(\frac{7\sqrt{2}}
 {10} \)}{\(- \frac{7} 
{5\sqrt{2}}\)}{\(\frac{\sqrt{2}}
 {5} \)}{\(\frac{\sqrt{2}}
{10} \)}{}{}}
\LANGpl{
\Question{
Określ  \(\cos \varphi \) gdzie
\(\varphi \) jest kątem pomiędzy 
prostymi \(p\) 
i \(q\).
\[ 
\begin{aligned}[t] p\colon x& = 1 + 4t, &
\\y & = 3 - 3t;\ t\in \mathbb{R};
\\ \end{aligned} \quad q\colon x + y - 3 = 0
\]}
{\(\frac{7\sqrt{2}}
 {10} \)}{\(- \frac{7} 
{5\sqrt{2}}\)}{\(\frac{\sqrt{2}}
 {5} \)}{\(\frac{\sqrt{2}}
{10} \)}{}{}}
\LANGcs{
\Question{
Kosinus odchylek přímek \(p\colon x = 1 + 4t;\ y = 3 - 3t;\ t\in \mathbb{R}\) 
a \(q\colon x + y - 3 = 0\) se
rovná:}
{\(\frac{7\sqrt{2}}
 {10} \)}{\(- \frac{7} 
{5\sqrt{2}}\)}{\(\frac{\sqrt{2}}
 {5} \)}{\(\frac{\sqrt{2}}
{10} \)}{}{}}
\LANGsk{
\Question{
Kosínus odchýlok priamok \(p\colon x = 1 + 4t;\ y = 3 - 3t;\ t\in \mathbb{R}\) 
a \(q\colon x + y - 3 = 0\) sa
rovná:}
{\(\frac{7\sqrt{2}}
 {10} \)}{\(- \frac{7} 
{5\sqrt{2}}\)}{\(\frac{\sqrt{2}}
 {5} \)}{\(\frac{\sqrt{2}}
{10} \)}{}{}}
\LANGes{
\Question{
Halla \(\cos \varphi \) donde
\(\varphi \) es el ángulo que forman las rectas \(p\) 
y \(q\).
\[ 
\begin{aligned}[t] p\colon x& = 1 + 4t, &
\\y & = 3 - 3t;\ t\in \mathbb{R};
\\ \end{aligned} \quad q\colon x + y - 3 = 0
\]}
{\(\frac{7\sqrt{2}}
 {10} \)}{\(- \frac{7} 
{5\sqrt{2}}\)}{\(\frac{\sqrt{2}}
 {5} \)}{\(\frac{\sqrt{2}}
{10} \)}{}{}}
\End

\Data\ID{9000107506}
\Updated{2020-07-17 09:07:05}
\Level{B}
\SubArea{Analytic geometry in a plane}
\LANGen{
\Question{
Find \(\cos \varphi \) where
\(\varphi \) is the angle
between the lines \(p\) 
and \(q\).
\[ 
p\colon y = 2x - 11;\quad q\colon y = \frac{1} 
{4}x
\]}
{\(\frac{6\sqrt{85}}
 {85} \)}{\(\frac{1}
{\sqrt{22}}\)}{\(\frac{\sqrt{6}}
{85} \)}{\(\frac{\sqrt{17}}
 {30} \)}{}{}}
\LANGpl{
\Question{
Określ  \(\cos \varphi \), gdzie
\(\varphi \) jest kątem 
pomiędzy prostymi \(p\) 
i \(q\).
\[ 
p\colon y = 2x - 11;\quad q\colon y = \frac{1} 
{4}x
\]}
{\(\frac{6\sqrt{85}}
 {85} \)}{\(\frac{1}
{\sqrt{22}}\)}{\(\frac{\sqrt{6}}
{85} \)}{\(\frac{\sqrt{17}}
 {30} \)}{}{}}
\LANGcs{
\Question{
Kosinus odchylky přímek \(p\colon y = 2x - 11\) 
a \(q\colon y = \frac{1} 
{4}x\) se
rovná:}
{\(\frac{6\sqrt{85}}
 {85} \)}{\(\frac{1}
{\sqrt{22}}\)}{\(\frac{\sqrt{6}}
{85} \)}{\(\frac{\sqrt{17}}
 {30} \)}{}{}}
\LANGsk{
\Question{
Kosínus odchýlky priamok \(p\colon y = 2x - 11\) 
a \(q\colon y = \frac{1} 
{4}x\) sa
rovná:}
{\(\frac{6\sqrt{85}}
 {85} \)}{\(\frac{1}
{\sqrt{22}}\)}{\(\frac{\sqrt{6}}
{85} \)}{\(\frac{\sqrt{17}}
 {30} \)}{}{}}
\LANGes{
\Question{
Halla \(\cos \varphi \) donde
\(\varphi \)  es el ángulo que forman las rectas \(p\) 
y \(q\).
\[ 
p\colon y = 2x - 11;\quad q\colon y = \frac{1} 
{4}x
\]}
{\(\frac{6\sqrt{85}}
 {85} \)}{\(\frac{1}
{\sqrt{22}}\)}{\(\frac{\sqrt{6}}
{85} \)}{\(\frac{\sqrt{17}}
 {30} \)}{}{}}
\End

\Data\ID{9000107507}
\Updated{2020-07-17 09:07:54}
\Level{B}
\SubArea{Analytic geometry in a plane}
\LANGen{
\Question{
Find \(\mathop{\mathrm{tg}}\nolimits \varphi \) where
\(\varphi \) is the angle
between the lines \(p\) 
and \(q\).
\[ 
\begin{aligned}[t] p\colon x& = 1 + t, &
\\y & = 3 + 2t;\ t\in \mathbb{R};
\\ \end{aligned}\quad q\colon y = 1
\]}
{\(2\)}{\(\frac{1}
{2}\)}{\(- 1\)}{\(0\)}{}{}}
\LANGpl{
\Question{
Oblicz  \(\mathop{\mathrm{tg}}\nolimits \varphi \) gdzie 
\(\varphi \) jest kątem
miedzy prostą   \(p\) 
i \(q\).
\[ 
\begin{aligned}[t] p\colon x& = 1 + t, &
\\y & = 3 + 2t;\ t\in \mathbb{R};
\\ \end{aligned}\quad q\colon y = 1
\]}
{\(2\)}{\(\frac{1}
{2}\)}{\(- 1\)}{\(0\)}{}{}}
\LANGcs{
\Question{
Je dána přímka \(p\colon x = 1 + t;\ y = 3 + 2t;\ t\in \mathbb{R}\) 
a přímka \(q\colon y = 1\). Tangens
odchylky přímek \(p,\ q\) 
je roven:}
{\(2\)}{\(\frac{1}
{2}\)}{\(- 1\)}{\(0\)}{}{}}
\LANGsk{
\Question{
Je daná priamka \(p\colon x = 1 + t;\ y = 3 + 2t;\ t\in \mathbb{R}\) 
a priamka \(q\colon y = 1\). Tangens
odchýlky priamok \(p,\ q\) 
je rovný:}
{\(2\)}{\(\frac{1}
{2}\)}{\(- 1\)}{\(0\)}{}{}}
\LANGes{
\Question{
Halla \(\mathop{\mathrm{tg}}\nolimits \varphi \) donde
\(\varphi \) es el ángulo que forman las rectas \(p\) 
y \(q\).
\[ 
\begin{aligned}[t] p\colon x& = 1 + t, &
\\y & = 3 + 2t;\ t\in \mathbb{R};
\\ \end{aligned}\quad q\colon y = 1
\]}
{\(2\)}{\(\frac{1}
{2}\)}{\(- 1\)}{\(0\)}{}{}}
\End

\Data\ID{9000107508}
\Updated{2020-07-17 09:07:19}
\Level{B}
\SubArea{Analytic geometry in a plane}
\LANGen{
\Question{
Find \(\cos \varphi \) where
\(\varphi \) is the angle
between the lines \(p\) 
and \(q\).

\[ 
\begin{aligned}[t] p\colon x& = t, &
\\y & = -3;\ t\in \mathbb{R};
\\ \end{aligned}\quad q\colon y = 1
\]}
{\(1\)}{\(\frac{1}
{\sqrt{2}}\)}{\(0\)}{\(\frac{\sqrt{10}}
 {10} \)}{}{}}
\LANGpl{
\Question{
Oblicz \(\cos \varphi \), gdzie
\(\varphi \) jest kątem
między prostymi \(p\) 
i \(q\).
\[ 
\begin{aligned}[t] p\colon x& = t, &
\\y & = -3;\ t\in \mathbb{R};
\\ \end{aligned}\quad q\colon y = 1
\]}
{\(1\)}{\(\frac{1}
{\sqrt{2}}\)}{\(0\)}{\(\frac{\sqrt{10}}
 {10} \)}{}{}}
\LANGcs{
\Question{
Kosinus odchylky přímek \(p\colon x = t;\ y = -3;\ t\in \mathbb{R}\) 
a \(q\colon y = 1\) je
roven:}
{\(1\)}{\(\frac{1}
{\sqrt{2}}\)}{\(0\)}{\(\frac{\sqrt{10}}
 {10} \)}{}{}}
\LANGsk{
\Question{
Kosínus odchýlky priamok \(p\colon x = t;\ y = -3;\ t\in \mathbb{R}\) 
a \(q\colon y = 1\) je
rovný:}
{\(1\)}{\(\frac{1}
{\sqrt{2}}\)}{\(0\)}{\(\frac{\sqrt{10}}
 {10} \)}{}{}}
\LANGes{
\Question{
Halla \(\cos \varphi \) donde
\(\varphi \) es el ángulo que forman las rectas\(p\) 
y \(q\).

\[ 
\begin{aligned}[t] p\colon x& = t, &
\\y & = -3;\ t\in \mathbb{R};
\\ \end{aligned}\quad q\colon y = 1
\]}
{\(1\)}{\(\frac{1}
{\sqrt{2}}\)}{\(0\)}{\(\frac{\sqrt{10}}
 {10} \)}{}{}}
\End

\Data\ID{9000107509}
\Updated{2021-12-05 10:12:32}
\Level{B}
\SubArea{Analytic geometry in a plane}
\LANGen{
\Question{
In the following list identify a parametric line such that the angle between this line and
the line \(q\) 
is \(0^{\circ }\).
\[ 
 q\colon x - 2y + 11 = 0
\]}
{\(\begin{aligned}[t] p\colon x& = 1 + 4t, &
\\y & = 3 + 2t;\ t\in \mathbb{R}
\\ \end{aligned}\)}{\(\begin{aligned}[t] p\colon x& = 1 + 2t, &
\\y & = 2 - t;\ t\in \mathbb{R}
\\ \end{aligned}\)}{\(\begin{aligned}[t] p\colon x& = 2 - t, &
\\y & = 3 + 2t;\ t\in \mathbb{R}
\\ \end{aligned}\)}{\(\begin{aligned}[t] p\colon x& = t, &
\\y & = 1 - 2t;\ t\in \mathbb{R}
\\ \end{aligned}\)}{}{}}
\LANGpl{
\Question{
Wskaż prostą w postaci parametrycznej tak, aby kąt między tą prostą, a prostą  \(q\) 
wynosił \(0^{\circ }\).
\[ 
 q\colon x - 2y + 11 = 0
\]}
{\(\begin{aligned}[t] p\colon x& = 1 + 4t, &
\\y & = 3 + 2t;\ t\in \mathbb{R}
\\ \end{aligned}\)}{\(\begin{aligned}[t] p\colon x& = 1 + 2t, &
\\y & = 2 - t;\ t\in \mathbb{R}
\\ \end{aligned}\)}{\(\begin{aligned}[t] p\colon x& = 2 - t, &
\\y & = 3 + 2t;\  t\in \mathbb{R}
\\ \end{aligned}\)}{\(\begin{aligned}[t] p\colon x& = t, &
\\y & = 1 - 2t;\ t\in \mathbb{R}
\\ \end{aligned}\)}{}{}}
\LANGcs{
\Question{
Z následujících přímek zadaných parametricky vyberte tu, která má s
přímkou \(q\colon x - 2y + 11 = 0\) 
odchylku \(0^{\circ }\):}
{\(p\colon x = 1 + 4t,\ y = 3 + 2t;\ t\in \mathbb{R}\)}{\(p\colon x = 1 + 2t,\ y = 2 - t;\ t\in \mathbb{R}\)}{\(p\colon x = 2 - t,\ y = 3 + 2t;\ t\in \mathbb{R}\)}{\(p\colon x = t,\ y = 1 - 2t;\ t\in \mathbb{R}\)}{}{}}
\LANGsk{
\Question{
Z následujúcich priamok zadaných parametricky vyberte tú, ktorá má s
priamkou \(q\colon x - 2y + 11 = 0\) 
odchýlku \(0^{\circ }\):}
{\(p\colon x = 1 + 4t,\ y = 3 + 2t;\ t\in \mathbb{R}\)}{\(p\colon x = 1 + 2t,\ y = 2 - t;\ t\in \mathbb{R}\)}{\(p\colon x = 2 - t,\ y = 3 + 2t;\ t\in \mathbb{R}\)}{\(p\colon x = t,\ y = 1 - 2t;\ t\in \mathbb{R}\)}{}{}}
\LANGes{
\Question{
En la siguiente lista, identifica una recta paramétrica de manera el ángulo entre esta recta y
la recta \(q\) 
sea \(0^{\circ }\).
\[ 
 q\colon x - 2y + 11 = 0
\]}
{\(\begin{aligned}[t] p\colon x& = 1 + 4t, &
\\y & = 3 + 2t;\ t\in \mathbb{R}
\\ \end{aligned}\)}{\(\begin{aligned}[t] p\colon x& = 1 + 2t, &
\\y & = 2 - t;\ t\in \mathbb{R}
\\ \end{aligned}\)}{\(\begin{aligned}[t] p\colon x& = 2 - t, &
\\y & = 3 + 2t;\ t\in \mathbb{R}
\\ \end{aligned}\)}{\(\begin{aligned}[t] p\colon x& = t, &
\\y & = 1 - 2t;\ t\in \mathbb{R}
\\ \end{aligned}\)}{}{}}
\End

\Data\ID{9000149401}
\Updated{2022-04-23 05:04:00}
\Level{B}
\SubArea{Analytic geometry in a plane}
\LANGen{
\Question{
Find the distance from the point \(P = [-4;2]\) 
to the line \(3x - 4y - 5 = 0\).}
{\(5\)}{\(1\)}{The line contains \(P\).}{\(\sqrt{5}\)}{}{}}
\LANGpl{
\Question{
Wyznacz odległość punktu \(P = [-4;2]\) 
do prostej \(p\colon 3x - 4y - 5 = 0\).}
{\(5\)}{\(1\)}{Punkt leży na prostej \(p\).}{\(\sqrt{5}\)}{}{}}
\LANGcs{
\Question{
Určete vzdálenost bodu \(P = [-4;2]\) 
od přímky \(p\colon 3x - 4y - 5 = 0\).}
{\(5\)}{\(1\)}{Bod leží na přímce.}{\(\sqrt{5}\)}{}{}}
\LANGsk{
\Question{
Určte vzdialenosť bodu \(P = [-4;2]\) 
od priamky \(p\colon 3x - 4y - 5 = 0\).}
{\(5\)}{\(1\)}{Bod leží na priamke.}{\(\sqrt{5}\)}{}{}}
\LANGes{
\Question{
Halla la distancia del punto \(P = [-4;2]\) 
a la recta \(p\colon 3x - 4y - 5 = 0\).}
{\(5\)}{\(1\)}{El punto está en la recta \(p\).}{\(\sqrt{5}\)}{}{}}
\End

\Data\ID{9000149402}
\Updated{2020-07-17 04:07:37}
\Level{B}
\SubArea{Analytic geometry in a plane}
\LANGen{
\Question{
Find the distance from the origin (i.e. the point
\([0.0]\)) to the
line \(p\colon x + 2y + 5 = 0\).}
{\(\sqrt{5}\)}{\(1\)}{The origin is on the line \(p\).}{\(8\)}{}{}}
\LANGpl{
\Question{
Określ odległość od początku układu współrzędnych (tj. punktu
\([0{,}0]\)) do 
prostej \(p\colon x + 2y + 5 = 0\).}
{\(\sqrt{5}\)}{\(1\)}{Początek układu współrzędnych leży na prostej \(p\).}{\(8\)}{}{}}
\LANGcs{
\Question{
Určete vzdálenost počátku kartézské soustavy souřadnic od přímky
\(p\colon x + 2y + 5 = 0\).}
{\(\sqrt{5}\)}{\(1\)}{Přímka prochází počátkem kartézské soustavy souřadnic.}{\(8\)}{}{}}
\LANGsk{
\Question{
Určte vzdialenosť začiatku karteziánskej sústavy súradníc od priamky
\(p\colon x + 2y + 5 = 0\).}
{\(\sqrt{5}\)}{\(1\)}{Priamka prechádza začiatkom karteziánskej sústavy súradníc.}{\(8\)}{}{}}
\LANGes{
\Question{
Halla la distancia del origen (el punto
\([0.0]\)) a la recta \(p\colon x + 2y + 5 = 0\).}
{\(\sqrt{5}\)}{\(1\)}{El origen está en la recta \(p\).}{\(8\)}{}{}}
\End

\Data\ID{9000149403}
\Updated{2020-07-17 04:07:49}
\Level{B}
\SubArea{Analytic geometry in a plane}
\LANGen{
\Question{
Find the distance from the point \(M = [1;1]\) 
to the line \(p\).
\[ 
\begin{aligned}p\colon x& = 3 + t, &
\\y & = 1 + t;\ t\in \mathbb{R}
\\ \end{aligned}
\]}
{\(\sqrt{2}\)}{\(2\)}{\(1\)}{\(0\) (the point is
on the line \(p\) )}{}{}}
\LANGpl{
\Question{
Wyznacz odległość od punktu \(M = [1;1]\) 
do prostej \(p\).
\[ 
\begin{aligned}p\colon x& = 3 + t, &
\\y & = 1 + t;\  t\in \mathbb{R}
\\ \end{aligned}
\]}
{\(\sqrt{2}\)}{\(2\)}{\(1\)}{\(0\) (punkt
leży na prostej \(p\) )}{}{}}
\LANGcs{
\Question{
Určete vzdálenost bodu \(M = [1;1]\) 
od přímky \(p\colon x = 3 + t\),
\(y = 1 + t\),
\(t\in \mathbb{R}\).}
{\(\sqrt{2}\)}{\(2\)}{\(1\)}{Bod leží na přímce.}{}{}}
\LANGsk{
\Question{
Určte vzdialenosť bodu \(M = [1;1]\) 
od priamky \(p\colon x = 3 + t\),
\(y = 1 + t\),
\(t\in \mathbb{R}\).}
{\(\sqrt{2}\)}{\(2\)}{\(1\)}{Bod leží na priamke.}{}{}}
\LANGes{
\Question{
Halla la distancia del punto \(M = [1;1]\) 
a la recta \(p\).
\[ 
\begin{aligned}p\colon x& = 3 + t, &
\\y & = 1 + t;\ t\in \mathbb{R}
\\ \end{aligned}
\]}
{\(\sqrt{2}\)}{\(2\)}{\(1\)}{\(0\) (el punto está en la recta \(p\) )}{}{}}
\End

\Data\ID{9000149404}
\Updated{2020-07-17 05:07:59}
\Level{B}
\SubArea{Analytic geometry in a plane}
\LANGen{
\Question{
Given points \(A = [-3;13]\),
\(K = [0;4]\),
\(L = [-5;-6]\), find the distance
from the point \(A\) 
to the line \(KL\).}
{\(3\sqrt{5}\)}{\(3\)}{\(5\)}{\(\sqrt{5}\)}{}{}}
\LANGpl{
\Question{
Dane są punkty \(A = [-3;13]\),
\(K = [0;4]\),
\(L = [-5;-6]\), wyznacz odległość
od punktu \(A\) 
do prostej \(KL\).}
{\(3\sqrt{5}\)}{\(3\)}{\(5\)}{\(\sqrt{5}\)}{}{}}
\LANGcs{
\Question{
Určete vzdálenost bodu \(A = [-3;13]\) 
od přímky \(KL\),
kde \(K = [0;4]\),
\(L = [-5;-6]\).}
{\(3\sqrt{5}\)}{\(3\)}{\(5\)}{\(\sqrt{5}\)}{}{}}
\LANGsk{
\Question{
Určte vzdialenosť bodu \(A = [-3;13]\) 
od priamky \(KL\),
kde \(K = [0;4]\),
\(L = [-5;-6]\).}
{\(3\sqrt{5}\)}{\(3\)}{\(5\)}{\(\sqrt{5}\)}{}{}}
\LANGes{
\Question{
Dados los puntos \(A = [-3;13]\),
\(K = [0;4]\),
\(L = [-5;-6]\), halla la distancia del punto \(A\) 
a la recta \(KL\).}
{\(3\sqrt{5}\)}{\(3\)}{\(5\)}{\(\sqrt{5}\)}{}{}}
\End

\Data\ID{9000149405}
\Updated{2022-04-23 05:04:00}
\Level{B}
\SubArea{Analytic geometry in a plane}
\LANGen{
\Question{
Find all the values of the parameter
\(c\) so that the distance from the point \(M = [2;-1]\) 
to the line \(3x + 4y + c = 0\) is \(5\).}
{\(c\in \{ - 27;23\}\)}{\(c\in \{25\}\)}{\(c\in \{5;25\}\)}{\(c\in \{ - 25;25\}\)}{}{}}
\LANGpl{
\Question{
Określ wartość 
\(c\) tak, aby
odległość od punktu \(M = [2;-1]\) 
do prostej \(p\) 
wynosiła \(5\). Prosta
 \(p\) spełnia równanie
\[ 
 p\colon 3x + 4y + c = 0.
\]}
{\(c\in \{ - 27;23\}\)}{\(c\in \{25\}\)}{\(c\in \{5;25\}\)}{\(c\in \{ - 25;25\}\)}{}{}}
\LANGcs{
\Question{
Určete všechny hodnoty parametru \(c\) 
tak, aby bod \(M = [2;-1]\) měl
od přímky \(p\colon 3x + 4y + c = 0\) 
vzdálenost \(5\).}
{\(c\in \{ - 27;23\}\)}{\(c\in \{25\}\)}{\(c\in \{5;25\}\)}{\(c\in \{ - 25;25\}\)}{}{}}
\LANGsk{
\Question{
Určte všetky hodnoty parametra \(c\) 
tak, aby bod \(M = [2;-1]\) mal
od priamky \(p\colon 3x + 4y + c = 0\) 
vzdialenosť \(5\).}
{\(c\in \{ - 27;23\}\)}{\(c\in \{25\}\)}{\(c\in \{5;25\}\)}{\(c\in \{ - 25;25\}\)}{}{}}
\LANGes{
\Question{
Halla el valor (valores) del parámetro
\(c\) suponiendo que la distancia del punto \(M = [2;-1]\) 
a la recta \(p\) 
es \(5\). La recta \(p\) está
definida por la ecuación 
\[ 
 p\colon 3x + 4y + c = 0.
\]}
{\(c\in \{ - 27;23\}\)}{\(c\in \{25\}\)}{\(c\in \{5;25\}\)}{\(c\in \{ - 25;25\}\)}{}{}}
\End

\Data\ID{9000149406}
\Updated{2022-04-23 05:04:00}
\Level{B}
\SubArea{Analytic geometry in a plane}
\LANGen{
\Question{
Given points \(A = [2;-5]\),
\(B = [2;3]\) and 
\(C = [-4;-1]\), find the length of the altitude of the triangle \(ABC\) through the point \(C\). Hint: The altitude through the point 
\(C\) of a triangle
\(ABC\) is the perpendicular line segment drawn from the vertex \(C\) to the line containing the side  \(AB\).}
{\(6\)}{\(\sqrt{2}\)}{\(\frac{3}
{2}\)}{The points \(A\),
\(B\),
\(C\) do
not define a triangle.}{}{}}
\LANGpl{
\Question{
Dane są punkty\(A = [2;-5]\),
\(B = [2;3]\),
\(C = [-4;-1]\), określ długość wysokości trójkąta \(ABC\)
przechodzącej przez punkt \(C\)}
{\(6\)}{\(\sqrt{2}\)}{\(\frac{3}
{2}\)}{Punkty \(A\),
\(B\),
\(C\) nie tworzą trójkąta.}{}{}}
\LANGcs{
\Question{
V trojúhelníku \(ABC\),
kde \(A = [2;-5]\),
\(B = [2;3]\),
\(C = [-4;-1]\), určete velikost
výšky na stranu \(AB\).}
{\(6\)}{\(\sqrt{2}\)}{\(\frac{3}
{2}\)}{Body \(A\),
\(B\),
\(C\) 
netvoří trojúhelník.}{}{}}
\LANGsk{
\Question{
V trojuholníku \(ABC\),
kde \(A = [2;-5]\),
\(B = [2;3]\),
\(C = [-4;-1]\), určte veľkosť
výšky na stranu \(AB\).}
{\(6\)}{\(\sqrt{2}\)}{\(\frac{3}
{2}\)}{Body \(A\),
\(B\),
\(C\) 
netvorí trojuholník.}{}{}}
\LANGes{
\Question{
Dados los puntos \(A = [2;-5]\),
\(B = [2;3]\),
\(C = [-4;-1]\), halla la longitud de la altura al punto \(C\) 
del triángulo \(ABC\). Pista: En geometría, la altura al punto 
\(C\) del triángulo
\(ABC\) es un segmento que une un vértice \(C\) con un punto de su lado opuesto y es perpendicular al lado \(AB\) 
del triángulo.}
{\(6\)}{\(\sqrt{2}\)}{\(\frac{3}
{2}\)}{Los puntos \(A\),
\(B\),
\(C\) 
no definen un triángulo.}{}{}}
\End

\Data\ID{9000149407}
\Updated{2020-07-17 05:07:07}
\Level{B}
\SubArea{Analytic geometry in a plane}
\LANGen{
\Question{
Find the distance between the lines \(p\colon 3x - 4y + 1 = 0\) 
and \(q\colon 3x - 4y + 4 = 0\).}
{\(\frac{3}
{5}\)}{\(1\)}{\(4\)}{\(0\) (the
lines have an intersection)}{}{}}
\LANGpl{
\Question{
Oblicz odległość między prostymi \(p\colon 3x - 4y + 1 = 0\) 
i \(q\colon 3x - 4y + 4 = 0\).}
{\(\frac{3}
{5}\)}{\(1\)}{\(4\)}{\(0\) (proste mają punkt wspólny)}{}{}}
\LANGcs{
\Question{
Určete vzdálenost přímky \(p\colon 3x - 4y + 1 = 0\) 
od přímky \(q\colon 3x - 4y + 4 = 0\).}
{\(\frac{3}
{5}\)}{\(1\)}{\(4\)}{Přímky \(p\) a
\(q\) mají společný bod
a jejich vzdálenost je \(0\).}{}{}}
\LANGsk{
\Question{
Určte vzdialenosť priamky \(p\colon 3x - 4y + 1 = 0\) 
od priamky \(q\colon 3x - 4y + 4 = 0\).}
{\(\frac{3}
{5}\)}{\(1\)}{\(4\)}{Priamky \(p\) a
\(q\) majú spoločný bod
a ich vzdialenosť je \(0\).}{}{}}
\LANGes{
\Question{
Halla la distancia entre las rectas \(p\colon 3x - 4y + 1 = 0\) 
y \(q\colon 3x - 4y + 4 = 0\).}
{\(\frac{3}
{5}\)}{\(1\)}{\(4\)}{\(0\) (las rectas tienen una intersección)}{}{}}
\End

\Data\ID{9000149408}
\Updated{2021-12-05 10:12:56}
\Level{B}
\SubArea{Analytic geometry in a plane}
\LANGen{
\Question{
On the \(x\)-axis
find the points such that the distance from these points to the line
\(p\colon x - 2y + 2 = 0\) is
\(\sqrt{5}\).}
{\([3;0]\),
\([-7;0]\)}{\([5;0]\)}{\(\left [\sqrt{5};0\right ]\),
\(\left [-\sqrt{5};0\right ]\)}{\([3;7]\)}{}{}}
\LANGpl{
\Question{
Wskaż wszystkie punkty na osi \(x\) tak, aby odległość od tych punktów do prostej 
\(p\colon x - 2y + 2 = 0\) była równa
\(\sqrt{5}\).}
{\([3;0]\),
\([-7;0]\)}{\([5;0]\)}{\(\left [\sqrt{5};0\right ]\),
\(\left [-\sqrt{5};0\right ]\)}{\([3;7]\)}{}{}}
\LANGcs{
\Question{
Na ose \(x\) 
najděte všechny body, které mají od přímky
\(p\colon x - 2y + 2 = 0\) vzdálenost
\(\sqrt{5}\).}
{\([3;0]\),
\([-7;0]\)}{\([5;0]\)}{\(\left [\sqrt{5};0\right ]\),
\(\left [-\sqrt{5};0\right ]\)}{\([3;7]\)}{}{}}
\LANGsk{
\Question{
Na osy \(x\) 
nájdite všetky body, ktoré majú od priamky
\(p\colon x - 2y + 2 = 0\) vzdialenosť
\(\sqrt{5}\).}
{\([3;0]\),
\([-7;0]\)}{\([5;0]\)}{\(\left [\sqrt{5};0\right ]\),
\(\left [-\sqrt{5};0\right ]\)}{\([3;7]\)}{}{}}
\LANGes{
\Question{
Halla los puntos del eje \(x\) tales que la distancia de estos puntos a la recta
\(p\colon x - 2y + 2 = 0\) sea
\(\sqrt{5}\).}
{\([3;0]\),
\([-7;0]\)}{\([5;0]\)}{\(\left [\sqrt{5};0\right ]\),
\(\left [-\sqrt{5};0\right ]\)}{\([3;7]\)}{}{}}
\End

\Data\ID{9000149409}
\Updated{2021-12-05 07:12:27}
\Level{B}
\SubArea{Analytic geometry in a plane}
\LANGen{
\Question{
Find all lines which are parallel to \(p\colon x - 3y + 2 = 0\) 
and the distance from every of these lines to
\(p\) is
\(\sqrt{10}\).}
{\(p_{1}\colon x - 3y + 12 = 0\),
\(p_{2}\colon x - 3y - 8 = 0\)}{\(p\colon x - 3y = 0\)}{\(p\colon x - 3y + \sqrt{10} = 0\)}{\(p_{1}\colon x - 3y + \sqrt{10} = 0\),
\(p_{2}\colon x - 3y -\sqrt{10} = 0\)}{}{}}
\LANGpl{
\Question{
Wyznacz wszystkie proste, które są równoległe do  \(p\colon x - 3y + 2 = 0\)
tak, aby odległość tych prostych do 
\(p\) była równa
\(\sqrt{10}\).}
{\(p_{1}\colon x - 3y + 12 = 0\),
\(p_{2}\colon x - 3y - 8 = 0\)}{\(p\colon x - 3y = 0\)}{\(p\colon x - 3y + \sqrt{10} = 0\)}{\(p_{1}\colon x - 3y + \sqrt{10} = 0\),
\(p_{2}\colon x - 3y -\sqrt{10} = 0\)}{}{}}
\LANGcs{
\Question{
Najděte všechny přímky, které jsou rovnoběžné s přímkou
\(p\colon x - 3y + 2 = 0\) a mají od ní
vzdálenost \(\sqrt{10}\).}
{\(p_{1}\colon x - 3y + 12 = 0\),
\(p_{2}\colon x - 3y - 8 = 0\)}{\(p\colon x - 3y = 0\)}{\(p\colon x - 3y + \sqrt{10} = 0\)}{\(p_{1}\colon x - 3y + \sqrt{10} = 0\),
\(p_{2}\colon x - 3y -\sqrt{10} = 0\)}{}{}}
\LANGsk{
\Question{
Nájdite všetky priamky, ktoré sú rovnobežné s priamkou
\(p\colon x - 3y + 2 = 0\) a majú od nej
vzdialenosť \(\sqrt{10}\).}
{\(p_{1}\colon x - 3y + 12 = 0\),
\(p_{2}\colon x - 3y - 8 = 0\)}{\(p\colon x - 3y = 0\)}{\(p\colon x - 3y + \sqrt{10} = 0\)}{\(p_{1}\colon x - 3y + \sqrt{10} = 0\),
\(p_{2}\colon x - 3y -\sqrt{10} = 0\)}{}{}}
\LANGes{
\Question{
Determina todas las rectas que son paralelas a la recta \(p\colon x - 3y + 2 = 0\) 
y la distancia de cada una de estas rectas a la recta \(p\) es
\(\sqrt{10}\).}
{\(p_{1}\colon x - 3y + 12 = 0\),
\(p_{2}\colon x - 3y - 8 = 0\)}{\(p\colon x - 3y = 0\)}{\(p\colon x - 3y + \sqrt{10} = 0\)}{\(p_{1}\colon x - 3y + \sqrt{10} = 0\),
\(p_{2}\colon x - 3y -\sqrt{10} = 0\)}{}{}}
\End

\Data\ID{9000149410}
\Updated{2020-07-17 05:07:39}
\Level{B}
\SubArea{Analytic geometry in a plane}
\LANGen{
\Question{
Find all lines passing through the point
\(A = [-2;-6]\) such that the distance
from the point \([0.0]\) 
to these lines is \(2\sqrt{2}\).}
{\(p_{1}\colon 7x + y + 20 = 0\),
\(p_{2}\colon x - y - 4 = 0\)}{\(p\colon 7x - y = 0\)}{\(p\colon x + y + 2\sqrt{2} = 0\)}{\(p_{1}\colon x - y + 2\sqrt{2} = 0\),
\(p_{2}\colon x + y - 2\sqrt{2} = 0\)}{}{}}
\LANGpl{
\Question{
Wyznacz wszystkie proste przechodzące przez punkt 
\(A = [-2;-6]\) tak, aby odległość
od punktu \([0{,}0]\) 
do tych prostych  była równa \(2\sqrt{2}\).}
{\(p_{1}\colon 7x + y + 20 = 0\),
\(p_{2}\colon x - y - 4 = 0\)}{\(p\colon 7x - y = 0\)}{\(p\colon x + y + 2\sqrt{2} = 0\)}{\(p_{1}\colon x - y + 2\sqrt{2} = 0\),
\(p_{2}\colon x + y - 2\sqrt{2} = 0\)}{}{}}
\LANGcs{
\Question{
Určete rovnice všech přímek, které prochází bodem
\(A = [-2;-6]\) 
a jejichž vzdálenost od počátku soustavy souřadnic je
\(2\sqrt{2}\).}
{\(p_{1}\colon 7x + y + 20 = 0\),
\(p_{2}\colon x - y - 4 = 0\)}{\(p\colon 7x - y = 0\)}{\(p\colon x + y + 2\sqrt{2} = 0\)}{\(p_{1}\colon x - y + 2\sqrt{2} = 0\),
\(p_{2}\colon x + y - 2\sqrt{2} = 0\)}{}{}}
\LANGsk{
\Question{
Určte rovnice všetkých priamok, ktoré prechádzajú bodom
\(A = [-2;-6]\) 
a ich vzdialenosť od začiatku sústavy súradníc je
\(2\sqrt{2}\).}
{\(p_{1}\colon 7x + y + 20 = 0\),
\(p_{2}\colon x - y - 4 = 0\)}{\(p\colon 7x - y = 0\)}{\(p\colon x + y + 2\sqrt{2} = 0\)}{\(p_{1}\colon x - y + 2\sqrt{2} = 0\),
\(p_{2}\colon x + y - 2\sqrt{2} = 0\)}{}{}}
\LANGes{
\Question{
Halla todas las rectas que pasan por el punto
\(A = [-2;-6]\) suponiendo que la distancia del punto \([0.0]\) 
a las rectas es \(2\sqrt{2}\).}
{\(p_{1}\colon 7x + y + 20 = 0\),
\(p_{2}\colon x - y - 4 = 0\)}{\(p\colon 7x - y = 0\)}{\(p\colon x + y + 2\sqrt{2} = 0\)}{\(p_{1}\colon x - y + 2\sqrt{2} = 0\),
\(p_{2}\colon x + y - 2\sqrt{2} = 0\)}{}{}}
\End

\Data\ID{9000151301}
\Updated{2022-04-23 05:04:00}
\Level{B}
\SubArea{Analytic geometry in a plane}
\LANGen{
\Question{
Find the angle \(\varphi \) 
between the lines \(3x - 7 = 0\) 
and \(x + y + 13 = 0\).}
{\(45^{\circ }\)}{\(90^{\circ }\)}{\(60^{\circ }\)}{\(30^{\circ }\)}{}{}}
\LANGpl{
\Question{
Wyznacz kąt  \(\varphi \) 
między  prostymi \(3x - 7 = 0\) 
i \(x + y + 13 = 0\).}
{\(45^{\circ }\)}{\(90^{\circ }\)}{\(60^{\circ }\)}{\(30^{\circ }\)}{}{}}
\LANGcs{
\Question{
Určete odchylku \(\varphi \) 
přímek zadaných obecnými rovnicemi
\(3x - 7 = 0\) a
\(x + y + 13 = 0\).}
{\(45^{\circ }\)}{\(90^{\circ }\)}{\(60^{\circ }\)}{\(30^{\circ }\)}{}{}}
\LANGsk{
\Question{
Určte odchýlku \(\varphi \) 
priamok zadaných všeobecnými rovnicami
\(3x - 7 = 0\) a
\(x + y + 13 = 0\).}
{\(45^{\circ }\)}{\(90^{\circ }\)}{\(60^{\circ }\)}{\(30^{\circ }\)}{}{}}
\LANGes{
\Question{
Determina el ángulo \(\varphi \) 
entre las rectas \(3x - 7 = 0\) 
y \(x + y + 13 = 0\).}
{\(45^{\circ }\)}{\(90^{\circ }\)}{\(60^{\circ }\)}{\(30^{\circ }\)}{}{}}
\End

\Data\ID{9000151302}
\Updated{2020-07-17 05:07:32}
\Level{B}
\SubArea{Analytic geometry in a plane}
\LANGen{
\Question{
Find the angle \(\varphi \) 
between the lines \(p\) 
and \(q\) 
given by their parametric equations. 
\[ 
p\colon \begin{aligned}[t] x& = 1 + 2t, &
\\y& = 3 - 3t;\ t\in \mathbb{R},
\\ \end{aligned}\qquad q\colon \begin{aligned}[t] x& = 2 - k, &
\\y& = 3 + k;\ k\in \mathbb{R}
\\ \end{aligned}
\]}
{\(11^{\circ }19'\)}{\(88^{\circ }41'\)}{\(45^{\circ }45'\)}{\(54^{\circ }12'\)}{}{}}
\LANGpl{
\Question{
Wyznacz kąt \(\varphi \) 
między prostymi \(p\) 
i \(q\) 
spełniającymi równania parametryczne.
\[ 
p\colon \begin{aligned}[t] x& = 1 + 2t, &
\\y& = 3 - 3t;\ t\in \mathbb{R},
\\ \end{aligned}\qquad q\colon \begin{aligned}[t] x& = 2 - k, &
\\y& = 3 + k;\ k\in \mathbb{R}
\\ \end{aligned}
\]}
{\(11^{\circ }19'\)}{\(88^{\circ }41'\)}{\(45^{\circ }45'\)}{\(54^{\circ }12'\)}{}{}}
\LANGcs{
\Question{
Určete odchylku \(\varphi \) 
přímek zadaných parametricky 
\[ 
p\colon \begin{aligned}[t] x& = 1 + 2t, &
\\y& = 3 - 3t;\ t\in \mathbb{R},
\\ \end{aligned}\qquad q\colon \begin{aligned}[t] x& = 2 - k, &
\\y& = 3 + k;\ k\in \mathbb{R}
\\ \end{aligned}
\]}
{\(11^{\circ }19'\)}{\(88^{\circ }41'\)}{\(45^{\circ }45'\)}{\(54^{\circ }12'\)}{}{}}
\LANGsk{
\Question{
Určte odchýlku \(\varphi \) 
priamok zadaných parametricky 
\[ 
p\colon \begin{aligned}[t] x& = 1 + 2t, &
\\y& = 3 - 3t;\ t\in \mathbb{R},
\\ \end{aligned}\qquad q\colon \begin{aligned}[t] x& = 2 - k, &
\\y& = 3 + k;\ k\in \mathbb{R}
\\ \end{aligned}
\]}
{\(11^{\circ }19'\)}{\(88^{\circ }41'\)}{\(45^{\circ }45'\)}{\(54^{\circ }12'\)}{}{}}
\LANGes{
\Question{
Determina el ángulo \(\varphi \) 
entre las rectas paramétricas \(p\) 
y \(q\) . 
\[ 
p\colon \begin{aligned}[t] x& = 1 + 2t, &
\\y& = 3 - 3t;\ t\in \mathbb{R},
\\ \end{aligned}\qquad q\colon \begin{aligned}[t] x& = 2 - k, &
\\y& = 3 + k;\ k\in \mathbb{R}
\\ \end{aligned}
\]}
{\(11^{\circ }19'\)}{\(88^{\circ }41'\)}{\(45^{\circ }45'\)}{\(54^{\circ }12'\)}{}{}}
\End

\Data\ID{9000151303}
\Updated{2021-12-05 07:12:07}
\Level{B}
\SubArea{Analytic geometry in a plane}
\LANGen{
\Question{
Find the angle \(\varphi \) 
between the lines given by their slope-intercept equations
\(y = 6\) and
\(y = \frac{3} 
{4}x\).}
{\(36^{\circ }52'\)}{\(45^{\circ }59'\)}{\(64^{\circ }33'\)}{\(76^{\circ }11'\)}{}{}}
\LANGpl{
\Question{
Wyznacz kąt\(\varphi \) 
między prostymi o równaniach w postaci kierunkowej
\(y = 6\) i
\(y = \frac{3} 
{4}x\).}
{\(36^{\circ }52'\)}{\(45^{\circ }59'\)}{\(64^{\circ }33'\)}{\(76^{\circ }11'\)}{}{}}
\LANGcs{
\Question{
Určete odchylku \(\varphi \) 
přímek zadaných rovnicemi ve směrnicovém tvaru
\(y = 6\) a
\(y = \frac{3} 
{4}x\).}
{\(36^{\circ }52'\)}{\(45^{\circ }59'\)}{\(64^{\circ }33'\)}{\(76^{\circ }11'\)}{}{}}
\LANGsk{
\Question{
Určte odchýlku \(\varphi \) 
priamok zadaných rovnicami v smernicovom tvare
\(y = 6\) a
\(y = \frac{3} 
{4}x\).}
{\(36^{\circ }52'\)}{\(45^{\circ }59'\)}{\(64^{\circ }33'\)}{\(76^{\circ }11'\)}{}{}}
\LANGes{
\Question{
Determina el ángulo \(\varphi \)
entre las rectas dadas por sus ecuaciones explícitas
\(y = 6\) y
\(y = \frac{3} 
{4}x\).}
{\(36^{\circ }52'\)}{\(45^{\circ }59'\)}{\(64^{\circ }33'\)}{\(76^{\circ }11'\)}{}{}}
\End

\Data\ID{9000151304}
\Updated{2020-07-17 05:07:50}
\Level{B}
\SubArea{Analytic geometry in a plane}
\LANGen{
\Question{
Find the angle \(\varphi \) 
between the lines \(y = 0\) 
and \(\frac{x}
{2} + \frac{y} 
{3} = 1\).}
{\(56^{\circ }19'\)}{\(13^{\circ }45'\)}{\(26^{\circ }46'\)}{\(81^{\circ }23'\)}{}{}}
\LANGpl{
\Question{
Wyznacz kąt\(\varphi \) 
pomiędzy prostymi \(y = 0\) 
i \(\frac{x}
{2} + \frac{y} 
{3} = 1\).}
{\(56^{\circ }19'\)}{\(13^{\circ }45'\)}{\(26^{\circ }46'\)}{\(81^{\circ }23'\)}{}{}}
\LANGcs{
\Question{
Určete odchylku \(\varphi \) 
přímky zadané rovnicí ve směrnicovém tvaru
\(y = 0\) 
a přímky zadané rovnicí v úsekovém tvaru
\(\frac{x}
{2} + \frac{y} 
{3} = 1\).}
{\(56^{\circ }19'\)}{\(13^{\circ }45'\)}{\(26^{\circ }46'\)}{\(81^{\circ }23'\)}{}{}}
\LANGsk{
\Question{
Určte odchýlku \(\varphi \) 
priamky zadanej rovnicou v smernicovom tvare
\(y = 0\) 
a priamky zadanej rovnicou v úsekovom tvare
\(\frac{x}
{2} + \frac{y} 
{3} = 1\).}
{\(56^{\circ }19'\)}{\(13^{\circ }45'\)}{\(26^{\circ }46'\)}{\(81^{\circ }23'\)}{}{}}
\LANGes{
\Question{
Determina el ángulo \(\varphi \) 
entre las rectas \(y = 0\) 
y \(\frac{x}
{2} + \frac{y} 
{3} = 1\).}
{\(56^{\circ }19'\)}{\(13^{\circ }45'\)}{\(26^{\circ }46'\)}{\(81^{\circ }23'\)}{}{}}
\End

\Data\ID{9000151305}
\Updated{2021-12-05 07:12:28}
\Level{B}
\SubArea{Analytic geometry in a plane}
\LANGen{
\Question{
Find the angle \(\varphi \) 
between the lines \(x + y + 1 = 0\) 
and \(x - y - 1 = 0\).}
{\(90^{\circ }\)}{\(0^{\circ }\)}{\(30^{\circ }\)}{\(60^{\circ }\)}{}{}}
\LANGpl{
\Question{
Wyznacz kąt\(\varphi \) 
pomiędzy prostymi \(x + y + 1 = 0\) 
i \(x - y - 1 = 0\).}
{\(90^{\circ }\)}{\(0^{\circ }\)}{\(30^{\circ }\)}{\(60^{\circ }\)}{}{}}
\LANGcs{
\Question{
Určete odchylku \(\varphi \) 
přímek zadaných obecnými rovnicemi
\(x + y + 1 = 0\) a
\(x - y - 1 = 0\).}
{\(90^{\circ }\)}{\(0^{\circ }\)}{\(30^{\circ }\)}{\(60^{\circ }\)}{}{}}
\LANGsk{
\Question{
Určte odchýlku \(\varphi \) 
priamok zadaných všeobecnými rovnicami
\(x + y + 1 = 0\) a
\(x - y - 1 = 0\).}
{\(90^{\circ }\)}{\(0^{\circ }\)}{\(30^{\circ }\)}{\(60^{\circ }\)}{}{}}
\LANGes{
\Question{
Determina el ángulo \(\varphi \) 
entre las rectas \(x + y + 1 = 0\) 
y  \(x - y - 1 = 0\).}
{\(90^{\circ }\)}{\(0^{\circ }\)}{\(30^{\circ }\)}{\(60^{\circ }\)}{}{}}
\End

\Data\ID{9000151306}
\Updated{2021-12-05 07:12:05}
\Level{B}
\SubArea{Analytic geometry in a plane}
\LANGen{
\Question{
Find the angle \(\varphi \) 
between the lines \(p\) 
and \(q\) 
given by their parametric equations. 

\[ 
p\colon \begin{aligned}[t] x& = 1 - t, &
\\y& = 2 + t;\ t\in \mathbb{R},
\\ \end{aligned}\qquad q\colon \begin{aligned}[t] x& = 4 - k, &
\\y& = 5 + k;\ k\in \mathbb{R}.
\\ \end{aligned}
\]}
{\(0^{\circ }\)}{\(90^{\circ }\)}{\(60^{\circ }\)}{\(30^{\circ }\)}{}{}}
\LANGpl{
\Question{
Wyznacz kąt \(\varphi \) między prostymi \(p\) 
i \(q\) 
spełniającymi równania parametryczne 

\[ 
p\colon \begin{aligned}[t] x& = 1 - t, &
\\y& = 2 + t;\ t\in \mathbb{R},
\\ \end{aligned}\qquad q\colon \begin{aligned}[t] x& = 4 - k, &
\\y& = 5 + k;\ k\in \mathbb{R}.
\\ \end{aligned}
\]}
{\(0^{\circ }\)}{\(90^{\circ }\)}{\(60^{\circ }\)}{\(30^{\circ }\)}{}{}}
\LANGcs{
\Question{
Určete odchylku \(\varphi \) 
přímek zadaných parametricky 

\[ 
p\colon \begin{aligned}[t] x& = 1 - t, &
\\y& = 2 + t;\ t\in \mathbb{R},
\\ \end{aligned}\qquad q\colon \begin{aligned}[t] x& = 4 - k, &
\\y& = 5 + k;\ k\in \mathbb{R}.
\\ \end{aligned}
\]}
{\(0^{\circ }\)}{\(90^{\circ }\)}{\(60^{\circ }\)}{\(30^{\circ }\)}{}{}}
\LANGsk{
\Question{
Určte odchýlku \(\varphi \) 
priamok zadaných parametricky 
\[ 
p\colon \begin{aligned}[t] x& = 1 - t, &
\\y& = 2 + t;\ t\in \mathbb{R},
\\ \end{aligned}\qquad q\colon \begin{aligned}[t] x& = 4 - k, &
\\y& = 5 + k;\ k\in \mathbb{R}.
\\ \end{aligned}
\]}
{\(0^{\circ }\)}{\(90^{\circ }\)}{\(60^{\circ }\)}{\(30^{\circ }\)}{}{}}
\LANGes{
\Question{
Determina el ángulo \(\varphi \) 
entre las rectas cuyas ecuaciones paramétricas son \(p\) 
y \(q\).

\[ 
p\colon \begin{aligned}[t] x& = 1 - t, &
\\y& = 2 + t;\ t\in \mathbb{R},
\\ \end{aligned}\qquad q\colon \begin{aligned}[t] x& = 4 - k, &
\\y& = 5 + k;\ k\in \mathbb{R}.
\\ \end{aligned}
\]}
{\(0^{\circ }\)}{\(90^{\circ }\)}{\(60^{\circ }\)}{\(30^{\circ }\)}{}{}}
\End

\Data\ID{9000151307}
\Updated{2021-12-05 07:12:00}
\Level{B}
\SubArea{Analytic geometry in a plane}
\LANGen{
\Question{
Find the angle \(\varphi \) 
between the line \(x + \sqrt{3}y - 6 = 0\) 
and the line \(p\) 
given by it's parametric equations. 
\[ 
p\colon \begin{aligned}[t] x& = 2 + t,&
\\y& = 5;\ t\in \mathbb{R}
\\ \end{aligned}
\]}
{\(30^{\circ }\)}{\(90^{\circ }\)}{\(60^{\circ }\)}{\(45^{\circ }\)}{}{}}
\LANGpl{
\Question{
Wyznacz kąt  \(\varphi \) 
pomiędzy prostymi \(x + \sqrt{3}y - 6 = 0\) 
i \(p\) spełniającymi równania parametryczne. 

\[ 
p\colon \begin{aligned}[t] x& = 2 + t,&
\\y& = 5;\ t\in \mathbb{R}
\\ \end{aligned}
\]}
{\(30^{\circ }\)}{\(90^{\circ }\)}{\(60^{\circ }\)}{\(45^{\circ }\)}{}{}}
\LANGcs{
\Question{
Určete odchylku \(\varphi \) přímky
zadané obecnou rovnicí \(x + \sqrt{3}y - 6 = 0\) 
a přímky zadané parametrickými rovnicemi 

\[ 
p\colon \begin{aligned}[t] x& = 2 + t,&
\\y& = 5;\ t\in \mathbb{R}.
\\ \end{aligned}
\]}
{\(30^{\circ }\)}{\(90^{\circ }\)}{\(60^{\circ }\)}{\(45^{\circ }\)}{}{}}
\LANGsk{
\Question{
Určte odchýlku \(\varphi \) priamky
zadanej všeobecnou rovnicou \(x + \sqrt{3}y - 6 = 0\) 
a priamky zadanej parametrickými rovnicami 
\[ 
p\colon \begin{aligned}[t] x& = 2 + t,&
\\y& = 5;\ t\in \mathbb{R}.
\\ \end{aligned}
\]}
{\(30^{\circ }\)}{\(90^{\circ }\)}{\(60^{\circ }\)}{\(45^{\circ }\)}{}{}}
\LANGes{
\Question{
Calcula el ángulo \(\varphi \) 
entre la recta \(x + \sqrt{3}y - 6 = 0\) 
y la recta  \(p\) 
que viene dada en forma paramétrica. 
\[ 
p\colon \begin{aligned}[t] x& = 2 + t,&
\\y& = 5;\ t\in \mathbb{R}
\\ \end{aligned}
\]}
{\(30^{\circ }\)}{\(90^{\circ }\)}{\(60^{\circ }\)}{\(45^{\circ }\)}{}{}}
\End

\Data\ID{9000151308}
\Updated{2022-04-23 05:04:00}
\Level{B}
\SubArea{Analytic geometry in a plane}
\LANGen{
\Question{
Given points \(A = [-1;4]\),
\(B = [2,-2]\) and 
\(C = [5,-1]\), find the angle
\(\beta \) (the interior angle
at the vertex \(B\))
in the triangle \(ABC\).}
{\(98^{\circ }08'\)}{\(81^{\circ }53'\)}{\(76^{\circ }17'\)}{\(68^{\circ }27'\)}{}{}}
\LANGpl{
\Question{
Dany jest trójkąt \(ABC\), gdzie \(A = [-1{,}4]\),
\(B = [2,-2]\),
\(C = [5,-1]\), wyznacz miarę kąta 
\(\beta \) przy wierzchołku 
 \(B\))}
{\(98^{\circ }08'\)}{\(81^{\circ }53'\)}{\(76^{\circ }17'\)}{\(68^{\circ }27'\)}{}{}}
\LANGcs{
\Question{
Je dán trojúhelník \(ABC\),
\(A = [-1{,}4]\),
\(B = [2,-2]\),
\(C = [5,-1]\). Vypočítejte velikost
vnitřního úhlu \(\beta \) 
u vrcholu \(B\) v
trojúhelníku \(ABC\).}
{\(98^{\circ }08'\)}{\(81^{\circ }53'\)}{\(76^{\circ }17'\)}{\(68^{\circ }27'\)}{}{}}
\LANGsk{
\Question{
Je daný trojuholník \(ABC\),
\(A = [-1{,}4]\),
\(B = [2,-2]\),
\(C = [5,-1]\). Vypočítajte veľkosť
vnútorného uhla \(\beta \) 
u vrchola \(B\) v
trojuholníku \(ABC\).}
{\(98^{\circ }08'\)}{\(81^{\circ }53'\)}{\(76^{\circ }17'\)}{\(68^{\circ }27'\)}{}{}}
\LANGes{
\Question{
Dados los puntos \(A = [-1.4]\),
\(B = [2,-2]\),
\(C = [5,-1]\), halla el ángulo
\(\beta \) (el ángulo interior del vértice \(B\))
en el triángulo \(ABC\).}
{\(98^{\circ }08'\)}{\(81^{\circ }53'\)}{\(76^{\circ }17'\)}{\(68^{\circ }27'\)}{}{}}
\End

\Data\ID{9000151309}
\Updated{2020-07-17 05:07:07}
\Level{B}
\SubArea{Analytic geometry in a plane}
\LANGen{
\Question{
Given points \(A = [-1.4]\),
\(B = [2,-2]\),
\(C = [5,-1]\), find the angle
\(\varphi \) between
the lines \(AB\) 
and \(BC\).}
{\(81^{\circ }52'\)}{\(98^{\circ }08'\)}{\(61^{\circ }22'\)}{\(45^{\circ }32'\)}{}{}}
\LANGpl{
\Question{
Dane są punkty \(A = [-1{,}4]\),
\(B = [2,-2]\),
\(C = [5,-1]\), wyznacz miarę kąta
\(\varphi \) pomiędzy 
prostymi  \(AB\) 
i \(BC\).}
{\(81^{\circ }52'\)}{\(98^{\circ }08'\)}{\(61^{\circ }22'\)}{\(45^{\circ }32'\)}{}{}}
\LANGcs{
\Question{
Je dán trojúhelník \(ABC\),
\(A = [-1{,}4]\),
\(B = [2,-2]\),
\(C = [5,-1]\). Vypočítejte
odchylku \(\varphi \) 
přímek \(AB\),
\(BC\).}
{\(81^{\circ }52'\)}{\(98^{\circ }08'\)}{\(61^{\circ }22'\)}{\(45^{\circ }32'\)}{}{}}
\LANGsk{
\Question{
Je daný trojuholník \(ABC\),
\(A = [-1{,}4]\),
\(B = [2,-2]\),
\(C = [5,-1]\). Vypočítajte
odchýlku \(\varphi \) 
priamok \(AB\),
\(BC\).}
{\(81^{\circ }52'\)}{\(98^{\circ }08'\)}{\(61^{\circ }22'\)}{\(45^{\circ }32'\)}{}{}}
\LANGes{
\Question{
Dados los puntos \(A = [-1.4]\),
\(B = [2,-2]\),
\(C = [5,-1]\), halla el ángulo
\(\varphi \) entre las rectas \(AB\) 
y \(BC\).}
{\(81^{\circ }52'\)}{\(98^{\circ }08'\)}{\(61^{\circ }22'\)}{\(45^{\circ }32'\)}{}{}}
\End

\Data\ID{1103109101}
\Updated{2021-12-05 10:12:27}
\Level{C}
\SubArea{Analytic geometry in a plane}
\IMG{1103109101.pdf}
\LANGen{
\Question{
Find equations of all lines at the distance \( \sqrt{10} \) from the point \( M=[5;4] \) which are perpendicular to the line \( p \) with the equation \( 2x+6y-3=0 \) (see the picture).}
{\( 3x-y-1=0;\ 3x-y-21=0 \)}{\( 3x-y+1=0;\ 3x-y-18=0 \)}{\( x+3y+1=0;\ x+3y+21=0 \)}{\( x+3y-1=0;\ x+3y-18=0 \)}{}{}}
\LANGpl{
\Question{
Wyznacz równania wszystkich prostych w odległości \( \sqrt{10} \) od punktu \( M=[5;4] \), prostopadłych do prostej  \( p \) o równaniu  \( 2x+6y-3=0 \) (spójrz na rysunek).}
{\( 3x-y-1=0;\ 3x-y-21=0 \)}{\( 3x-y+1=0;\ 3x-y-18=0 \)}{\( x+3y+1=0;\ x+3y+21=0 \)}{\( x+3y-1=0;\ x+3y-18=0 \)}{}{}}
\LANGcs{
\Question{
Určete obecné rovnice všech přímek, které jsou kolmé k přímce \( p \): \( 2x+6y-3=0 \) a mají od bodu \( M=[5;4] \) vzdálenost rovnou \( \sqrt{10} \) (viz obrázek).}
{\( 3x-y-1=0;\ 3x-y-21=0 \)}{\( 3x-y+1=0;\ 3x-y-18=0 \)}{\( x+3y+1=0;\ x+3y+21=0 \)}{\( x+3y-1=0;\ x+3y-18=0 \)}{}{}}
\LANGsk{
\Question{
Nájdite rovnice všetkých priamok so vzdialenosťou \( \sqrt{10} \) z bodu \(M = [5; 4] \), ktoré sú kolmé na priamku \( p \) s rovnicou \( 2x+6y-3=0  \) (pozri obrázok).}
{\( 3x-y-1=0;\ 3x-y-21=0 \)}{\( 3x-y+1=0;\ 3x-y-18=0 \)}{\( x+3y+1=0;\ x+3y+21=0 \)}{\( x+3y-1=0;\ x+3y-18=0 \)}{}{}}
\LANGes{
\Question{
Determina las ecuaciones generales de todas las rectas que son perpendiculares a la recta \( p \): \( 2x+6y-3=0 \) y cuya distancia al punto \( M=[5;4] \) es  \( \sqrt{10} \) (mira la imagen).}
{\( 3x-y-1=0;\ 3x-y-21=0 \)}{\( 3x-y+1=0;\ 3x-y-18=0 \)}{\( x+3y+1=0;\ x+3y+21=0 \)}{\( x+3y-1=0;\ x+3y-18=0 \)}{}{}}
\End

\Data\ID{1103109102}
\Updated{2021-12-05 10:12:31}
\Level{C}
\SubArea{Analytic geometry in a plane}
\IMG{1103109102.pdf}
\LANGen{
\Question{
Let \( p \) and \( q \) be intersecting lines with the equations \( y=\frac{\sqrt3}3x \) and \( x=0 \) respectively. Find equations of lines \( o_1 \) and \( o_2 \) that are lines of symmetry of the angles contained between \( p \) and \( q \) (see the picture).}
{\( y=\sqrt3x;\  y=-\frac{\sqrt3}3x \)}{\( y=2x;\ y=-\frac12x \)}{\( y=\sqrt2x;\ y=-\frac{\sqrt2}2x \)}{\( y=3x;\ y=-\frac13x \)}{}{}}
\LANGpl{
\Question{
Dane są proste przecinające się  \( p \) i \( q \), proste są określone równaniami \( y=\frac{\sqrt3}3x \) i  \( x=0 \).  Wyznacz równania prostych  \( o_1 \) i \( o_2 \), które są osiami symetrii kątów pomiędzy prostymi  \( p \) i \( q \) (spójrz na rysunek).}
{\( y=\sqrt3x;\  y=-\frac{\sqrt3}3x \)}{\( y=2x;\ y=-\frac12x \)}{\( y=\sqrt2x;\ y=-\frac{\sqrt2}2x \)}{\( y=3x;\ y=-\frac13x \)}{}{}}
\LANGcs{
\Question{
Jsou dány přímky \( p \): \( y=\frac{\sqrt3}3x \) a \( q \): \( x=0 \). Určete rovnice přímek  \( o_1 \) a \( o_2 \), které jsou osami úhlů různoběžek \( p \) a \( q \) (viz obrázek).}
{\( y=\sqrt3x;\  y=-\frac{\sqrt3}3x \)}{\( y=2x;\ y=-\frac12x \)}{\( y=\sqrt2x;\ y=-\frac{\sqrt2}2x \)}{\( y=3x;\ y=-\frac13x \)}{}{}}
\LANGsk{
\Question{
Sú dané priamky \( p \): \( y=\frac{\sqrt3}3x \) a \( q \): \( x=0 \). Určte rovnice priamok  \( o_1 \) a \( o_2 \), ktoré sú osami uhlov rôznobežiek \( p \) a \( q \) (viď obrázok).}
{\( y=\sqrt3x;\  y=-\frac{\sqrt3}3x \)}{\( y=2x;\ y=-\frac12x \)}{\( y=\sqrt2x;\ y=-\frac{\sqrt2}2x \)}{\( y=3x;\ y=-\frac13x \)}{}{}}
\LANGes{
\Question{
Sea \( p \)  y \( q \) rectas secantes con ecuaciones \( y=\frac{\sqrt3}3x \) y \( x=0 \) respectivamente. Determina la ecuación de las rectas \( o_1 \) y \( o_2 \) que son ejes de simetría de los ángulos  entre \( p \) y \( q \) (mira la imagen).}
{\( y=\sqrt3x;\  y=-\frac{\sqrt3}3x \)}{\( y=2x;\ y=-\frac12x \)}{\( y=\sqrt2x;\ y=-\frac{\sqrt2}2x \)}{\( y=3x;\ y=-\frac13x \)}{}{}}
\End

\Data\ID{1103109103}
\Updated{2021-12-05 10:12:57}
\Level{C}
\SubArea{Analytic geometry in a plane}
\IMG{1103109103.pdf}
\LANGen{
\Question{
Let \( y=-\frac{\sqrt3}3x+1 \) be the equation of the line \( p \) and let \( M \) be the point \( [0;-3] \). Find equations of all lines passing through \( M \) and intersecting \( p \) at an angle of \( 60^{\circ} \) (see the picture).}
{\( x=0;\ y=\frac{\sqrt3}3x-3 \)}{\( y=0;\ y=\frac{\sqrt3}3x-3 \)}{\( y=0;\ y=x-3 \)}{\( x=0;\  y=\sqrt3x-3 \)}{}{}}
\LANGpl{
\Question{
Dane jest równanie \( y=-\frac{\sqrt3}3x+1 \) prostej  \( p \) oraz punkt \( M \) o współrzędnych \( [0;-3] \). Wyznacz równania wszystkich prostych przechodzących przez punkt \( M \) oraz przecinających \( p \) pod kątem \( 60^{\circ} \) (spójrz na rysunek).}
{\( x=0;\ y=\frac{\sqrt3}3x-3 \)}{\( y=0;\ y=\frac{\sqrt3}3x-3 \)}{\( y=0;\ y=x-3 \)}{\( x=0;\  y=\sqrt3x-3 \)}{}{}}
\LANGcs{
\Question{
Najděte rovnice všech přímek, které procházejí bodem \( M \) = \( [0;-3] \) a mají od přímky  \( p \): \( y=-\frac{\sqrt3}3x+1 \) odchylku \( 60^{\circ} \) (viz obrázek).}
{\( x=0;\ y=\frac{\sqrt3}3x-3 \)}{\( y=0;\ y=\frac{\sqrt3}3x-3 \)}{\( y=0;\ y=x-3 \)}{\( x=0;\  y=\sqrt3x-3 \)}{}{}}
\LANGsk{
\Question{
Nájdite rovnice všetkých priamok, ktoré prechádzajú bodom \( M \) = \( [0;-3] \) a majú od priamky  \( p \): \( y=-\frac{\sqrt3}3x+1 \) odchýlku \( 60^{\circ} \) (viď obrázok).}
{\( x=0;\ y=\frac{\sqrt3}3x-3 \)}{\( y=0;\ y=\frac{\sqrt3}3x-3 \)}{\( y=0;\ y=x-3 \)}{\( x=0;\  y=\sqrt3x-3 \)}{}{}}
\LANGes{
\Question{
Sea \( y=-\frac{\sqrt3}3x+1 \) la ecuación de la recta \( p \) y sea \( M \) el punto \( [0;-3] \). Determina las ecuaciones de todas las rectas que pasan por el punto \( M \) y cortan a la recta \( p \) formando un ángulo de \( 60^{\circ} \) (mira la imagen).}
{\( x=0;\ y=\frac{\sqrt3}3x-3 \)}{\( y=0;\ y=\frac{\sqrt3}3x-3 \)}{\( y=0;\ y=x-3 \)}{\( x=0;\  y=\sqrt3x-3 \)}{}{}}
\End

\Data\ID{1103109104}
\Updated{2021-12-05 10:12:20}
\Level{C}
\SubArea{Analytic geometry in a plane}
\IMG{1103109104.pdf}
\LANGen{
\Question{
Let \( 2x-3y+6=0 \) be the equation of the line \( p \) and let \( M \) be the point \( [5;3] \). Find equations of all lines passing through \( M \) and intersecting \( p \) at an angle of \( 45^{\circ} \) (see the picture).}
{\( x+5y-20=0;\ 5x-y-22=0 \)}{\( x+6y-23=0;\ 6x-y-27=0 \)}{\( x+4y-17=0;\ 4x-y-16=0 \)}{\( x+5y-28=0;\ 5x-y-10=0 \)}{}{}}
\LANGpl{
\Question{
\( 2x-3y+6=0 \) to równanie prostej \( p \),  \( M \) to punkt o współrzędnych  \( [5;3] \). Wyznacz równania prostych przechodzących przez punkt \( M \) oraz przecinających  \( p \) pod kątem  \( 45^{\circ} \) (spójrz na rysunek).}
{\( x+5y-20=0;\ 5x-y-22=0 \)}{\( x+6y-23=0;\ 6x-y-27=0 \)}{\( x+4y-17=0;\ 4x-y-16=0 \)}{\( x+5y-28=0;\ 5x-y-10=0 \)}{}{}}
\LANGcs{
\Question{
Najděte rovnice všech přímek, které procházejí bodem \( M \)\( = [5;3] \) a mají od přímky \( p \): \( 2x-3y+6=0 \) odchylku \( 45^{\circ} \) (viz obrázek).}
{\( x+5y-20=0;\ 5x-y-22=0 \)}{\( x+6y-23=0;\ 6x-y-27=0 \)}{\( x+4y-17=0;\ 4x-y-16=0 \)}{\( x+5y-28=0;\ 5x-y-10=0 \)}{}{}}
\LANGsk{
\Question{
Nájdite rovnice všetkých priamok, ktoré prechádzajú bodom \( M \)\( = [5;3] \) a majú od priamky \( p \): \( 2x-3y+6=0 \) odchýlku \( 45^{\circ} \) (viď obrázok).}
{\( x+5y-20=0;\ 5x-y-22=0 \)}{\( x+6y-23=0;\ 6x-y-27=0 \)}{\( x+4y-17=0;\ 4x-y-16=0 \)}{\( x+5y-28=0;\ 5x-y-10=0 \)}{}{}}
\LANGes{
\Question{
Sea \( 2x-3y+6=0 \) la ecuación de la recta \( p \) y sea \( M \) el punto \( [5;3] \).  Determina las ecuaciones de todas las rectas que pasan por el punto \( M \) y cortan a la recta  \( p \) formando un ángulo de \( 45^{\circ} \) (mira la imagen).}
{\( x+5y-20=0;\ 5x-y-22=0 \)}{\( x+6y-23=0;\ 6x-y-27=0 \)}{\( x+4y-17=0;\ 4x-y-16=0 \)}{\( x+5y-28=0;\ 5x-y-10=0 \)}{}{}}
\End

\Data\ID{1103109105}
\Updated{2021-12-05 10:12:29}
\Level{C}
\SubArea{Analytic geometry in a plane}
\IMG{1103109105.pdf}
\LANGen{
\Question{
Let \( p \) and \( q \) be the lines with the equations \( x-2y-1=0 \) and \( 2x+y-12=0 \) respectively. Find all the points at the same distance of \( \sqrt5 \) from \( p \) and \( q \)  (see the picture).}
{\([2;3] \), \([6;5] \), \([8;1] \), \([4;-1] \)}{\([2;3] \), \([6;5] \), \([8.5;1] \), \([4.5;-1] \)}{\([2;3.5] \), \([6;5.5] \), \([8;1] \), \([4;-1] \)}{\([2;3] \), \([6;5.5] \), \([8;1.5] \), \([4;-1] \)}{}{}}
\LANGpl{
\Question{
Dane są proste \( p \) i \( q \) określone równaniami \( x-2y-1=0 \) i \( 2x+y-12=0 \) . Wyznacz wszystkie punkty w odległości \( \sqrt5 \) od \( p \) i \( q \)  (spójrz na rysunek).}
{\( [2;3] \), \( [6;5] \), \( [8;1] \), \( [4;-1] \)}{\( [2;3] \), \( [6;5] \), \( [8{,}5;1] \), \( [4{,}5;-1] \)}{\(  [2;3{,}5] \), \(  [6;5{,}5] \), \(  [8;1] \), \( [4;-1] \)}{\( [2;3] \), \( [6;5{,}5] \), \( [8;1{,}5] \), \( [4;-1] \)}{}{}}
\LANGcs{
\Question{
Jsou dány přímky \( p \): \( x-2y-1=0 \) a \( q \): \( 2x+y-12=0 \). Určete souřadnice všech takových bodů, které mají od obou daných přímek vzdálenost \( \sqrt5 \) (viz obrázek).}
{\([2;3] \), \([6;5] \), \([8;1] \), \([4;-1] \)}{\([2;3] \), \([6;5] \), \([8{,}5;1] \), \([4{,}5;-1] \)}{\([2;3{,}5] \), \([6;5{,}5] \), \([8;1] \), \([4;-1] \)}{\([2;3] \), \([6;5{,}5] \), \([8;1{,}5] \), \([4;-1] \)}{}{}}
\LANGsk{
\Question{
Sú dané priamky \( p \): \( x-2y-1=0 \) a \( q \): \( 2x+y-12=0 \). Určte súradnice všetkých takých bodov, ktoré majú od oboch daných priamok vzdialenosť \( \sqrt5 \) (viď obrázok).}
{\( [2;3] \), \( [6;5] \), \( [8;1] \), \( [4;-1] \)}{\( [2;3] \), \( [6;5] \), \( [8{,}5;1] \), \( [4{,}5;-1] \)}{\( [2;3{,}5] \), \( [6;5{,}5] \), \( [8;1] \), \( [4;-1] \)}{\( [2;3] \), \( [6;5{,}5] \), \( [8;1{,}5] \), \( [4;-1] \)}{}{}}
\LANGes{
\Question{
Sean \( p \) y \( q \)  rectas con ecuaciones \( x-2y-1=0 \) y \( 2x+y-12=0 \) respectivamente. Halla todos los puntos que están a la misma distancia \( \sqrt5 \) de \( p \) y de \( q \)  (mira la imagen).}
{\([2;3] \), \([6;5] \), \([8;1] \), \([4;-1] \)}{\([2;3] \), \([6;5] \), \([8.5;1] \), \([4.5;-1] \)}{\([2;3.5] \), \([6;5.5] \), \([8;1] \), \([4;-1] \)}{\([2;3] \), \([6;5.5] \), \([8;1.5] \), \([4;-1] \)}{}{}}
\End

\Data\ID{1103109106}
\Updated{2021-12-05 07:12:52}
\Level{C}
\SubArea{Analytic geometry in a plane}
\IMG{1103109106.pdf}
\LANGen{
\Question{
Find general form equations of all the lines passing through the point \( M=[-2;4] \) at the distance of \( 2 \) from the origin \( O \) (see the picture).}
{\( x+2=0;\ 3x+4y-10=0 \)}{\( x-2=0;\ 3x+4y-10=0 \)}{\( x+2=0;\ 4x-3y+20=0 \)}{\( x-2=0;\ 4x-3y+20=0 \)}{}{}}
\LANGpl{
\Question{
Wskaż równania wszystkich prostych przechodzących przez punkt \( M=[-2;4] \) w odległości \( 2 \) od początku układu współrzędnych \( O \) (spójrz na rysunek).}
{\( x+2=0;\ 3x+4y-10=0 \)}{\( x-2=0;\ 3x+4y-10=0 \)}{\( x+2=0;\ 4x-3y+20=0 \)}{\( x-2=0;\ 4x-3y+20=0 \)}{}{}}
\LANGcs{
\Question{
Určete obecné rovnice všech přímek, které procházejí bodem \( M=[-2;4] \) a mají od počátku souřadné soustavy \( O \) vzdálenost \( 2 \) (viz obrázek).}
{\( x+2=0;\ 3x+4y-10=0 \)}{\( x-2=0;\ 3x+4y-10=0 \)}{\( x+2=0;\ 4x-3y+20=0 \)}{\( x-2=0;\ 4x-3y+20=0 \)}{}{}}
\LANGsk{
\Question{
Určte všeobecné rovnice všetkých priamok, ktoré prechádzajú bodom \( M=[-2;4] \) a majú od začiatku súradnicovej sústavy \( O \) vzdialenosť \( 2 \) (viď obrázok).}
{\( x+2=0;\ 3x+4y-10=0 \)}{\( x-2=0;\ 3x+4y-10=0 \)}{\( x+2=0;\ 4x-3y+20=0 \)}{\( x-2=0;\ 4x-3y+20=0 \)}{}{}}
\LANGes{
\Question{
Halla las ecuaciones generales de todas las rectas que pasan por el punto \( M=[-2;4] \) y están a una  distancia \( 2 \)  del origen \( O \) (mira la imagen).}
{\( x+2=0;\ 3x+4y-10=0 \)}{\( x-2=0;\ 3x+4y-10=0 \)}{\( x+2=0;\ 4x-3y+20=0 \)}{\( x-2=0;\ 4x-3y+20=0 \)}{}{}}
\End

\Data\ID{1103109107}
\Updated{2022-06-24 08:06:51}
\Level{C}
\SubArea{Analytic geometry in a plane}
\IMG{1103109107.pdf}
\LANGen{
\Question{
Let \( ABC \) be a triangle (see the picture). Determine the angle \( \varphi \) between the height \( v_c \) and the median \( t_c \). Give the angle rounded to minutes.}
{\( \varphi\doteq 21^{\circ}48' \)}{\( \varphi\doteq 21^{\circ}24' \)}{\( \varphi\doteq 21^{\circ}36' \)}{\( \varphi\doteq 21^{\circ}52' \)}{}{}}
\LANGpl{
\Question{
Dany jest trójkąt  \( ABC \)  (spójrz na rysunek). Wyznacz kąt  \( \varphi \) pomiędzy wysokością  \( v_c \) a środkową \( t_c \). Zaokrągli miarę kąta do minut.}
{\( \varphi\doteq 21^{\circ}48' \)}{\( \varphi\doteq 21^{\circ}24' \)}{\( \varphi\doteq 21^{\circ}36' \)}{\( \varphi\doteq 21^{\circ}52' \)}{}{}}
\LANGcs{
\Question{
Je dán trojúhelník \( ABC \) (viz obrázek). Určete odchylku \( \varphi \) jeho  výšky \( v_c \) a těžnice \( t_c \). Odchylku zaokrouhlete na minuty.}
{\( \varphi\doteq 21^{\circ}48' \)}{\( \varphi\doteq 21^{\circ}24' \)}{\( \varphi\doteq 21^{\circ}36' \)}{\( \varphi\doteq 21^{\circ}52' \)}{}{}}
\LANGsk{
\Question{
Je daný trojuholník \( ABC \) (viď obrázok). Určte odchýlku \( \varphi \) jeho  výšky \( v_c \) a ťažnice \( t_c \). Odchýlku zaokrúhlite na minúty.}
{\( \varphi\doteq 21^{\circ}48' \)}{\( \varphi\doteq 21^{\circ}24' \)}{\( \varphi\doteq 21^{\circ}36' \)}{\( \varphi\doteq 21^{\circ}52' \)}{}{}}
\LANGes{
\Question{
Sea \( ABC \) un triángulo  (mira la imagen). Determina el ángulo \( \varphi \) entre la altura \( v_c \) y la mediana \( t_c \). Redondea el ángulo a minutos.}
{\( \varphi\doteq 21^{\circ}48' \)}{\( \varphi\doteq 21^{\circ}24' \)}{\( \varphi\doteq 21^{\circ}36' \)}{\( \varphi\doteq 21^{\circ}52' \)}{}{}}
\End

\Data\ID{1103109108}
\Updated{2022-06-24 08:06:51}
\Level{C}
\SubArea{Analytic geometry in a plane}
\IMG{1103109108.pdf}
\LANGen{
\Question{
Let \( ABC \) be a triangle (see the picture). Determine the angle \( \varphi \) between the height \( v_b \) and the angle bisector \( o_\alpha \). Give the angle rounded to minutes.}
{\( \varphi\doteq 71^{\circ}34' \)}{\( \varphi\doteq 71^{\circ}33' \)}{\( \varphi\doteq 71^{\circ}40' \)}{\( \varphi\doteq 71^{\circ}38' \)}{}{}}
\LANGpl{
\Question{
Dany jest trójkąt  \( ABC \)  (spójrz na rysunek). Określ kąt \( \varphi \) pomiędzy wysokością  \( v_b \), a dwusieczną kąta \( o_\alpha \). Zaokrągli miarę kąta do minut.}
{\( \varphi\doteq 71^{\circ}34' \)}{\( \varphi\doteq 71^{\circ}33' \)}{\( \varphi\doteq 71^{\circ}40' \)}{\( \varphi\doteq 71^{\circ}38' \)}{}{}}
\LANGcs{
\Question{
Je dán trojúhelník \( ABC \) (viz obrázek). Určete odchylku \( \varphi \) jeho výšky \( v_b \) a osy úhlu \( o_\alpha \). Odchylku zaokrouhlete na minuty.}
{\( \varphi\doteq 71^{\circ}34' \)}{\( \varphi\doteq 71^{\circ}33' \)}{\( \varphi\doteq 71^{\circ}40' \)}{\( \varphi\doteq 71^{\circ}38' \)}{}{}}
\LANGsk{
\Question{
Je daný trojuholník \( ABC \) (viď obrázok). Určte odchýlku \( \varphi \) jeho výšky \( v_b \) a osy uhlov \( o_\alpha \). Odchýlku zaokrúhlite na minúty.}
{\( \varphi\doteq 71^{\circ}34' \)}{\( \varphi\doteq 71^{\circ}33' \)}{\( \varphi\doteq 71^{\circ}40' \)}{\( \varphi\doteq 71^{\circ}38' \)}{}{}}
\LANGes{
\Question{
Sea \( ABC \) un triángulo  (mira la imagen). Determina el ángulo \( \varphi \) entre la altura \( v_b \) y la bisectriz de un ángulo \( o_\alpha \). Redondea el ángulo a minutos.}
{\( \varphi\doteq 71^{\circ}34' \)}{\( \varphi\doteq 71^{\circ}33' \)}{\( \varphi\doteq 71^{\circ}40' \)}{\( \varphi\doteq 71^{\circ}38' \)}{}{}}
\End

\Data\ID{2010014207}
\Updated{2022-05-16 11:05:40}
\Level{C}
\SubArea{Analytic geometry in a plane}
\LANGen{
\Question{
Given two points \(A = [2;1]\) 
and \(B = [4;-2]\), identify a
number \(m\in \mathbb{R}\) such
that the point \(C = [1;m]\) 
is on the line \(AB\).}
{\( m=\frac52 \)}{\( m=-\frac12 \)}{\( m=2 \)}{\( m=\frac13 \)}{}{}}
\LANGpl{
\Question{
Rozważając dwa punkty \(A = [2;1]\)
i \(B = [4;-2]\). Wyznacz wartość 
\(m\in \mathbb{R}\) taką, że punkt
\(C = [1;m]\)  znajduje się na prostej \(AB\).}
{\( m=\frac52 \)}{\( m=-\frac12 \)}{\( m=2 \)}{\( m=\frac13 \)}{}{}}
\LANGcs{
\Question{
Jsou dány body \(A = [2;1]\)  a \(B = [4;-2]\). Určete \(m\in \mathbb{R}\) tak,
aby bod \(C = [1;m]\) 
ležel na přímce \(AB\).}
{\( m=\frac52 \)}{\( m=-\frac12 \)}{\( m=2 \)}{\( m=\frac13 \)}{}{}}
\LANGsk{
\Question{
Dané sú body \(A = [2;1]\) 
a \(B = [4;-2]\), určte
\(m\in \mathbb{R}\) tak, aby
bod \(C = [1;m]\) 
patril priamke \(AB\).}
{\( m=\frac52 \)}{\( m=-\frac12 \)}{\( m=2 \)}{\( m=\frac13 \)}{}{}}
\LANGes{
\Question{
Dados dos puntos \(A = [2;1]\) 
y \(B = [4;-2]\).  Identifica un número \(m\in \mathbb{R}\) suponiendo que el punto \(C = [1;m]\) 
está en la recta \(AB\).}
{\( m=\frac52 \)}{\( m=-\frac12 \)}{\( m=2 \)}{\( m=\frac13 \)}{}{}}
\End

\Data\ID{2010014208}
\Updated{2022-05-16 11:05:40}
\Level{C}
\SubArea{Analytic geometry in a plane}
\LANGen{
\Question{
Given lines \(p\) 
and \(q\), find
\(m\in \mathbb{R}\) such that
the lines \(p\) 
and \(q\) 
are parallel. 
\[ 
 p\colon x +3y + 4= 0,\qquad q\colon mx -2y - 7= 0
\]}
{\( m=-\frac23 \)}{\( m=6 \)}{\( m=-\frac13 \)}{\( m=\frac13 \)}{}{}}
\LANGpl{
\Question{
Dane są proste \(p\)
i \(q\). Znajdź wartość 
\(m\in \mathbb{R}\) taką, że proste 
\(p\) i \(q\)
są równoległe.
\[ 
 p\colon x +3y + 4= 0,\qquad q\colon mx -2y - 7= 0
\]}
{\( m=-\frac23 \)}{\( m=6 \)}{\( m=-\frac13 \)}{\( m=\frac13 \)}{}{}}
\LANGcs{
\Question{
Určete \(m\in \mathbb{R}\) tak, aby byly přímky \(p\) a \(q\) rovnoběžné. 
\[ 
 p\colon x +3y + 4= 0,\qquad q\colon mx -2y - 7= 0
\]}
{\( m=-\frac23 \)}{\( m=6 \)}{\( m=-\frac13 \)}{\( m=\frac13 \)}{}{}}
\LANGsk{
\Question{
Dané sú priamky \(p\) 
a \(q\). Určte
\(m\in \mathbb{R}\) tak, aby 
priamky \(p\) 
a \(q\) 
boli rovnobežné.
\[ 
 p\colon x +3y + 4= 0,\qquad q\colon mx -2y - 7= 0
\]}
{\( m=-\frac23 \)}{\( m=6 \)}{\( m=-\frac13 \)}{\( m=\frac13 \)}{}{}}
\LANGes{
\Question{
Dadas las rectas \(p\) 
y \(q\), determina  el punto
\(m\in \mathbb{R}\) suponiendo que las rectas \(p\) 
y \(q\) 
son paralelas. 
\[ 
 p\colon x +3y + 4= 0,\qquad q\colon mx -2y - 7= 0
\]}
{\( m=-\frac23 \)}{\( m=6 \)}{\( m=-\frac13 \)}{\( m=\frac13 \)}{}{}}
\End

\Data\ID{9000090901}
\Updated{2022-04-23 05:04:00}
\Level{C}
\SubArea{Analytic geometry in a plane}
\LANGen{
\Question{
Given two points \(A = [2;5]\) 
and \(B = [-3;2]\), identify a
number \(m\in \mathbb{R}\) such
that the point \(C = [1;m]\) 
is on the line \(AB\).}
{\(m = \frac{22} 
 {5} \)}{\(m = 20\)}{\(m = -3\)}{\(m = \frac{2} 
{3}\)}{\(m = -\frac{5}
{2}\)}{}}
\LANGpl{
\Question{
Dane są punkty \(A = [2;5]\) 
i \(B = [-3;2]\), określ
liczbę  \(m\in \mathbb{R}\) tak, 
aby punkt \(C = [1;m]\) 
leżał na prostej \(AB\).}
{\(m = \frac{22} 
 {5} \)}{\(m = 20\)}{\(m = -3\)}{\(m = \frac{2} 
{3}\)}{\(m = -\frac{5}
{2}\)}{}}
\LANGcs{
\Question{
Určete \(m\in \mathbb{R}\) tak,
aby bod \(C = [1;m]\) ležel
na přímce \(AB\),
kde \(A = [2;5]\),
\(B = [-3;2]\).}
{\(m = \frac{22} 
 {5} \)}{\(m = 20\)}{\(m = -3\)}{\(m = \frac{2} 
{3}\)}{\(m = -\frac{5}
{2}\)}{}}
\LANGsk{
\Question{
Určte \(m\in \mathbb{R}\) tak,
aby bod \(C = [1;m]\) ležal
na priamke \(AB\),
kde \(A = [2;5]\),
\(B = [-3;2]\).}
{\(m = \frac{22} 
 {5} \)}{\(m = 20\)}{\(m = -3\)}{\(m = \frac{2} 
{3}\)}{\(m = -\frac{5}
{2}\)}{}}
\LANGes{
\Question{
Dados dos puntos \(A = [2;5]\) 
y \(B = [-3;2]\). Identifica un número \(m\in \mathbb{R}\) suponiendo que el punto \(C = [1;m]\) 
está en la recta\(AB\).}
{\(m = \frac{22} 
 {5} \)}{\(m = 20\)}{\(m = -3\)}{\(m = \frac{2} 
{3}\)}{\(m = -\frac{5}
{2}\)}{}}
\End

\Data\ID{9000090902}
\Updated{2021-12-05 07:12:08}
\Level{C}
\SubArea{Analytic geometry in a plane}
\LANGen{
\Question{
Given the parametric line \(p\),
find \(m\in \mathbb{R}\) such that
the point \(C = [m;3]\) 
is on the line \(p\).
\[ 
\begin{aligned}p\colon x& = 1 - t, &
\\y & = -3 + 2t;\  t\in \mathbb{R}
\\ \end{aligned}
\]}
{\(m = -2\)}{\(m = 4\)}{\(m = 11\)}{\(m = -\frac{11}
 {3} \)}{\(m = \frac{3} 
{2}\)}{}}
\LANGpl{
\Question{
Dana jest prosta \(p\),
wyznacz \(m\in \mathbb{R}\) tak, aby
punkt \(C = [m;3]\) 
leżał na prostej \(p\).
\[ 
\begin{aligned}p\colon x& = 1 - t, &
\\y & = -3 + 2t;\  t\in \mathbb{R}
\\ \end{aligned}
\]}
{\(m = -2\)}{\(m = 4\)}{\(m = 11\)}{\(m = -\frac{11}
 {3} \)}{\(m = \frac{3} 
{2}\)}{}}
\LANGcs{
\Question{
Určete \(m\in \mathbb{R}\) tak, aby
bod \(C = [m;3]\) ležel na
přímce \(p\). \begin{align*} p\colon x &= 1 - t, \\
y &= -3 + 2t;\ t\in \mathbb{R} \end{align*}}
{\(m = -2\)}{\(m = 4\)}{\(m = 11\)}{\(m = -\frac{11}
 {3} \)}{\(m = \frac{3} 
{2}\)}{}}
\LANGsk{
\Question{
Určte \(m\in \mathbb{R}\) tak, aby
bod \(C = [m;3]\) ležal na
priamke \(p\).\begin{align*} p\colon x &= 1 - t, \\
y &= -3 + 2t;\ t\in \mathbb{R} \end{align*}}
{\(m = -2\)}{\(m = 4\)}{\(m = 11\)}{\(m = -\frac{11}
 {3} \)}{\(m = \frac{3} 
{2}\)}{}}
\LANGes{
\Question{
Dada la recta paramétrica \(p\).
Halla \(m\in \mathbb{R}\) suponiendo que el punto \(C = [m;3]\) 
está en la recta \(p\).
\[ 
\begin{aligned}p\colon x& = 1 - t, &
\\y & = -3 + 2t;\  t\in \mathbb{R}
\\ \end{aligned}
\]}
{\(m = -2\)}{\(m = 4\)}{\(m = 11\)}{\(m = -\frac{11}
 {3} \)}{\(m = \frac{3} 
{2}\)}{}}
\End

\Data\ID{9000090903}
\Updated{2021-12-05 07:12:19}
\Level{C}
\SubArea{Analytic geometry in a plane}
\LANGen{
\Question{
Given the line 
\[ 
 p\colon 3x - 2y + 11 = 0,
\]
find \(m\in \mathbb{R}\) such that
the point \(C = [m;0]\) 
is on the line \(p\).}
{\(m = -\frac{11}
 {3} \)}{\(m = -1\)}{\(m = 11\)}{\(m = -\frac{1}
{11}\)}{\(m = 2\)}{}}
\LANGpl{
\Question{
Dana jest prosta  
\[ 
 p\colon 3x - 2y + 11 = 0,
\]
określ \(m\in \mathbb{R}\) tak, aby
punkt  \(C = [m;0]\) 
leżał na prostej  \(p\).}
{\(m = -\frac{11}
 {3} \)}{\(m = -1\)}{\(m = 11\)}{\(m = -\frac{1}
{11}\)}{\(m = 2\)}{}}
\LANGcs{
\Question{
Určete \(m\in \mathbb{R}\) tak, aby
bod \(C = [m;0]\) ležel na
přímce \(p\). \[ p\colon 3x - 2y + 11 = 0\]}
{\(m = -\frac{11}
 {3} \)}{\(m = -1\)}{\(m = 11\)}{\(m = -\frac{1}
{11}\)}{\(m = 2\)}{}}
\LANGsk{
\Question{
Určte \(m\in \mathbb{R}\) tak, aby
bod \(C = [m;0]\) ležal na
priamke \(p\). \[ p\colon 3x - 2y + 11 = 0\]}
{\(m = -\frac{11}
 {3} \)}{\(m = -1\)}{\(m = 11\)}{\(m = -\frac{1}
{11}\)}{\(m = 2\)}{}}
\LANGes{
\Question{
Dada la recta
\[ 
 p\colon 3x - 2y + 11 = 0,
\]
Determina el punto \(m\in \mathbb{R}\) suponiendo que el punto \(C = [m;0]\) 
está en la recta  \(p\).}
{\(m = -\frac{11}
 {3} \)}{\(m = -1\)}{\(m = 11\)}{\(m = -\frac{1}
{11}\)}{\(m = 2\)}{}}
\End

\Data\ID{9000090904}
\Updated{2022-04-23 05:04:00}
\Level{C}
\SubArea{Analytic geometry in a plane}
\LANGen{
\Question{
Given lines \(p\) 
and \(q\), find
\(m\in \mathbb{R}\) such that
the lines \(p\) 
and \(q\) 
are parallel. 
\[ 
 p\colon x - 2y + 7 = 0,\qquad q\colon mx + 3y - 11 = 0
\]}
{\(m = -\frac{3}
{2}\)}{\(m = \frac{2} 
{3}\)}{\(m = \frac{3} 
{2}\)}{\(m = -\frac{2}
{3}\)}{another solution}{}}
\LANGpl{
\Question{
Dane są proste \(p\) 
i \(q\), określ
\(m\in \mathbb{R}\) tak, aby
proste \(p\) 
i \(q\) 
były równoległe.
\[ 
 p\colon x - 2y + 7 = 0,\qquad q\colon mx + 3y - 11 = 0
\]}
{\(m = -\frac{3}
{2}\)}{\(m = \frac{2} 
{3}\)}{\(m = \frac{3} 
{2}\)}{\(m = -\frac{2}
{3}\)}{inne rozwiązanie}{}}
\LANGcs{
\Question{
Určete \(m\in \mathbb{R}\) tak, aby přímka
\(p\colon x - 2y + 7 = 0\) byla rovnoběžná
s přímkou \(q\colon mx + 3y - 11 = 0\).}
{\(m = -\frac{3}
{2}\)}{\(m = \frac{2} 
{3}\)}{\(m = \frac{3} 
{2}\)}{\(m = -\frac{2}
{3}\)}{jiná možnost}{}}
\LANGsk{
\Question{
Určte \(m\in \mathbb{R}\) tak, aby priamka
\(p\colon x - 2y + 7 = 0\) bola rovnobežná
s priamkou \(q\colon mx + 3y - 11 = 0\).}
{\(m = -\frac{3}
{2}\)}{\(m = \frac{2} 
{3}\)}{\(m = \frac{3} 
{2}\)}{\(m = -\frac{2}
{3}\)}{iná možnosť}{}}
\LANGes{
\Question{
Dadas las rectas \(p\) 
y \(q\). Determina el punto
\(m\in \mathbb{R}\) suponiendo que las rectas \(p\) 
y \(q\) 
son paralelas. 
\[ 
 p\colon x - 2y + 7 = 0,\qquad q\colon mx + 3y - 11 = 0
\]}
{\(m = -\frac{3}
{2}\)}{\(m = \frac{2} 
{3}\)}{\(m = \frac{3} 
{2}\)}{\(m = -\frac{2}
{3}\)}{another solution}{}}
\End

\Data\ID{9000090905}
\Updated{2020-07-15 03:07:37}
\Level{C}
\SubArea{Analytic geometry in a plane}
\LANGen{
\Question{
Given lines \(p\) 
and \(q\), find
\(m\in \mathbb{R}\) such that
the lines \(p\) 
and \(q\) 
are parallel. 
\[ 
 p\colon x - 2y + 7 = 0,\qquad q\colon x + 3y + m = 0
\]}
{does not exist}{\(m = -7\)}{\(m = -\frac{2}
{7}\)}{\(m = 2\)}{\(m = 7\)}{}}
\LANGpl{
\Question{
Dana jest proste \(p\) 
i \(q\),określ
\(m\in \mathbb{R}\) tak, aby
proste \(p\) 
i \(q\) 
były równoległe. 
\[ 
 p\colon x - 2y + 7 = 0,\qquad q\colon x + 3y + m = 0
\]}
{nie istnieje}{\(m = -7\)}{\(m = -\frac{2}
{7}\)}{\(m = 2\)}{\(m = 7\)}{}}
\LANGcs{
\Question{
Určete \(m\in \mathbb{R}\) tak, aby přímka
\(p\colon x - 2y + 7 = 0\) byla rovnoběžná
s přímkou \(q\colon x + 3y + m = 0\).}
{takové \(m\) neexistuje}{\(m = -7\)}{\(m = -\frac{2}
{7}\)}{\(m = 2\)}{\(m = 7\)}{}}
\LANGsk{
\Question{
Určte \(m\in \mathbb{R}\) tak, aby priamka
\(p\colon x - 2y + 7 = 0\) bola rovnobežná
s priamkou \(q\colon x + 3y + m = 0\).}
{také \(m\) 
neexistuje}{\(m = -7\)}{\(m = -\frac{2}
{7}\)}{\(m = 2\)}{\(m = 7\)}{}}
\LANGes{
\Question{
Dadas las rectas \(p\) 
y \(q\). Determina el punto
\(m\in \mathbb{R}\) suponiendo que las rectas \(p\) 
y \(q\) 
son paralelas. 
\[ 
 p\colon x - 2y + 7 = 0,\qquad q\colon x + 3y + m = 0
\]}
{no existe}{\(m = -7\)}{\(m = -\frac{2}
{7}\)}{\(m = 2\)}{\(m = 7\)}{}}
\End

\Data\ID{9000090906}
\Updated{2020-07-15 03:07:37}
\Level{C}
\SubArea{Analytic geometry in a plane}
\LANGen{
\Question{
Given lines \(p\) 
and \(q\), find
\(m\in \mathbb{R}\) such that
the lines \(p\) 
and \(q\) 
are parallel. 
\[ 
\begin{aligned}p\colon x& = 1 + t, &
\\y & = -3t;\  t\in \mathbb{R},
\\ \end{aligned}\qquad \begin{aligned}q\colon x& = 3 - 2u, &
\\y & = 1 + mu;\ u\in \mathbb{R}
\\ \end{aligned}
\]}
{\(m = 6\)}{\(m = \frac{3} 
{2}\)}{\(m = -\frac{2}
{3}\)}{does not exist}{}{}}
\LANGpl{
\Question{
Dane są proste \(p\) 
i \(q\), określ
\(m\in \mathbb{R}\) tak, aby
proste \(p\) 
i \(q\) 
były równoległe. 

\[ 
\begin{aligned}p\colon x& = 1 + t, &
\\y & = -3t;\  t\in \mathbb{R},
\\ \end{aligned}\qquad \begin{aligned}q\colon x& = 3 - 2u, &
\\y & = 1 + mu;\ u\in \mathbb{R}
\\ \end{aligned}
\]}
{\(m = 6\)}{\(m = \frac{3} 
{2}\)}{\(m = -\frac{2}
{3}\)}{nie istnieje}{}{}}
\LANGcs{
\Question{
Určete \(m\in \mathbb{R}\) tak,
aby přímka \(p\)
\[ 
 p\colon x = 1 + t,\ y = -3t,\ t\in \mathbb{R}
\]
byla rovnoběžná s přímkou \(q\)
\[ 
 q\colon x = 3 - 2u,\ y = 1 + mu,\ u\in \mathbb{R}.
\]}
{\(m = 6\)}{\(m = \frac{3} 
{2}\)}{\(m = -\frac{2}
{3}\)}{takové \(m\) 
neexistuje}{}{}}
\LANGsk{
\Question{
Určte \(m\in \mathbb{R}\) tak,
aby priamka \(p\)
\[ 
 p\colon x = 1 + t,\ y = -3t,\ t\in \mathbb{R}
\]
bola rovnobežná s priamkou \(q\)
\[ 
 q\colon x = 3 - 2u,\ y = 1 + mu,\ u\in \mathbb{R}.
\]}
{\(m = 6\)}{\(m = \frac{3} 
{2}\)}{\(m = -\frac{2}
{3}\)}{také \(m\) 
neexistuje}{}{}}
\LANGes{
\Question{
Dadas las rectas \(p\) 
y \(q\). Determina el punto 
\(m\in \mathbb{R}\) suponiendo que las rectas \(p\) 
y \(q\) 
son paralelas. 
\[ 
\begin{aligned}p\colon x& = 1 + t, &
\\y & = -3t;\  t\in \mathbb{R},
\\ \end{aligned}\qquad \begin{aligned}q\colon x& = 3 - 2u, &
\\y & = 1 + mu;\ u\in \mathbb{R}
\\ \end{aligned}
\]}
{\(m = 6\)}{\(m = \frac{3} 
{2}\)}{\(m = -\frac{2}
{3}\)}{no existe}{}{}}
\End

\Data\ID{9000090907}
\Updated{2020-07-15 03:07:37}
\Level{C}
\SubArea{Analytic geometry in a plane}
\LANGen{
\Question{
Given points \(A = [2;m]\) 
and \(B = [-1;0]\), find
\(m\in \mathbb{R}\) such that the line
\(p\) is parallel to the line passes
through the points \(A\), \(B\).
\[ 
\begin{aligned}p\colon x& = 3 + 2t, &
\\y & = 5 - t;\ t\in \mathbb{R}
\\ \end{aligned}
\]}
{\(m = -\frac{3}
{2}\)}{\(m = \frac{3} 
{2}\)}{\(m = -\frac{2}
{3}\)}{\(m = 2\)}{does not exist}{}}
\LANGpl{
\Question{
Dane są punkty \(A = [2;m]\) 
i \(B = [-1;0]\), określ
\(m\in \mathbb{R}\) tak, aby prosta
\(p\) była równoległa do prostej 
 \(AB\).
\[ 
\begin{aligned}p\colon x& = 3 + 2t, &
\\y & = 5 - t;\ t\in \mathbb{R}
\\ \end{aligned}
\]}
{\(m = -\frac{3}
{2}\)}{\(m = \frac{3} 
{2}\)}{\(m = -\frac{2}
{3}\)}{\(m = 2\)}{brak rozwiązania}{}}
\LANGcs{
\Question{
Určete \(m\in \mathbb{R}\) tak, aby přímka
\(p\colon x = 3 + 2t,\ y = 5 - t,\ t\in \mathbb{R}\) byla rovnoběžná
s přímkou \(AB\),
kde \(A = [2;m]\),
\(B = [-1;0]\).}
{\(m = -\frac{3}
{2}\)}{\(m = \frac{3} 
{2}\)}{\(m = -\frac{2}
{3}\)}{\(m = 2\)}{takové \(m\) 
neexistuje}{}}
\LANGsk{
\Question{
Určte \(m\in \mathbb{R}\) tak, aby priamka
\(p\colon x = 3 + 2t,\ y = 5 - t,\ t\in \mathbb{R}\) bola rovnobežná
s priamkou \(AB\),
kde \(A = [2;m]\),
\(B = [-1;0]\).}
{\(m = -\frac{3}
{2}\)}{\(m = \frac{3} 
{2}\)}{\(m = -\frac{2}
{3}\)}{\(m = 2\)}{také \(m\) 
neexistuje}{}}
\LANGes{
\Question{
Dados los puntos \(A = [2;m]\) 
y \(B = [-1;0]\). Determina el punto 
\(m\in \mathbb{R}\) suponiendo que la recta
\(p\) es paralela a la recta que pasa por los puntos \(A\), \(B\).
\[ 
\begin{aligned}p\colon x& = 3 + 2t, &
\\y & = 5 - t;\ t\in \mathbb{R}
\\ \end{aligned}
\]}
{\(m = -\frac{3}
{2}\)}{\(m = \frac{3} 
{2}\)}{\(m = -\frac{2}
{3}\)}{\(m = 2\)}{no existe}{}}
\End

\Data\ID{9000090908}
\Updated{2020-07-15 03:07:37}
\Level{C}
\SubArea{Analytic geometry in a plane}
\LANGen{
\Question{
Given points \(A = [2;1]\) 
and \(B = [m;0]\), find
\(m\in \mathbb{R}\) such that the line
\(p\) is parallel to the line passes
through the points \(A\), \(B\).
\[ 
 p\colon 3x - y + 17 = 0
\]}
{\(m = \frac{5} 
{3}\)}{\(m = 4\)}{\(m = \frac{5} 
{2}\)}{\(m = -1\)}{another solution}{}}
\LANGpl{
\Question{
Dane są punkty \(A = [2;1]\) 
i \(B = [m;0]\), określ
\(m\in \mathbb{R}\) tak, aby prosta
\(p\) była równoległa do prostej
 \(AB\).
\[ 
 p\colon 3x - y + 17 = 0
\]}
{\(m = \frac{5} 
{3}\)}{\(m = 4\)}{\(m = \frac{5} 
{2}\)}{\(m = -1\)}{inne rozwiązanie}{}}
\LANGcs{
\Question{
Určete \(m\in \mathbb{R}\) tak, aby přímka
\(p\colon 3x - y + 17 = 0\) byla rovnoběžná
s přímkou \(AB\),
kde \(A = [2;1]\),
\(B = [m;0]\).}
{\(m = \frac{5} 
{3}\)}{\(m = 4\)}{\(m = \frac{5} 
{2}\)}{\(m = -1\)}{jiná možnost}{}}
\LANGsk{
\Question{
Určte \(m\in \mathbb{R}\) tak, aby priamka
\(p\colon 3x - y + 17 = 0\) bola rovnobežná
s priamkou \(AB\),
kde \(A = [2;1]\),
\(B = [m;0]\).}
{\(m = \frac{5} 
{3}\)}{\(m = 4\)}{\(m = \frac{5} 
{2}\)}{\(m = -1\)}{iná možnosť}{}}
\LANGes{
\Question{
Dados los puntos  \(A = [2;1]\) 
y \(B = [m;0]\). Determina el punto 
\(m\in \mathbb{R}\) suponiendo que la recta
\(p\) es paralela a la recta que pasa por los puntos \(A\), \(B\).
\[ 
 p\colon 3x - y + 17 = 0
\]}
{\(m = \frac{5} 
{3}\)}{\(m = 4\)}{\(m = \frac{5} 
{2}\)}{\(m = -1\)}{otra solución}{}}
\End

\Data\ID{9000090909}
\Updated{2020-07-15 03:07:37}
\Level{C}
\SubArea{Analytic geometry in a plane}
\LANGen{
\Question{
Given lines \(p\) 
and \(q\), find
\(m\in \mathbb{R}\) such that
the lines \(p\) 
and \(q\) 
are parallel. 
\[ 
p\colon 2x+my-3 = 0,\qquad \begin{aligned}[t] q\colon x& = 1 + t, &
\\y & = 2 - t;\ t\in \mathbb{R}
\\ \end{aligned}
\]}
{\(m = 2\)}{\(m = -2\)}{\(m = 11\)}{\(m = -\frac{1}
{11}\)}{does not exist}{}}
\LANGpl{
\Question{
Dane są proste \(p\) 
i \(q\), określ
\(m\in \mathbb{R}\) tak, aby
proste \(p\) 
i \(q\) 
były równoległe
\[ 
p\colon 2x+my-3 = 0,\qquad \begin{aligned}[t] q\colon x& = 1 + t, &
\\y & = 2 - t;\ t\in \mathbb{R}
\\ \end{aligned}
\]}
{\(m = 2\)}{\(m = -2\)}{\(m = 11\)}{\(m = -\frac{1}
{11}\)}{nie istnieje}{}}
\LANGcs{
\Question{
Určete \(m\in \mathbb{R}\) tak, aby přímka
\(p\colon x = 1 + t,\ y = 2 - t,\ t\in \mathbb{R}\) byla rovnoběžná
s přímkou \(q\colon 2x + my - 3 = 0\).}
{\(m = 2\)}{\(m = -2\)}{\(m = 11\)}{\(m = -\frac{1}
{11}\)}{takové \(m\) 
neexistuje}{}}
\LANGsk{
\Question{
Určte \(m\in \mathbb{R}\) tak, aby priamka
\(p\colon x = 1 + t,\ y = 2 - t,\ t\in \mathbb{R}\) bola rovnobežná
s priamkou \(q\colon 2x + my - 3 = 0\).}
{\(m = 2\)}{\(m = -2\)}{\(m = 11\)}{\(m = -\frac{1}
{11}\)}{také \(m\) 
neexistuje}{}}
\LANGes{
\Question{
Dadas las rectas \(p\) 
y \(q\). Determina el punto 
\(m\in \mathbb{R}\) suponiendo que las rectas \(p\) 
y \(q\) 
son paralelas. 
\[ 
p\colon 2x+my-3 = 0,\qquad \begin{aligned}[t] q\colon x& = 1 + t, &
\\y & = 2 - t;\ t\in \mathbb{R}
\\ \end{aligned}
\]}
{\(m = 2\)}{\(m = -2\)}{\(m = 11\)}{\(m = -\frac{1}
{11}\)}{no existe}{}}
\End

\Data\ID{9000090910}
\Updated{2020-07-15 03:07:37}
\Level{C}
\SubArea{Analytic geometry in a plane}
\LANGen{
\Question{
Given lines \(p\) 
and \(q\), find
\(m\in \mathbb{R}\) such that
the line \(p\) is
parallel to \(q\).
\[ 
p\colon x+4y-3 = 0,\qquad \begin{aligned}[t] q\colon x& = 1 + mt,&
\\y & = 2 - 3t;\ t\in \mathbb{R}
\\ \end{aligned}
\]}
{\(m = 12\)}{\(m = -\frac{1}
{12}\)}{\(m = 4\)}{\(m = \frac{5} 
{2}\)}{\(m = -1\)}{}}
\LANGpl{
\Question{
Dane są proste \(p\) 
i \(q\), określ
\(m\in \mathbb{R}\) tak, aby
prosta \(p\) była
równoległa do prostej \(q\).
\[ 
p\colon x+4y-3 = 0,\qquad \begin{aligned}[t] q\colon x& = 1 + mt,&
\\y & = 2 - 3t;\ t\in \mathbb{R}
\\ \end{aligned}
\]}
{\(m = 12\)}{\(m = -\frac{1}
{12}\)}{\(m = 4\)}{\(m = \frac{5} 
{2}\)}{\(m = -1\)}{}}
\LANGcs{
\Question{
Určete \(m\in \mathbb{R}\) tak, aby přímka
\(p\colon x = 1 + mt,\ y = 2 - 3t,\ t\in \mathbb{R}\) byla rovnoběžná
s přímkou \(q\colon x + 4y - 3 = 0\).}
{\(m = 12\)}{\(m = -\frac{1}
{12}\)}{\(m = 4\)}{\(m = \frac{5} 
{2}\)}{\(m = -1\)}{}}
\LANGsk{
\Question{
Určte \(m\in \mathbb{R}\) tak, aby priamka
\(p\colon x = 1 + mt,\ y = 2 - 3t,\ t\in \mathbb{R}\) bola rovnobežná
s priamkou \(q\colon x + 4y - 3 = 0\).}
{\(m = 12\)}{\(m = -\frac{1}
{12}\)}{\(m = 4\)}{\(m = \frac{5} 
{2}\)}{\(m = -1\)}{}}
\LANGes{
\Question{
Dadas las rectas \(p\) 
y \(q\). Determina el punto 
\(m\in \mathbb{R}\) suponiendo que la recta \(p\) es
paralela a la recta \(q\).
\[ 
p\colon x+4y-3 = 0,\qquad \begin{aligned}[t] q\colon x& = 1 + mt,&
\\y & = 2 - 3t;\ t\in \mathbb{R}
\\ \end{aligned}
\]}
{\(m = 12\)}{\(m = -\frac{1}
{12}\)}{\(m = 4\)}{\(m = \frac{5} 
{2}\)}{\(m = -1\)}{}}
\End

\Data\ID{9000106805}
\Updated{2020-07-17 08:07:21}
\Level{C}
\SubArea{Analytic geometry in a plane}
\LANGen{
\Question{
Given points \(A = [0;5]\),
\(B = [6;1]\),
\(C = [7;9]\),
find the direction vector of the line passing through the point
\(A\) and the midpoint of
the segment \(BC\) (i.e. the
median of the triangle \(ABC\) 
through the vertex \(A\)).}
{\((1;0)\)}{\((1;8)\)}{\((1;9)\)}{\((6.5;5)\)}{}{}}
\LANGpl{
\Question{
Dane są punkty \(A = [0;5]\),
\(B = [6;1]\),
\(C = [7;9]\),
wyznacz wektor kierunkowy środkowej trójkąta wychodzącej z wierzchołka
\(A\).}
{\((1;0)\)}{\((1;8)\)}{\((1;9)\)}{\((6{,}5;5)\)}{}{}}
\LANGcs{
\Question{
Pro daný trojúhelník \(ABC\) 
z nabízených možností vyberte směrový vektor
přímky, na které leží jeho těžnice na stranu
\(BC\). Souřadnice vrcholů
trojúhelníka jsou: \(A = [0;5]\),
\(B = [6;1]\),
\(C = [7;9]\).}
{\((1;0)\)}{\((1;8)\)}{\((1;9)\)}{\((6{,}5;5)\)}{}{}}
\LANGsk{
\Question{
Pre daný trojuholník \(ABC\) 
z ponúknutých možností vyberte smerový vektor
priamky, na ktorej leží jeho ťažnica na stranu
\(BC\). Súradnice vrcholov
trojuholníka sú: \(A = [0;5]\),
\(B = [6;1]\),
\(C = [7;9]\).}
{\((1;0)\)}{\((1;8)\)}{\((1;9)\)}{\((6{,}5;5)\)}{}{}}
\LANGes{
\Question{
Para un triángulo dado \(ABC\) 
selecciona de la siguiente lista un vector director de una recta 
 en la que se encuentra la mediana al lado
\(BC\). Las coordenadas de los vértices del triángulo \(ABC\) son: \(A = [0;5]\),
\(B = [6;1]\),
\(C = [7;9]\).}
{\((1;0)\)}{\((1;8)\)}{\((1;9)\)}{\((6.5;5)\)}{}{}}
\End

\Data\ID{9000106806}
\Updated{2021-12-05 10:12:33}
\Level{C}
\SubArea{Analytic geometry in a plane}
\LANGen{
\Question{
Given points \(A = [0;5]\),
\(B = [6;1]\),
\(C = [7;9]\),
find the direction vector of the line passing through the point
\(A\) and perpendicular
to the segment \(BC\) 
(i.e. the line which contains the altitude of the triangle
\(ABC\) through
the point \(A\)).}
{\((8;-1)\)}{\((1;8)\)}{\((1;9)\)}{\((-9;1)\)}{}{}}
\LANGpl{
\Question{
Dane są punkty \(A = [0;5]\),
\(B = [6;1]\),
\(C = [7;9]\), wyznacz  wektor kierunkowy  wysokości trójkąta  opuszczonej na bok   BC.}
{\((8;-1)\)}{\((1;8)\)}{\((1;9)\)}{\((-9;1)\)}{}{}}
\LANGcs{
\Question{
Pro daný trojúhelník \(ABC\) 
z nabízených možností vyberte směrový vektor
přímky, na které leží jeho výška na stranu
\(BC\). Souřadnice vrcholů
trojúhelníka jsou: \(A = [0;5]\),
\(B = [6;1]\),
\(C = [7;9]\).}
{\((8;-1)\)}{\((1;8)\)}{\((1;9)\)}{\((-9;1)\)}{}{}}
\LANGsk{
\Question{
Pre daný trojuholník \(ABC\) 
z ponúknutých možností vyberte smerový vektor
priamky, na ktorej leží jeho výška na stranu
\(BC\). Súradnice vrcholov
trojuholníka sú: \(A = [0;5]\),
\(B = [6;1]\),
\(C = [7;9]\).}
{\((8;-1)\)}{\((1;8)\)}{\((1;9)\)}{\((-9;1)\)}{}{}}
\LANGes{
\Question{
Para un triángulo dado \(ABC\) 
selecciona de la siguiente lista un vector director de una recta 
 en la que se encuentra la altura del lado
\(BC\). Las coordenadas de los vértices del triángulo son: \(A = [0;5]\),
\(B = [6;1]\),
\(C = [7;9]\).}
{\((8;-1)\)}{\((1;8)\)}{\((1;9)\)}{\((-9;1)\)}{}{}}
\End

\Data\ID{9000106807}
\Updated{2021-12-05 10:12:28}
\Level{C}
\SubArea{Analytic geometry in a plane}
\LANGen{
\Question{
Consider the points \(A = [0;5]\),
\(B = [6;1]\),
\(C = [7;9]\) and the
triangle \(ABC\).
Find the direction vector of the line which is the perpendicular bisector of the side
\(b\) 
(i.e. the line through the midpoint of the side
\(AC\) which is perpendicular
to the segment \(AC\)).}
{\((4;-7)\)}{\((7;4)\)}{\((7;9)\)}{\((7;-9)\)}{}{}}
\LANGpl{
\Question{
Dane są wierzchołki trójkąta  \(ABC\), gdzie  \(A = [0;5]\),
\(B = [6;1]\),
\(C = [7;9]\).
Wyznacz wektor kierunkowy symetralnej  boku AC.}
{\((4;-7)\)}{\((7;4)\)}{\((7;9)\)}{\((7;-9)\)}{}{}}
\LANGcs{
\Question{
Pro daný trojúhelník \(ABC\) 
z nabízených možností vyberte směrový vektor osy strany
\(AC\).  Souřadnice vrcholů
trojúhelníka jsou: \(A = [0;5]\),
\(B = [6;1]\),
\(C = [7;9]\).}
{\((4;-7)\)}{\((7;4)\)}{\((7;9)\)}{\((7;-9)\)}{}{}}
\LANGsk{
\Question{
Pre daný trojuholník \(ABC\) 
z ponúknutých možností vyberte smerový vektor osy strany
\(AC\). Súradnice vrcholov
trojuholníka sú: \(A = [0;5]\),
\(B = [6;1]\),
\(C = [7;9]\).}
{\((4;-7)\)}{\((7;4)\)}{\((7;9)\)}{\((7;-9)\)}{}{}}
\LANGes{
\Question{
Para un triángulo dado  \(ABC\) 
 selecciona de la siguiente lista un vector director de la recta en la que se encuentra el lado 
\(AC\). Las coordenadas de los vértices del triángulo son: \(A = [0;5]\),
\(B = [6;1]\),
\(C = [7;9]\).}
{\((4;-7)\)}{\((7;4)\)}{\((7;9)\)}{\((7;-9)\)}{}{}}
\End

\Data\ID{9000106808}
\Updated{2021-12-05 10:12:34}
\Level{C}
\SubArea{Analytic geometry in a plane}
\LANGen{
\Question{
Consider the points \(A = [0;5]\),
\(B = [6;1]\),
\(C = [7;9]\) and the
triangle \(ABC\).
Find the direction vector of the line which is the bisector of the angle
\(ACB\) 
(i.e. the line which splits the internal angle at the point
\(C\) into
two angles with equal measures).}
{\((2;3)\)}{\((6;-4)\)}{\((7;9)\)}{\((7;8)\)}{}{}}
\LANGpl{
\Question{
Dane są wierzchołki trójkąta \(ABC\), gdzie \(A = [0;5]\),
\(B = [6;1]\),
\(C = [7;9]\)  
wyznacz wektor kierunkowy  dwusieczną kąta
\(ACB\).}
{\((2;3)\)}{\((6;-4)\)}{\((7;9)\)}{\((7;8)\)}{}{}}
\LANGcs{
\Question{
Pro daný trojúhelník \(ABC\) 
z nabízených možností vyberte směrový vektor osy úhlu při vrcholu
\(C\). Souřadnice vrcholů
trojúhelníka jsou: \(A = [0;5]\),
\(B = [6;1]\),
\(C = [7;9]\).}
{\((2;3)\)}{\((6;-4)\)}{\((7;9)\)}{\((7;8)\)}{}{}}
\LANGsk{
\Question{
Pre daný trojuholník \(ABC\) 
z ponúknutých možností vyberte smerový vektor osy uhlov pri vrchole
\(C\). Súradnice vrcholov
trojuholníka sú: \(A = [0;5]\),
\(B = [6;1]\),
\(C = [7;9]\).}
{\((2;3)\)}{\((6;-4)\)}{\((7;9)\)}{\((7;8)\)}{}{}}
\LANGes{
\Question{
Para un triángulo dado \(ABC\) 
selecciona de la siguiente lista un vector director de la bisectriz del ángulo del vértice
\(C\). Las coordenadas de los vértices del triángulo son: \(A = [0;5]\),
\(B = [6;1]\),
\(C = [7;9]\).}
{\((2;3)\)}{\((6;-4)\)}{\((7;9)\)}{\((7;8)\)}{}{}}
\End
